\documentclass[11pt]{amsart}
\usepackage[T1]{fontenc}
\usepackage{lmodern}
\usepackage{amsmath,amssymb,amsthm}
\usepackage[hyphens]{url}
\usepackage[hidelinks]{hyperref}

\ifdefined\pdfinfoomitdate
  \pdfinfoomitdate=1
\fi
\ifdefined\pdftrailerid
  \pdftrailerid{}
\fi

\newtheorem{theorem}{Theorem}[section]
\newtheorem{proposition}[theorem]{Proposition}
\newtheorem{lemma}[theorem]{Lemma}
\newtheorem{corollary}[theorem]{Corollary}
\theoremstyle{definition}
\newtheorem{definition}[theorem]{Definition}
\theoremstyle{remark}
\newtheorem{remark}[theorem]{Remark}
\numberwithin{equation}{section}
\newcommand{\PA}{\mathsf{PA}}
\newcommand{\N}{\mathbb N}
\newcommand{\Q}{\mathbb Q}
\newcommand{\QQi}{\mathbb Q[\mathrm i]}
\newcommand{\pair}[2]{\langle #1,#2\rangle}
\newcommand{\disk}{\mathsf{Disk}}
\newcommand{\Probe}{\mathsf{Probe}}
\newcommand{\Zero}{\mathsf{Zero}}
\newcommand{\Good}{\mathsf{Good}}
\newcommand{\Gph}{\mathsf{Graph}}
\newcommand{\Cert}{\mathsf{Cert}}
\newcommand{\Row}{\mathsf{Row}}
\newcommand{\Trace}{\mathsf{Trace}}
\newcommand{\Axioms}{\mathsf{A}}
\newcommand{\lz}{\ell}
\newcommand{\ip}[2]{[#1,#2]}
\newcommand{\ff}{\operatorname{fl}}
\newcommand{\re}{\operatorname{Re}}
\newcommand{\im}{\operatorname{Im}}
\newcommand{\code}[1]{\nolinkurl{#1}}
\setlength{\emergencystretch}{2em}
\raggedbottom
\setcounter{tocdepth}{1}

\begin{document}
\title[The windowed Landau reduction in arithmetic]
  {Landau's windowed ceiling reduction in arithmetic}
\author{The \texttt{landau-windows} project}
\begin{abstract}
We formulate the windowed Landau ceiling reduction in first-order
Peano arithmetic using natural-number codes for rational data and
finite certificates. Ordinary induction gives the signal enclosures,
window search, power-sum bounds, uniform Euler polynomial estimates,
rational factorization, and paired contours with stable counts.
Finite Mellin quadrature transports search failure to the contour
coefficients, where it contradicts a finite power-sum bound.
We verify the computed threshold and the returned rational witness.
We also give the classical interpretation for the Riemann zeta
function.
\end{abstract}
\maketitle
\tableofcontents
\clearpage

\section{Statement of the arithmetic theorem}

The classical ceiling step converts a bound on a smoothed prime-counting
error throughout recurring intervals of the form $[Y,Y^{3/2}]$ into a
zero-free region for the Riemann zeta function. The companion manuscript
\cite{companion} treats a single-zero certificate. Here the input is
a rational disk containing arbitrarily fine small-value probes. This
form handles clusters and repeated zeros through the size of the
function value. It leaves the threshold independent of the probe
point.

The ACL2 development certifies an end-to-end reduction from these probes
to positive rational lower endpoints. Its final theorem is
\code{lw-euler-window-encoded-zero-reduction}. We prove its first-order
arithmetic translation by direct derivations on finite codes.
The ordinary mathematical interpretation is given in
Section~\ref{sec:comparison}.

\subsection{The target formula}

Write $D=(c,b,h,C,p,q)$ for a tuple of rational data and two naturals.
The predicate $\Good(D)$ asserts
\[
 q>0,\quad 0\le p<q,\quad C\ge 0,\quad
 251<c<252-p/q,\quad 0<b,\quad h\in\Q .
\]
The source allows \(p\le q\), and the strict interval
for \(c\) forces \(p<q\).
The physical disk is centered at $252+\mathrm i h$ and has radius
$c-b$ when this is nonnegative. The probe predicate is false for
$c-b<0$, as in the source disk test.
A point code $u$ holds an offset in $\QQi$. For positive $e\in\Q$,
$\Probe(D,u,e)$ says that $u$ is in the offset disk, and that the
finite Euler evaluator with request $e$ at the point
$252+\mathrm i h+u$ returns a pair whose squared norm is at most $e^2$.
Here and throughout, the evaluator means the program specified by
\code{lw-complex-euler-continuation-value}, with graph notation as in
Section~\ref{sec:graphs}. Define
\[
 \Zero(D) \ \equiv\
 \forall e\bigl((e\in\Q\ \wedge\ e>0)\ \longrightarrow\
                  \exists u\,\Probe(D,u,e)\bigr).
\]

Let $\mathsf{Prepared}(D,T,r)$ express that the prescribed threshold and
request programs return $T$ and $r$ from $D$. The threshold is computed
by \code{lw-euler-window-ceiling-threshold}. The request is computed by
\code{lw-euler-window-zero-sample-request}, which reads only the disk
data. Let $\Cert(D,Y,x,L,U)$ assert that the window search returns the
address $x$ and endpoints $L,U$, with
\[
 1\le Y\le x,\qquad x^2\le Y^3,\qquad 0<L\le U
\]
and identifies $(L,U)$ as the signal interval for $x,C,p,q$ at
request $1/4$. The square comparison is the entire occurrence of the
right edge of the window in the arithmetic formula.

Let $\mathsf{Eligible}(D,T,r,Y,u,e)$ abbreviate
\[
  Y\in\Q\ \wedge\ T\le Y\ \wedge\ \Probe(D,u,e)\ \wedge\ 4e\le r.
\]

\begin{definition}[Arithmetic target]\label{def:target}
The target sentence $\mathcal W$ is
\begin{equation}\label{eq:target}
\begin{gathered}
\forall D\,\bigl(\Good(D)\longrightarrow\exists T,r\,[\\
 T\in\N_{\ge1} \ \wedge\ r\in\Q_{>0} \ \wedge\ \mathsf{Prepared}(D,T,r)\\
 {}\wedge\ \forall Y,u,e\bigl(\mathsf{Eligible}(D,T,r,Y,u,e)
       \longrightarrow\\
 \exists x,L,U\ \Cert(D,Y,x,L,U)\bigr)]\bigr).
\end{gathered}
\end{equation}
\end{definition}

Quantifiers in \eqref{eq:target} range over natural numbers. Quantifiers
displayed over rational domains below abbreviate a code-validity guard.
The \emph{graph} of each relevant program is given by a formula even
before its totality is proved. In particular, \eqref{eq:target} asserts
that the threshold and request exist uniformly. When $\Zero(D)$ holds,
ordinary existential instantiation at $e=r/4$ suffices for every later
window. This step uses one witness for a fixed request and uses the same
$T$ for all $Y$.

For a rational $Y\ge1$, call the arithmetic ceiling test satisfied
when the lower endpoint at request $1/4$ is nonpositive at every
rational $x$ with $Y\le x$ and $x^2\le Y^3$. If $\mathcal W$ holds,
arbitrarily large starts satisfying this test imply $\neg\Zero(D)$.
Indeed, the prepared threshold gives one such start, and $\Zero(D)$
supplies a probe at $r/4$. Applying \eqref{eq:target} yields a lower
endpoint that is positive at a valid address, a contradiction.
This corollary to the target is itself a first-order implication.

\section{The arithmetic language of the proof}
\label{sec:graphs}

We work in first-order $\PA$ with symbols $0,S,+,\cdot$ and equality.
Order, divisibility, exponentiation, finite sums, rational operations,
and bounded loops are definitional abbreviations. Induction may be
applied to formulas with parameters.

\subsection{Numbers, pairs, lists, and interval comparisons}

\begin{lemma}[Coding pairs and signed integers]
The map
\[
 \pair{a}{b}=1+(a+b)^2+a,\qquad a,b\in\N,
\]
is injective, exceeds both coordinates, and has a PA-definable
bounded decoding operation. Finite lists are representable by
$0$ for the empty list and $\pair{a}{t}$ for cons. Signed integers are
representable by $2n$ for $n\ge0$ and $2n+1$ for $-(n+1)$.
\end{lemma}
\begin{proof}
At fixed $s=a+b$ the range lies in $[1+s^2,1+s^2+s]$. The lower
endpoint for $s+1$ exceeds the preceding upper endpoint since
$(s+1)^2-s^2-s=s+1$. For two different sums the ranges are
disjoint by transitivity and monotonicity of squaring. At a fixed sum,
equality determines $a$ and then $b$. Both coordinates are at most
$s$, and $1+s^2>s$ for $s=0,1$ and by multiplication for $s\ge2$.
For a given input $w$ search over $a,b<w$ and verify the equation.
An unsuccessful search gives the distinguished failure value.
Cons decoding decreases the second argument along the list, so ordinary
induction on the code proves termination and unique list contents.
Existence and uniqueness of quotient and remainder on division by~2
follow by induction on the input. The signed operations are given by
the usual cases on the two signs, with natural addition, subtraction
of a smaller magnitude, and multiplication in each case.
\end{proof}

A rational code represents $(a,d)$ with an integer $a$ and a positive
natural $d$. One may either reduce the fraction or work with
representatives and cross-multiplication. The latter is sufficient
in the direct arithmetic lemmas. Replaying an equality test on
canonical fractions uses reduction by the greatest common divisor.
The greatest divisor of both magnitudes can be found by a bounded
search when at least one magnitude is positive. Uniqueness of the
reduced representative follows from Euclid's lemma and the positive
denominator convention. Euclid's lemma has the ordinary integer proof
from the Euclidean division algorithm: the nonzero remainder decreases
at each step, and the last common divisor is an integer linear
combination of the original two numbers by substitution in those
division equations. All induction indices in this argument are naturals.
For $a/d$ and $a'/d'$, use $ad'\le a'd$ for order and replace the
numerator and denominator under $+$ or $\cdot$ by the signed integer
formulas. Cross-multiplication proves invariance under choice of
representatives.

A code for $z\in\QQi$ is a pair of rational codes. Its squared norm
is the rational $\re(z)^2+\im(z)^2$, and
\[
 \disk(z,R) \ \equiv\ R\ge0\ \wedge\
        \re(z)^2+\im(z)^2\le R^2.
\]
A disk estimate written $|z|\le R$ below abbreviates exactly this formula.
The inequalities $|u+v|\le A+B$ and $|uv|\le AB$ from
$|u|\le A, |v|\le B$, where $A,B$ are nonnegative, require only
rational algebra. If $u=(a,b)$ and $v=(c,d)$ then
\[
 (ac+bd)^2+(ad-bc)^2=(a^2+b^2)(c^2+d^2).
\]
The assumptions imply $ac+bd\le AB$ by a squared
comparison and nonnegativity of $AB$. Expanding the squared norm of
$u+v$ gives the sum rule. The same identity with the usual complex
product gives the product rule. Applying the sum rule to differences
gives the error-propagation rule used throughout.

\subsection{Finite computation graphs}

For any fixed finite program $F$, write $\Gph_F(\mathbf x,y)$ for
existence of a finite computation trace with input $\mathbf x$, output
$y$, and valid equations between each instruction and its arguments.
One possible trace code is a list of records in evaluation order.
A record contains a procedure tag, input and output codes, branch tags
when needed, and pointers to argument values and completed subcalls
at smaller positions. The check consists of bounded list walks and tests
on codes. A subcall record is checked by the equation for its procedure
and the stipulated input and output. Checks for conditionals require the
selected branch and its predicate. Consing a new record after valid
subcall traces, with pointer indices shifted by the preceding lengths,
is a finite list operation.

\begin{lemma}[Reduction to natural-number induction]\label{lem:recursion}
Suppose the primitive operations of a deterministic procedure are
PA-provably total and single-valued. For a recursive call let $G$ be
the guard for the branch containing the call. Suppose $\PA$ proves that
a measure $m$ gives a natural on the domain and that the subcall
stays in the domain and has measure strictly smaller whenever $G$
holds and preceding subcalls have returned their specified values.
Then $\PA$ proves totality and single-valuedness of $\Gph_F$ on
that domain.
\end{lemma}
\begin{proof}
Apply strong induction to the assertion that every input of measure
at most $k$ has a trace and a unique result. To treat an input of
measure $k+1$, evaluate its finite list of branch tests and
nonrecursive operations. In the order required by the selected
branch, the measure condition and induction hypothesis give each
subresult and its uniqueness. The trace constructor then gives the
parent record. Any valid trace has the same branch as the one chosen
by the primitive tests. Each justified subcall has smaller measure,
so all its results agree with those already obtained. The final
operation is deterministic. The case of measure zero uses the absence
of a strictly smaller natural measure. The same argument treats a finite
group of mutually recursive procedures by including the procedure tag
in the inputs being quantified.
\end{proof}

For later use, suppose a proof has induction cases with guards
$G_i(\mathbf x)$ and substitutions $\sigma_{i,j}$, and suppose $\PA$
proves all the decreases
\[
 G_i(\mathbf x)\ \longrightarrow\
 m(\mathbf x\sigma_{i,j})<m(\mathbf x).
\]
If the base formula follows when no guard holds and each guarded
step proves the formula from its listed hypotheses, apply strong
induction to $\forall\mathbf x(m(\mathbf x)\le k\to\varphi(\mathbf x))$.
The guard in force selects the step, and the smaller indices justify
all its hypotheses. This proves $\forall\mathbf x\,\varphi(\mathbf x)$
using a single ordinary-induction schema instance after expansion of
strong induction. The proof uses arithmetically coded inputs and
applies to formulas with unbounded first-order quantifiers.

\section{Finite interval estimates}

For rational intervals $I=[l,u]$ and $J=[l',u']$, use endpoint addition,
$[l-u',u-l']$ for subtraction, and $[cl,cu]$ for multiplication by
$c\ge0$. Their widths are additive for addition and subtraction and
multiply by $c$ for scaling. All nesting statements in this section
are inequalities between finite rational expressions.

\subsection{Logarithm rectangles and the harmonic correction}
\label{sec:log-rectangles}

For $x\ge1$, let $n$ be any positive integer and $t=(x-1)/n$.
Define
\[
 L_n(x)=t\sum_{j=1}^n (1+jt)^{-1},\qquad
 U_n(x)=t\sum_{j=0}^{n-1}(1+jt)^{-1}.
\]
Their difference is exactly $t(1-1/x)$, by cancellation of the
common terms. The lower sum can only increase and the upper sum
can only decrease when each cell is subdivided. On a cell
$[\alpha,\beta]$, this is the fact that the reciprocal at each
subinterval's right endpoint is at least $1/\beta$ and the
reciprocal at each left endpoint is at most $1/\alpha$.
Summing the estimates proves nesting for refinements.
For any two meshes, the mesh with the product of their subdivision
counts refines both. In particular, each lower sum is at most each
upper sum. This uses only induction on the finite sums.

If $x-1=A/B$ with $A\ge0$, $B\ge1$, and the request is $e=a/d>0$,
the choice $n=1+Ad$ gives $(x-1)/n\le e$, since
$Ad\le anB$. Thus there is a computed interval of width at most $e$.
These are the rectangle formulas in the arithmetic evaluator.
The use of a common refinement also lets $\PA$ compare enclosures
at independently chosen precisions.

For a finite harmonic sum write $H_n=\sum_{j=1}^n 1/j$.
For log intervals $I(n)$ and $I(m)$, we may compare their refinements
using the cells from $n$ to $m$, provided $1\le n\le m$. On each
$[k,k+1]$ reciprocal rectangles are between $1/(k+1)$ and $1/k$.
For a subdivision of $[1,m]$ that includes $n$ and all these integer
breakpoints, this gives
\[
 H_m-H_n\ \le\ (\hbox{lower or upper sum from }n\hbox{ to }m)
        \ \le\ H_m-H_n+\frac1n-\frac1m .
\]
The right side uses the exact telescoping sum of $1/k-1/(k+1)$.
Uniform meshes may be refined to include the finitely many
breakpoints, and a common further refinement gives the same
comparison for independently chosen endpoint enclosures. An exact
arithmetic version is that for any specified finite accuracy
$\eta>0$, two refined approximations to $H_m-\log(m)$ and
$H_n-\log(n)$ have difference in
$[-1/n+1/m-2\eta,2\eta]$.
Here the notation $\log$ in the finite approximations denotes
either endpoint of an interval of width at most $\eta$.
Every comparison follows by adding the interval errors after the
finite common refinement.

The computed interval used to name the Euler constant is
\[
 \ip{H_n-U(n)-1/n}{H_n-L(n)},\qquad
 I(n)=[L(n),U(n)],
\]
where the log request is $e/2$ and $n$ is chosen to satisfy
$1/n\le e/2$. Its width is at most $e$. The cell comparisons give
pairwise overlap of these naming intervals: for $m\ge n$, compare
the two refined sums from $1$ to $n$ with the sum from $n$ to $m$.
For the lower endpoint at $n$ versus the upper endpoint at $m$,
use the upper cell bound above and nesting under refinement. For the
lower endpoint at $m$ versus the upper endpoint at $n$, use the lower
cell bound and nesting. Thus the arithmetic argument uses a consistent
family of intervals and bounds its widths by explicit rationals.

\subsection{Exponentiation and rational powers}
\label{sec:finite-exp}

Integer powers are given by repeated multiplication. For a positive
integer $q$ and $x\ge1$, $0\le p\le q$, the bracket $[1,x]$ satisfies
$1^q\le x^p\le x^q$. Bisect the bracket and compare the $q$th power of
the rational midpoint with $x^p$. Monotonicity of powers on nonnegative
rationals, proved by induction on $q$, preserves both inequalities.
Induction on the bisection count proves halving of the width until
a prescribed rational precision is reached. Termination may be
bounded by any $k$ satisfying $(1+k)e\ge x-1$, since
$2^k\ge1+k$ follows by induction. The same argument brackets a square
root below and above using only squared inequalities.

For completeness, a finite exponential tail comparison uses
$E_N(z)=\sum_{j=0}^N z^j/j!$ in $\QQi$. With
$D=|\re(z)|+|\im(z)|$, each term has squared norm at most
$(D^j/j!)^2$, by the disk product rule. Once $N+2\ge2D$, successive
terms beyond degree $N$ have upper radii in ratio at most $1/2$.
Their sum through any finite degree is bounded by $2D^{N+1}/(N+1)!$.
To choose $N$ at request $e$, start beyond $2D$, bound the first
term radius by a rational computation, and use the halving estimate
and $2^k\ge1+k$ to find a sufficient additional count.
This proves finite agreement of partial sums at independently
specified degrees using only rational arithmetic and induction.

\subsection{The signal interval}

For $x\ge1$ and $k\ge1$, set $w_x(k)=\max(0,x/k-1)$.
For $x=a/d$ with positive numerator and denominator, the bounded
prefix through $a$ includes every integer $k\le x$. Beyond this,
the weights vanish. Test primality of any $l\le a$ by trial division
at $2,\ldots,l-1$, and for each prime test $l^j=k$ over
$1\le j\le k$. Since $l\ge2$ gives $l^{k+1}>k$, this bound
covers every possible exponent. If there are two positive prime
powers for the same $k$, Euclid's lemma implies that the two primes
coincide. The positive exponent is then unique by monotonicity.
The interval for $\Lambda(k)$ is the log rectangle interval
for the prime base or the zero interval in the absence of such a power.

Let $\mathcal I(x,C,p,q,e)$ be the signal interval assembled
from these terms with
\[
 \epsilon=\frac{e}{1+(a+2)x+C}.
\]
Use $\epsilon$ for all primitive widths. Form the weighted finite
sum, subtract $x$ times the log interval at $x$, add $x$ times one
plus the Euler interval, and subtract $C$ times the power bracket.
The first width is at most $ax\epsilon$, since each weight
is at most $x$. The next two widths are at most $x\epsilon$ each.
The power contribution has width at most $C\epsilon$.
Consequently $\mathcal I$ has width at most $e$.
This proves a PA statement about the constructed finite endpoints.
Section~\ref{sec:comparison} identifies their external interpretation.

\section{The finite window and search}

\subsection{The interval-search book in PA}
\label{sec:interval-search}

The following derivations cover the eleven theorem events in
\code{acl2/interval-search.lisp}. The search theorems quantify over
unrestricted ACL2 objects. Their interpretation therefore includes
atomic inputs and cons chains with atomic tails.

Use $0$ for \texttt{nil} and implement the nonempty objects with
disjoint tags. In particular, the cons code is
$\pair{0}{\pair{a}{b}}$, and the rational code is
$\pair{1}{r}$ for a canonical rational code $r$. A complex number
with nonzero imaginary part has tag $2$ and a pair of canonical
rational codes. Reserve other atom codes for the other ACL2 atoms,
including \texttt{t}. On atoms the projections \texttt{car} and
\texttt{cdr} return \texttt{nil}, and \texttt{endp} tests whether
the input is an atom. The truth test is inequality to $0$.
For logical \texttt{<}, view numbers as pairs of real and imaginary
rationals and other objects as $(0,0)$. Compare the two coordinates
lexicographically, as in the ACL2 completion of \texttt{<}.
For \texttt{+}, \texttt{-} and \texttt{*}, use the corresponding
complex operations with the same default $(0,0)$ and reduce results
with zero imaginary part to the rational tag. These operations are
PA-total by the coding in Section~\ref{sec:graphs}.

Write $n(L)$ for the number of successive cons cells reached by
\texttt{cdr} on $L$. Its recursion stops on any atom and decreases
the object code at each cons. This is the graph interpretation of
\texttt{len}. Write $L[i]$ for the \texttt{car} of the $i$th tail,
with $i=0$ at the input. On natural indices the ACL2 \texttt{nth}
recursion has
$(a\mathbin{::}B)[0]=a$ and $(a\mathbin{::}B)[i+1]=B[i]$,
and returns \texttt{nil} at an atom. Let $P_t(a)$ be the truth of
\texttt{(< threshold (car a))} at threshold $t$.
For an interval object $a$ its first projection is the lower
endpoint and its second projection is the upper endpoint.

\begin{lemma}[The first successful index]
\label{lem:first-above}
Let $F(t,L)$ denote \code{lw-first-above} and $H(t,L)$ denote
\code{lw-some-lower-above}. In $\PA$, if no $i<n(L)$ satisfies
$P_t(L[i])$, then $F(t,L)=H(t,L)=\texttt{nil}$. Otherwise the
least such $i$ satisfies
\[
  H(t,L)=\texttt{t},\qquad
  F(t,L)=\operatorname{cons}(i,L[i]),
\]
where the first cons field is the rational tag for the natural $i$.
\end{lemma}
\begin{proof}
Lemma~\ref{lem:recursion} gives both graphs since the recursive
calls replace $L$ by its cdr in the cons case. Apply ordinary
induction to the assertion for all $t,L$ of length $n$.
At $n=0$ the input is atomic and both defining equations return
\texttt{nil}. At $n+1$ write $L=a\mathbin{::}B$. If $P_t(a)$ holds,
the minimum is $0$ and the equations give \texttt{t} and
$\operatorname{cons}(0,a)$. If $P_t(a)$ fails and there is no
successful tail index, the induction hypothesis gives two
\texttt{nil} results and both equations preserve them.
In the remaining branch the first successful tail index is $j$.
Every index from $1$ to $j$ in $L$ corresponds to an index
below $j$ in $B$, so all these checks fail. The check at
$j+1$ succeeds by $L[j+1]=B[j]$. The return rule for $F$ increments
the tail result's natural index and preserves its sample, and
the rule for $H$ preserves truth.
\end{proof}

The equivalence of the truth tests for $F$ and $H$ is
\code{lw-first-above-complete}. The index in a successful return
is a natural below $n(L)$, giving \code{lw-first-above-index-natural}
and \code{lw-first-above-index-bound}. The returned sample equals
$L[i]$, giving \code{lw-first-above-nth}. Its first projection passes
the threshold test by the choice of $i$, giving
\code{lw-first-above-lower-gt}.

For the other interval lemmas write $[a,b]$ for an ordered pair of
rational endpoints satisfying $a\le b$. The \code{lw-contains}
condition for $v$ is exactly $a\le v\le b$ with rational $v$.
If $a\le v\le b$ and $c\le w\le d$, adding the lower and upper
inequalities gives $a+c\le v+w\le b+d$. Adding the inequalities
$-d\le -w\le -c$ to the bounds on $v$ gives
$a-d\le v-w\le b-c$. For rational scaling, multiplying by a
nonnegative $k$ gives $ka\le kv\le kb$, and multiplying by a
negative $k$ gives $kb\le kv\le ka$. In each case the new
endpoints are ordered by transitivity. These are the three
\code{lw-contains} arithmetic lemmas. All order rules here are
integer inequalities after multiplication by positive denominators.
If $b-a\le e$ and $t+e<v\le b$, then $t+e<b\le a+e$ and
cancellation gives $t<a$. This proves \code{lw-lower-gt-of-margin}.

The graph for \code{lw-encloses-values} is defined by descent on
the sample list. Its truth entails equal cons-chain lengths and
containment of the corresponding value at every index below that
length. Induct on the sample length. The base branch requires
the value input to be an atom. The cons branch requires a cons
value input, tests containment at index $0$, and tests
\code{lw-encloses-values} on both tails. At each successor index
the claim follows from the two \texttt{nth} equations and the
induction hypothesis. The same head and successor equations
give \code{lw-widths-at-most} as the conjunction of the width
checks up to $n(L)$, and \code{lw-some-value-above} as the
existence of a successful check up to the value list's length.
The first two recursions decrease the sample code and the last
decreases the value code.

For \code{lw-first-above-value-gt}, the returned index $i$ has a
lower endpoint above $t$. Containment gives a corresponding
value at least that endpoint, so the value at index $i$ exceeds
$t$. For \code{lw-search-complete-of-margin}, take the index
whose value exceeds $t+e$. The width check and containment
imply $t<a$ at this index by the margin inequality above.
Lemma~\ref{lem:first-above} gives a successful search. This
discharges the two remaining theorem events with the rational
threshold and error premises stated in the book.

\subsection{Geometry of the rational window}

Let $Y\ge1$. Bisect a bracket for the square root of $Y^3$ starting
from $[Y,Y^2]$ and stop with a rational bracket $[B,B^+]$ of
width at most $\delta>0$. Thus $Y\le B\le B^+$,
$B^2\le Y^3\le(B^+)^2$, and $B^+-B\le\delta$.
Take the increasing list $Y,Y+\delta,\ldots$ up to $B$ and
append $B$ if needed. The source bisects after the halving depth of
$Y^2-Y$ at $\delta$ and uses precision fuel for width $B-Y$ to
construct the mesh. For width $w=A/D$ and step
$\delta=P/Q$ in reduced form with $w\ge0$ and $\delta>0$,
the fuel $AQ$ satisfies $w\le AQ\delta$. It therefore
suffices for all mesh steps. The halving computation with the same
fuel either stops at the width test or uses enough halvings, since
$2^k\ge k+1$ by induction. These tests give the stated bracketing
width and the indicated list at the source's exact outputs.
There is a mesh point at distance at most
$\delta$ from each rational $x\ge Y$ with $x^2\le Y^3$.
For $x\le B$, compare $x$ to the adjacent mesh points, using a
finite scan. For $x>B$, use $B$ and $x\le B^+$, the latter
following by squared comparison. All mesh addresses are valid
window addresses. Each loop is bounded by a natural majorizing
$(B-Y)/\delta$, and subsequent induction is over that finite mesh.
This proves the geometry used in \code{rational-window.lisp}.

\subsection{Failure and the interval comparison}

Fix a rational upper address $B\ge1$ and a common integer
prefix $N\ge B$, which can be the numerator of $B$. For rational
addresses $x,y\in[1,B]$, set $d=|x-y|$.
The lower and upper brackets at each occurrence can have different
primitive requests. Refining the reciprocal partitions to their common
breakpoints gives a rational number in each original log interval,
using the same refined sum on the shared part of the two ranges.
In particular the two chosen log representatives differ by at most
$d$. They are nonnegative and at most $B$, since each reciprocal in
the sum on $[1,B]$ is at most one. The products by $x,y$ differ
by at most $2Bd$ by writing $xa-yb=x(a-b)+(x-y)b$.

There is a common representative for each pair of prime log
intervals, by the same refinement. For the index $k$ this number
is at most $k$, and the ramp difference is at most $d/k$.
The latter follows by the three cases for $x/k-1$ and $y/k-1$
being below or above zero. Thus each weighted term differs by
at most $d$, and the common finite prefix contributes at most
$Nd$. Padding to the common prefix has value zero beyond the
address. For the harmonic correction choose a common representative
of the two intervals and the interval $[0,1]$ from $n=1$.
Three intervals that overlap pairwise have common intersection
because the largest lower endpoint is at most every upper endpoint.
Multiplication by $1$ plus this representative bounds the difference
in the two harmonic terms by $2d$, which we allow to be $3d$.

The power brackets also compare directly. If $x=y+d$, let $u_y$
be the upper endpoint of the bracket at $y$. Since $u_y\le y$,
we have $u_y+d\ge u_y(1+d/y)$. Raising to the $q$th power
and using $q\ge p$ gives
\[
 (u_y+d)^q\ge y^p(1+d/y)^q\ge y^p(1+d/y)^p=x^p .
\]
Thus the lower power endpoint at $x$ is at most $u_y+d$.
The lower endpoint at $y$ is at most the upper endpoint at $x$
by $y^p\le x^p$. These two comparisons bound the one-sided endpoint
differences in both directions. After multiplying by $C$ and
applying the sign of that contribution, it adds at most $Cd$.
Consequently, if $[l(x),u(x)]$ denotes the full signal interval
and each width is at most $\varepsilon$,
\[
     l(x)-u(y)\le L|x-y|,\qquad L=N+2B+3+C.
\]
Reversing the roles of the points yields the same estimate
with $x,y$ reversed. Adding the widths gives
\begin{equation}\label{eq:transport}
 u(x)\le l(y)+L|x-y|+2\varepsilon .
\end{equation}
If the search finds a mesh address with $l(y)>0$, the data at that
address give the required positive bracket. If it fails, induction on
the search list yields $l(y)\le0$ at all mesh addresses. For each
rational address in the window choose a mesh point at distance at most
$\delta$. Then $u(x)\le L\delta+2\varepsilon$ by
\eqref{eq:transport}. In the implemented search use transport cap
$Y^2$ and its numerator for the prefix bound $N$.
The step is
$\delta=1/(4(L+1))$ and $\varepsilon=1/4$, so this is less than~1.
If $U_G(x)$ is the upper endpoint before the ceiling subtraction
and $P_l(x)$ is the power lower endpoint, this gives
$U_G(x)\le CP_l(x)+1\le(C+1)P_l(x)$ because $P_l(x)\ge1$.
This is the arithmetic input to the coefficient estimates. The
corresponding transport proof is in \code{signal-continuity.lisp}.

\section{Finite estimates for the coefficient argument}

\subsection{Coefficient extraction by a rational orbit}

\begin{lemma}[Arithmetic Cauchy bound]\label{lem:finite-cauchy}
Let $P(X)=\sum_{k=0}^n a_kX^k$ have coefficients in $\QQi$.
Suppose $R>0$ and $M\ge0$ are rational and
$|P(z)|\le M$ at every rational pair with squared norm $R^2$.
Then $|a_jR^j|\le M$ for every $0\le j\le n$.
\end{lemma}
\begin{proof}
Take $\omega=(3+4\mathrm i)/5$. Its squared norm is one.
It has no positive power equal to one. If $\omega^d=1$
for $d\ge1$, then $(3+4\mathrm i)^d=5^d$ in the Gaussian
integers. Reducing modulo five and substituting $\mathrm i=2$
respects both addition and multiplication because $2^2+1=5$.
It carries the left side to $1^d=1$ and the right side to zero.
Preservation under a positive power follows by induction on $d$.
Negative powers are handled by inversion.

For a positive natural $N$ take the exact average
\[
 A_N=\frac1N\sum_{s=0}^{N-1} P(R\omega^s)\omega^{-js}.
\]
The disk hypothesis and finite triangle rule give $|A_N|\le M$.
The term of order $j$ in the expanded average is $a_jR^j$.
For each other order the geometric factor is
$(1-\omega^{(k-j)N})/\bigl(N(1-\omega^{k-j})\bigr)$.
This is a finite geometric sum identity with a nonzero denominator.
Writing $|b|_1=|\re(b)|+|\im(b)|$, we obtain the rational bound
\[
 |A_N-a_jR^j|\le \frac{D}{N},
 \qquad
 D=2\sum_{k\ne j}|a_k|_1 R^k
               \bigl|(1-\omega^{k-j})^{-1}\bigr|_1 .
\]
If $G=|a_jR^j|^2-M^2>0$, set
$\eta=\min(1,G/(2(2M+1)))$ and take an integer $N>D/\eta$.
Then $(M+D/N)^2\le(M+\eta)^2\le M^2+G/2$,
contrary to the sum rule applied to $A_N$ and its error.
This is a finite contradiction using only a computed average.
\end{proof}

Applying the lemma to a difference of polynomials gives coefficient
comparison from any uniform boundary error bound. It applies to
polynomial models at individually chosen degrees. A use on quotient
or factor data still requires a finite polynomial approximation and
the uniform boundary inequality on that approximation.

\subsection{Choosing a separated circle for a finite factor list}

\begin{lemma}[A finite radius selection]\label{lem:radius-selection}
Let $v_1,\ldots,v_m\in\QQi$ satisfy $|v_\nu|\le B$ for a rational
$B\ge0$, and let $0<a<b$ be rational. There is a computed
rational radius $R\in(a,b)$ and a computed $\xi>0$ such that
every rational point $z$ on the circle satisfies
$|z-v_\nu|\ge\xi$ for every $\nu$.
\end{lemma}
\begin{proof}
Use the $m+1$ trial radii $a+i\Delta$ for $1\le i\le m+1$,
where $\Delta=(b-a)/(m+2)$. Consecutive squared radii differ by
at least $2a\Delta$. Put $\eta=a\Delta/2$. For fixed $\nu$,
at most one trial radius has squared distance from $|v_\nu|^2$
at most $\eta$ on the rational line. If every radius were rejected,
assign to it the first index rejecting it. This would give an
injection from a list of length $m+1$ into one of length $m$.
The finite pigeonhole argument is induction on $m$, deleting
one image and its preimage at each step. Thus a bounded scan
picks a radius with
\[
   \bigl| R^2-|v_\nu|^2\bigr|>\eta\quad\text{for every }\nu.
\]
For $|z|^2=R^2$, the expression on the left is
$|\re((z-v_\nu)\overline{(z+v_\nu)})|$.
The disk rules give $|z+v_\nu|\le R+B$. Taking
$\xi=\eta/(R+B)>0$, membership of $z-v_\nu$ in the disk of
radius $\xi$ would put that coordinate in $[-\eta,\eta]$.
This proves even the strict squared separation. Every radius,
index and comparison is given by integer and rational arithmetic.
\end{proof}

If $P(z)=A\prod_\nu(z-v_\nu)$ and $|A|^2\ge A_0^2>0$
for a specified rational $A_0>0$, the norm product identity
proves $|P(z)|^2\ge(A_0\xi^m)^2$ on the chosen circle.
Any uniformly approximating evaluator within $A_0\xi^m/2$
inherits a rational lower bound of $A_0\xi^m/2$. This gives a
finite lower-bound proof for a supplied factor list.
The contour construction of Section~\ref{sec:contour} also
requires a factor list from the evaluator and estimates valid
as its precision changes.

\subsection{Power sums without a separation assumption}

For $m\ge1$ let
\[
 S(n)=\sum_{\nu=1}^m w_\nu z_\nu^n
\]
with $z_\nu,w_\nu\in\QQi$. The entries may repeat. We first prove
the bound for unit-disk nodes and then reverse the sequence.
\begin{lemma}[Prefix domination]\label{lem:prefix}
If $|z_\nu|\le1$ for each $\nu$, and $|S(n)|\le A$ for
$0\le n<m$, then, for every $N\ge0$ and $0\le n\le N$,
\[
 |S(n)|\le A[3(N+1)]^{m-1}.
\]
Here $A$ is a nonnegative rational.
\end{lemma}
\begin{proof}
Induct on $m$ with the list of entries and the cap as parameters.
For $m=1$, the constant term is $w_1$ and the disk product rule
bounds every power. For $m\ge2$, put
$T(n)=S(n+1)-z_1S(n)$. Expanding the two finite sums removes
the first node and replaces $w_\nu$, for $\nu\ge2$, by
$w_\nu(z_\nu-z_1)$. For $0\le n<m-1$ we have $|T(n)|\le2A$
by the triangle and product rules. Apply the induction hypothesis
with $m-1$ nodes and cap $2A$. For $n\le N$, repeated substitution
in $S(n+1)=z_1S(n)+T(n)$ gives
$|S(n)|\le A+2nA[3(N+1)]^{m-2}$.
Since $[3(N+1)]^{m-2}\ge1$ and $1+2N\le3(N+1)$,
the desired bound follows. This proof eliminates one entry even
when the corresponding nodes coincide.
\end{proof}

\begin{corollary}[A large late term]\label{cor:late}
Suppose now that $|z_\nu|\ge1$ for every $\nu$ and $S(0)\ne0$.
For any $K\ge0$, some $K<n\le K+m$ satisfies
\[
 |S(n)|>\frac{|\re S(0)|+|\im S(0)|}
                {8[3(K+m+1)]^{m-1}} .
\]
\end{corollary}
\begin{proof}
Set $B=K+m$ and take the new nodes $z_\nu^{-1}$ with
weights $w_\nu z_\nu^B$. They lie in the unit disk by the
squared norm formula for reciprocals. Their $j$th sum is
$S(B-j)$ for $0\le j\le B$. If all the $m$ late terms lay
in the stated disk, Lemma~\ref{lem:prefix} would put $S(0)$
in the disk of radius $(|\re S(0)|+|\im S(0)|)/8$.
At least one coordinate has absolute value at least half this
sum before division by eight, which contradicts the coordinate
bound for the disk. A bounded scan tests the squared norms
to locate a witness.
\end{proof}

For a list of $m$ nonzero rational pole offsets $v_\nu$ in a disk
of radius $d$, take $z_\nu=d/(-v_\nu)$ and unit weights.
For $P_j=\sum_\nu(-v_\nu)^{-j-1}$ the scaled sum at order $n$
is $d^nP_{n-1}$. In particular,
\begin{equation}\label{eq:poleblock}
 \exists j\in[J,J+m-1]\quad
 |P_{j+2}|>
 \frac{m}{8[3(J+m+3)]^{m-1}d^{j+3}},
\end{equation}
using $K=J+2$. Formula \eqref{eq:poleblock} is a pure PA
consequence of the finite list hypothesis. It is the
index shift used by the window coefficient search. The proof
also supplies a common positive floor for the entire block:
replace $d^{j+3}$ in the denominator by
$\max(1,d)^{J+m+2}$.

\subsection{A geometric tail defeats a polynomial loss}

\begin{lemma}[Computed domination]\label{lem:geometric}
Let $0<d<R$, $m\ge1$, $a\ge1$, and $H\ge0$ be
rational parameters with $m$ a natural. Given $\rho>0$,
there is a computed natural $k_0$ such that, for every $k\ge k_0$,
\[
 H[3(k+a)]^{m-1}(d/R)^k<\rho .
\]
\end{lemma}
\begin{proof}
Set $\lambda=R/d-1>0$. For $k\ge2m$,
\[
 (R/d)^k=(1+\lambda)^k
 \ \ge\ \binom{k}{m}\lambda^m
 \ \ge\ \frac{(k/2)^m}{m!}\lambda^m .
\]
The first bound follows by binomial expansion with nonnegative terms.
For the second, each factor in $k(k-1)\cdots(k-m+1)$ is at least
$k/2$. For $k\ge a$ also $k+a\le2k$. Therefore the expression
in the lemma is at most
\[
  \frac{H\,6^{m-1}2^m m!}{k\lambda^m}.
\]
Choose an integer strictly above the quotient of the numerator
by $\rho\lambda^m$, and at least $2m$ and $a$. A ceiling in this
construction is a bounded integer division on signed rational codes.
The proof proves the inequality for every later index.
\end{proof}

In \eqref{eq:poleblock} the loss multiplying $d^{-(j+3)}$ is
an inverse polynomial in $J$ for fixed $m$. An upper bound with
base $R>d$, a uniform constant, and any fixed shift in the order
is therefore below any prescribed fraction of the pole bound
past an explicit threshold. In applying Lemma~\ref{lem:geometric}
one absorbs fixed factors such as $R^{-2}$ and $d^3$ into $H$,
and uses $J\le j$ to replace $J+m+3$ by $j+m+3$ in the upper bound.
No test for distinct nodes or for separation of their arguments is needed.

\section{Finite truncations of the Euler evaluator}
\label{sec:euler-truncations}

Write $s=a+\mathrm i h$ with $a>0$ and $h$ rational.
Let $u_j(r)$ be the complex rational value of
\code{lw-dirichlet-sample} at $(a,h,j,r,r)$, where $j\ge1$ and $r>0$.
This section derives the dyadic tail and fixed-truncation
comparisons from finite rational estimates on the sample program.

Set $H=a+|h|$, $C_s=H+2H^2$, and $A=32C_s/a^2$.
For the two inequalities involving $1+1/j$, use the same
sample program at that rational index and write its result as
$u_{1+1/j}(r)$. The inputs needed below are
\begin{align}
 |u_j(r)-u_j(r')|
    &\le (2H+32)(r+r'),
       \label{eq:sample-transport}\\
 |u_j(r)|&\le 2,\qquad
 |u_j(r)|\le \frac{32}{a^2 k^2}+3r,
       \label{eq:sample-decay}\\
 |s+(j+1)(u_{1+1/j}(r)-1)|
    &\le \frac{C_s}{j}+(H+1)(j+1)r,
       \label{eq:sample-cell-factor}\\
 (j+1)|u_{j+1}(r)-u_j(r)u_{1+1/j}(r)|
    &\le B_j r.
       \label{eq:sample-product}
\end{align}
In \eqref{eq:sample-transport}, assume $j\ge1$ and $r,r'>0$.
In \eqref{eq:sample-decay}, assume $0<r\le1$, $k\ge2$
and $j\ge2^k$. In \eqref{eq:sample-cell-factor} and
\eqref{eq:sample-product}, also assume $j\ge H$.
The positive rational $B_j$ abbreviates the product cost from
\code{lw-euler-dirichlet-product-cost}. Its formula is $(j+1)$
times the sum of two nonnegative node cost caps, $66$, and $6H$.
The node caps are $1$ plus positive exponential upper bounds
at nonnegative rational arguments. Inequalities involving absolute
values of complex rationals use the squared disk convention
from Section~\ref{sec:graphs}.

\subsection{The sample bounds}

For a rational complex $z$, put $d(z)=|\re z|+|\im z|$ and
$S_N(z)=\sum_{\nu=0}^N z^\nu/\nu!$. The native exponential
$U(z,r)$ is $S_N(z)$ at the degree chosen by \code{lw-exp-index}.
Here is the PA adequacy proof for that actual degree. For $D=d(z)$,
the starting degree is $N_0=1+\lceil2D\rceil$. The starting tail
budget is $w=2D^{N_0+1}/(N_0+1)!$. For a nonnegative rational
$w=A_0/B_0$ and positive $r=A_1/B_1$ in canonical form,
the source fuel is $f=A_0B_1$. Thus $w\le(f+1)r$.
The procedure stops at fuel zero or $w\le r$. Otherwise it
reduces the fuel by one and doubles $r$. In that branch
$w\le(f+1)r\le2fr$, and $f\ge1$. Induction on fuel shows that
if it takes $d$ steps then $w/2^d\le r$. This also justifies its
total single-valued graph by Lemma~\ref{lem:recursion}.
The tail comparison in Section~\ref{sec:finite-exp} at degree $N_0+d$
gives $|S_N(z)-U(z,r)|\le r$ for every $N\ge N_0+d$.
Factorials, powers, and finite sums have graphs by downward
natural recursion. The source ceiling loop for $v\ge0$ starts
with fuel $1+\operatorname{num}(v)\ge v$ and subtracts one from
both input and fuel until the input is nonpositive or fuel is
zero. For any natural fuel $f\ge v$, induction on fuel proves
the result lies in $[v,v+1)$. At $v=0$ it returns zero.
For $0<v\le1$ the initial fuel is positive and the next input
is nonpositive, so the result is one. For $v>1$ apply the
induction hypothesis to $v-1\le f-1$ and add one.

We will use fresh finite degrees in the following estimates.
Given finitely many rational points and a rational $\eta>0$,
choose $N\ge2$ beyond the prescribed degrees, with the
exponential tail at $N$ bounded by $\eta$ at each needed
radius. The radii include each $d(z)+d(w)$ used below. A
single maximum suffices, and the same starting-budget and
halving calculation computes $N$. The binomial formula,
proved by multiplication and induction on degree, yields
\[
 |S_N(z)S_N(w)-S_N(z+w)|\le\eta .
\]
In fact the coefficients of total degree at most $N$
coincide. The sum of upper radii of all terms of degree
$m>N$ in the product is at most
$(d(z)+d(w))^m/m!$ by the nonnegative binomial formula.
There are no terms of degree greater than $2N$.

For $x\ge0$ and rational $y$, use this product estimate for
the pairs $(\mathrm i y,-\mathrm i y)$, $(x,-x)$,
and $(-x,\mathrm i y)$. The first product is exactly the
squared norm of $S_N(\mathrm i y)$, because the two factors
are conjugate term by term. Its squared norm is at most
$1+\eta$, and hence $|S_N(\mathrm i y)|\le1+\eta$.
For the second pair, $S_N(x)\ge1+x+x^2/2$, and the
product differs from $1$ by at most $\eta$. Thus
$|S_N(-x)|\le (1+x+x^2/2)^{-1}+\eta$.
The third pair then gives, for $0<\eta\le1$,
\[
 |S_N(-x+\mathrm i y)|\le (1+x+x^2/2)^{-1}+4\eta .
\]
As $U(-x+\mathrm i y,r)$ differs from this polynomial by at most
$r$, the same estimate holds for $U$ with $r+4\eta$ in place
of $4\eta$. We use the following arithmetic inference here
and below. If a fixed rational complex point has disk bounds
$R+c\eta$ for each rational $0<\eta\le1$, with $R,c\ge0$,
then it has disk bound $R$. If its squared norm exceeds $R^2$
by $\Delta>0$, take
\[
 \delta=\min\{1/2,\Delta/(2(2R+1))\},\qquad
 \eta=\delta/(1+c).
\]
Then $(R+c\eta)^2\le (R+\delta)^2\le R^2+\Delta/2$,
by direct expansion. This supplies the contradiction using
only a rational test. In particular,
$|U(-x+\mathrm i y,r)|\le (1+x+x^2/2)^{-1}+r$.

If $\re z,\re w\le0$ and $v=d(z-w)\le1$, the nonconstant
terms of $S_N(z-w)$ have total upper radius at most $2v$.
Indeed $m!\ge 2^{m-1}$ for $m\ge1$, by induction, and the
geometric sum of the corresponding coefficient bounds is
at most $2$. Use the product estimate for $(w,z-w)$ and
the damped bound $|S_N(w)|\le1+4\eta$. They give
$|S_N(z)-S_N(w)|\le\eta+(1+4\eta)2v\le2v+9\eta$.
If $v>1$ the two damped bounds give $2+8\eta\le2v+9\eta$.
Comparison with the native exponential at two requests and
removal of $9\eta$ therefore give
\[
 |U(z,r)-U(w,r')|\le2d(z-w)+r+r'
 \qquad(\re z,\re w\le0).
\]

For rational $x\ge1$, set $\ell_x(r)=L_n(x)$ with the
native log count $n=1+\operatorname{num}(x-1)\operatorname{den}(r)$.
This is the rectangle count in Section~\ref{sec:log-rectangles}.
Both its step and its interval width are at most $r$. For $x>1$ the sample
definition is $u_x(r)=U(-s\ell_x(r),r)$, and $u_1(r)=1$.
Overlapping log intervals give
$|\ell_x(r)-\ell_x(r')|\le r+r'$. The spatial inequality
for $U$ gives \eqref{eq:sample-transport}, with the
stronger coefficient $2H+1$. The damped bound gives the
unit estimate for $r\le1$. For $j\ge2^k$, the lower sum of
the partition with endpoints $1,2,4,\ldots,2^k,j$ is at
least $k/2$, since every doubling cell contributes $1/2$.
If $j=2^k$ omit the duplicate endpoint. A common refinement
shows that the upper endpoint of the native log interval
is at least $k/2$, so $\ell_j(r)\ge k/2-1$.
For $k\ge4$ this is at least $k/4$. For $k=2,3$,
every reciprocal in the native lower sum is at least
$1/j$, giving $\ell_j(r)\ge(j-1)/j\ge3/4\ge k/4$.
Using $x=a\ell_j(r)$ in the damped bound and
$(1+x+x^2/2)^{-1}\le2/x^2$ proves
\eqref{eq:sample-decay} with error $r$ in place of $3r$.

Put $x=1+1/j$ and $\ell=\ell_x(r)$. The reciprocal bounds
on $[1,x]$ and the interval width give
$0\le\ell\le1/j$ and $\ell+r\ge 1/(j+1)$.
In particular $|1-(j+1)\ell|\le1/j+(j+1)r$.
At $z=-s\ell$ we have $|z|\le H/j\le1$.
The degree of $U(z,r)$ is at least one. Its terms of degrees
$m\ge2$ have total upper radius at most $(H/j)^2$ from
$m!\ge2^{m-1}$ and summing the geometric series from $2$.
Then $s+(j+1)(u_x(r)-1)$ differs from
$s(1-(j+1)\ell)$ by at most $(j+1)H^2/j^2\le2H^2/j$.
Adding the two estimates proves \eqref{eq:sample-cell-factor}.

For \eqref{eq:sample-product}, scale the interval $[1,x]$
and its rectangle partition by $j$. Each cell width and
endpoint are multiplied by $j$, preserving the value of each
rectangle. The image is a partition of $[j,jx]$.
Concatenate it with the partition of $[1,j]$ and use a common
refinement with that of $[1,jx]$. The sum of the first
two lower endpoints and the third lower endpoint are
within $3r$ of one another because the three widths are
at most $r$. Thus $z'=-s\ell_{jx}(r)$ and
$z+w=-s(\ell_j(r)+\ell_x(r))$ have $d(z'-(z+w))\le3Hr$,
where $z=-s\ell_j(r)$ and $w=-s\ell_x(r)$.
Use a fresh degree with the preceding product, spatial
and damped comparisons. The polynomials at $z'$ and
$z+w$ differ by at most $6Hr+9\eta$. The latter polynomial
differs from $S_N(z)S_N(w)$ by at most $\eta$.
Replacing $S_N(z)S_N(w)$ by $U(z,r)U(w,r)$ adds at most
$2r+(1+4\eta)r$, using $|U(w,r)|\le2$ and
$|S_N(z)|\le1+4\eta$. Replacing $S_N(z')$ by $U(z',r)$
adds at most $r$. The total bound is
$(6H+4)r+14\eta$, since $r\le1$.
Remove the margin and multiply by $j+1$. The result is
at most $B_j r$ by the stated product cost formula.
The proof before the multiplication gives
$|u_{jx}(r)-u_j(r)u_x(r)|\le(6H+4)r$ for arbitrary rational
$j,x>1$ and $0<r\le1$. This also holds if $j=1$ or $x=1$
because the difference then vanishes.

\subsection{The native prefix and its evaluator}

Let the ordinary prefix be
\[
 E_n(r)=(s-1)\sum_{j=1}^n u_j(r)+(n+1)u_{n+1}(r),
 \qquad P_b(r)=E_{2^b-1}(r).
\]
These are the defining Euler prefixes in
\code{euler-dirichlet-prefix.lisp} and
\code{euler-dirichlet-tail.lisp}. For a step at $j\ge1$,
\[
 E_j(r)-E_{j-1}(r)=s u_j(r)+(j+1)(u_{j+1}(r)-u_j(r)).
\]
Subtracting $u_j(r)[s+(j+1)(u_{1+1/j}(r)-1)]$ from this
step leaves $(j+1)(u_{j+1}(r)-u_j(r)u_{1+1/j}(r))$.
For $j$ in the block $2^k,\ldots,2^{k+1}-1$, bound the
first product by multiplying the smaller of the two bounds
in \eqref{eq:sample-decay} with \eqref{eq:sample-cell-factor}.
Use the decaying bound for the $C_s/j$ term and $2$ for
the term proportional to $r$. Then \eqref{eq:sample-product}
gives
\begin{equation}\label{eq:finite-euler-cell}
 |E_j(r)-E_{j-1}(r)|\le \frac{A}{k^2j}+J_j r,\qquad
 J_j=\frac{3C_s}{j}+2(H+1)(j+1)+B_j.
\end{equation}
The expression $J_j$ is the cell rounding cost in the source.

\begin{lemma}[Dyadic tail and change of request]
\label{lem:dyadic-prefix}
For $2\le b\le c$ and $H\le2^b$, define the nonnegative
rational
\[
 D_{b,c}=\sum_{k=b}^{c-1}\sum_{j=2^k}^{2^{k+1}-1}J_j .
\]
If $0<r\le1$, then
\begin{equation}\label{eq:finite-euler-tail}
 |P_c(r)-P_b(r)|\le A/(b-1)+D_{b,c}r .
\end{equation}
For $b\ge2$ and $r,r'>0$, put
$T_b=((H+1)(2^b-1)+2^b)(2H+32)$. Then
\begin{equation}\label{eq:finite-euler-precision}
 |P_b(r)-P_b(r')|\le T_b(r+r').
\end{equation}
\end{lemma}
\begin{proof}
The finite sums of the consecutive steps in
\eqref{eq:finite-euler-cell} telescope by induction on their
length. There are $2^k$ indices in each block, and every
$1/j$ there is at most $1/2^k$, so the block sum of reciprocals
is at most $1$. Summing over blocks therefore gives a main bound
$A\sum_{k=b}^{c-1}1/k^2$, and the error sum is $D_{b,c}r$.
For $k\ge2$,
$1/k^2\le 1/(k-1)-1/k$. The latter sum is
$1/(b-1)-1/(c-1)$ by induction. This proves
\eqref{eq:finite-euler-tail}, including the empty sum when $b=c$.

For \eqref{eq:finite-euler-precision}, sum
\eqref{eq:sample-transport} over $n=2^b-1$ samples.
The factor $s-1$ belongs to the disk of radius $H+1$,
and the endpoint term has weight $n+1$. The sum of the
two bounds is $((H+1)n+n+1)(2H+32)(r+r')$.
The argument uses induction for the finite sums, with the
disk addition and multiplication rules as its successor steps.
\end{proof}

For $x>0$ and $T\ge0$, the source uses the cutoffs
\[
 K(x)=2+\lceil H\rceil+\lceil 2A/x\rceil,\qquad
 Q(T,x)=\min\{1,x/(2+2T)\}.
\]
The defining formulas give $K(x)\ge2$, $K(x)\ge H$ and
$A/(K(x)-1)\le x/2$. For $b\ge K(x)$, induction gives
$2^b\ge b+1>H$, and $A/(b-1)\le x/2$.
Also $0<Q(T,x)\le1$ and $TQ(T,x)\le x/2$ by multiplying
the bound by $2+2T$. These are exactly the cutoff and request
estimates from \code{euler-dirichlet-tail-budget.lisp}.

\begin{lemma}[Fixed truncation]
\label{lem:fixed-truncation}
For positive $e$, set $d=K(e/4)$ and $r_e=Q(T_d,e/4)$.
The native evaluator returns $V(s,e)=P_d(r_e)$.
For $b\ge2$ with $H\le2^b$, and for $q>0$,
\begin{equation}\label{eq:finite-fixed-truncation}
 |V(s,e)-P_b(q)| \le e/4+A/(b-1)+T_bq .
\end{equation}
If $x>0$, $b\ge K(x)$ and $q\le Q(T_b,x)$, the right side
is at most $e/4+x$.
\end{lemma}
\begin{proof}
Take $c=\max\{b,d\}$ and a common request $0<r\le1$.
By Lemma~\ref{lem:dyadic-prefix}, the distance from
$P_d(r_e)$ to $P_c(r)$ is bounded by
$T_dr_e+A/(d-1)+(T_d+D_{d,c})r$.
The distance from $P_c(r)$ to $P_b(q)$ is at most
$A/(b-1)+T_bq+(D_{b,c}+T_b)r$.
The sum is at most $R+Mr$, where
\[
 R=e/4+A/(b-1)+T_bq,\qquad
 M=T_d+D_{d,c}+D_{b,c}+T_b.
\]
All terms in $M$ are nonnegative. Since $r$ is an arbitrary
rational in $(0,1]$, the squared disk argument in the preceding
subsection removes $Mr$. This gives
\eqref{eq:finite-fixed-truncation}. The asserted computed bound
follows from the cutoff and request inequalities.
\end{proof}

This proves the two comparisons in
\code{euler-fixed-truncation.lisp}. For each $x>0$ both values
at requests $e,e'$ compare to the prefix with $b=K(x)$ and
$q=Q(T_b,x)$. Their difference has disk radius
$(e+e')/4+2x$. Removing the margin gives the pointwise
comparison with radius $(e+e')/4$, which implies the comparison
in \code{euler-continuation.lisp}.

\subsection{Uniform continuity and the value at one}
\label{sec:native-at-one}

The pointwise bounds give a uniform comparison on a rational
region with $\re s,\re t\ge\alpha>0$ and
$d(s),d(t)\le H_*$ for a rational $H_*$. Put
\[
 A_*=\frac{32(H_*+2H_*^2)}{\alpha^2},\qquad
 N=2^b,\qquad L=(H_*+1)(N-1)+N.
\]
Take $b\ge2$ with $N\ge H_*$.
Keeping the same request $0<r\le1$ for the two prefixes,
write $P_b^s(r)$ for the prefix at $s$. For $1\le j\le N$
the log endpoint is at most $j-1$. The spatial bound for
$U$ in Section~\ref{sec:euler-truncations} therefore gives
a sample difference at most $2N d(s-t)+2r$. Each sample
has disk radius $2$. In the coefficient difference for the
prefix, use
$(s-1)u_j^s-(t-1)u_j^t=(s-1)(u_j^s-u_j^t)+(s-t)u_j^t$.
Including the endpoint term yields
\[
 |P_b^s(r)-P_b^t(r)| \le [2NL+2(N-1)]d(s-t)+2Lr .
\]
Compare the two evaluator outputs to these prefixes using
Lemma~\ref{lem:fixed-truncation}, with $A\le A_*$ and
$T_b\le L(2H_*+32)$ at either point. The squared disk
argument removes the resulting multiple of $r$. Thus,
with $M_b=2NL+2(N-1)$,
\begin{equation}\label{eq:euler-region-spatial}
 |V(s,e)-V(t,e')|
 \le (e+e')/4 + 2A_* /(b-1)+M_b d(s-t).
\end{equation}
For any rational $\delta>0$, the choice
$b=2+\lceil H_*\rceil+\lceil 4A_*/\delta\rceil$
gives $2A_* /(b-1)\le\delta/2$. Then the last two terms
in \eqref{eq:euler-region-spatial} sum to at most $\delta$
whenever $d(s-t)\le\delta/(2M_b)$.

\begin{lemma}[Native normalization at one]\label{lem:native-at-one}
In $\PA$, for every rational $e>0$ the evaluator satisfies
$|V(1,e)-1|\le e/8$.
\end{lemma}
\begin{proof}
For rational $x>1$, take a finer log rectangle request
$r'>0$, with step $\sigma\le r'$ and $m$ cells. The right
endpoints are $y_i=1+i\sigma$ for $0\le i\le m$.
Put $\delta_i=\sigma/y_i$ for $1\le i\le m$. Then
\[
 \prod_{i=1}^m(1-\delta_i)=\prod_{i=1}^m y_{i-1}/y_i=1/x,
 \qquad
 \sum_{i=1}^m \delta_i^2\le\sigma(x-1).
\]
The first identity uses cancellation at each step. The
second follows from $\delta_i\le\sigma$ and
$\sum_i\delta_i\le m\sigma=x-1$.
For each $i$ use the finite exponential product estimate
at the negative partial sum of the preceding $\delta_j$
and at $-\delta_i$, with error $\eta\le1$. All the needed
arguments and sums of radii are bounded by $2(x-1)$, so
one degree works at every step. By the damped estimate
the partial sum polynomial has radius at most $1+4\eta\le5$.
The bound on its terms of degree at least two at
$-\delta_i$ gives $|S_N(-\delta_i)-(1-\delta_i)|\le\delta_i^2$.
In the comparison to the product of $(1-\delta_i)$ the
previous error has factor $1-\delta_i\in(0,1)$, and the new
error is at most $5\delta_i^2+\eta$. Induction on $i$ gives
\[
 |S_N(-\ell_x(r'))-1/x|\le5\sigma(x-1)+m\eta.
\]
For ordinary request $r$ take $N$ also beyond the degree
of $U(-\ell_x(r),r)$. The comparison with
$S_N(-\ell_x(r'))$ has error at most $r+2(r+r')+9\eta$,
by the common log interval and finite spatial estimates.
Consequently the total error is at most
$3r+[2+5(x-1)]r'+(m+9)\eta$. For a prescribed margin
$0<g\le1$, first take $r'=g/(2(3+5(x-1)))$. Then, after
computing $m$, take $\eta=g/(2(m+10))$.
The total error is at most $3r+g$. Removing the margin
shows $|U(-\ell_x(r),r)-1/x|\le3r$.

At $s=1$ the first summand in $E_n(r)$ vanishes, so the
last bound at $x=n+1$ gives $|P_b(r)-1|\le3\cdot 2^b r$.
The native degree index and point request are
$b=K(e/4)$ and $r=Q(T_b,e/4)$. Here $T_b\ge3\cdot2^b$
and $T_b r\le e/8$, proving the claim.
\end{proof}

\subsection{A finite real-axis clearance bound}
\label{sec:real-axis}

For $a>0$ write $g(a)=1-(1+a/2)^{-1}$. This rational is
positive, is at most $1$, and is increasing by reciprocal
order. It is the first-term gap in \code{real-dirichlet-order.lisp}.

\begin{lemma}[Real-axis clearance]\label{lem:real-axis}
For rational $a>0$ and $e>0$,
\[
        |V(a,e)|+e/4\ge g(a)/4 .
\]
In particular for $a\ge\alpha>0$ the request $e\le g(\alpha)/64$
gives $|V(a,e)|\ge g(\alpha)/8$.
\end{lemma}
\begin{proof}
All the sample polynomials are real when the height is zero.
For $0\le x\le y$, use the finite product estimates at $(x,-x)$
and $(y,-y)$. At a common sufficiently large degree,
$1\le S_N(x)\le S_N(y)$ and each $S_N(-v)$, for $v=x,y$,
differs from $1/S_N(v)$ by at most $\eta$.
The tail bound for $U$ gives $U(-x,r)\ge U(-y,r)-2r-2\eta$,
and gives $U(-y,r)\ge-r-\eta$.
Removing the rational margins gives the two comparisons with
errors $2r$ and $r$.
For integer $j\ge2$, the common log partition from $1$ to
$j+1$ with a breakpoint at $j$ gives
$\ell_{j+1}(r)+r\ge\ell_j(r)$.
If $\ell_j(r)\le\ell_{j+1}(r)$, the order comparison just
proved gives $u_j(r)-u_{j+1}(r)\ge-2r$.
In the other case the spatial comparison gives
$u_j(r)-u_{j+1}(r)\ge-(2a+2)r$.
At $j=1$, $u_1(r)=1$ and $\ell_2(r)\ge1/2$, so the
damped bound gives $u_1(r)-u_2(r)\ge g(a)-r$.

For $n\ge1$ the odd alternating prefix
$Y_n=\sum_{j=1}^{2n+1}(-1)^{j+1}u_j(r)$ therefore satisfies
\[
 Y_n\ge g(a)-F_n r,\qquad F_n=2+(n-1)(2a+2).
\]
The estimate sums the first pair bound, the $n-1$ further
pair bounds, and $u_{2n+1}(r)\ge-r$.
Let $M=2u_2(r)$ and $d_0=|a-1|$.
The native sample comparison at $a=1$ in the preceding
subsection gives $|u_2^1(r)-1/2|\le3r$. Since
$\ell_2(r)\le1$, the spatial comparison of $u_2^a(r)$
and $u_2^1(r)$ gives a difference of at most $2d_0+2r$.
Consequently $|1-M|\le4d_0+10r$, and $|M|\le4$ for $r\le1$.

Separate the even indices in the finite sum for $E_{2n+1}(r)$.
Subtracting $M E_n(r)$ and $(a-1)Y_n$ leaves
\begin{align*}
 &2(a-1)\sum_{j=1}^n (u_{2j}(r)-u_2(r)u_j(r))\\
 &\qquad{}+2(n+1)(u_{2(n+1)}(r)-u_2(r)u_{n+1}(r)).
\end{align*}
The identity follows from the even and odd partition of the
prefix, which holds by induction on $n$, and from the
two endpoint terms. The general sample product estimate
bounds this residual by $C_n r$ for
$C_n=2(d_0 n+n+1)(6a+4)$.

Choose any positive rational $\tau$, take $b=K(\tau)$ and
$n=2^b-1$. The two ordinary prefixes above are $P_b(r)$ and
$P_{b+1}(r)$. For any
$0<r\le\min\{Q(T_b,\tau),Q(T_{b+1},\tau)\}$,
they each compare to $V(a,e')$ with error at most $e'/4+\tau$.
The triangle inequality applied to the residual gives
\[
 |(1-M)V(a,e')-(a-1)Y_n|\le C_n r+5(e'/4+\tau).
\]
Together with the lower bound on $Y_n$ and the bound on
$|1-M|$, it follows that
\[
 d_0g(a)\le
  4d_0|V(a,e')|+5e'/4+5\tau+
  [d_0F_n+10|V(a,e')|+C_n]r .
\]
For fixed $\tau$ the coefficient of $r$ is rational and fixed.
If the inequality with the last term omitted fails by a
positive rational, choose a positive $r$ below the common
request and small enough that this last term is less than
the gap. This removes the term. The resulting inequality
holds for every $\tau>0$, so the same rational comparison
removes $5\tau$. Comparing $V(a,e')$ with $V(a,e)$ at
radius $(e+e')/4$ now gives
\[
 d_0g(a)\le 4d_0(|V(a,e)|+e/4)+(d_0+5/4)e'.
\]
Remove $e'$ by the same comparison and divide by $d_0$
when $a\ne1$. At $a=1$ the claim follows from
Lemma~\ref{lem:native-at-one} and $g(1)\le1$.
\end{proof}

The request in \code{euler-real-axis-clearance.lisp} is at most
$\rho g(\alpha)/64$ and its floor
is at most $\rho g(\alpha)/32$, where the source anchor
radius $\rho$ is at most $1/4$ by its definition. Positive
requests in the theorem imply $\rho>0$, and the bound in
Lemma~\ref{lem:real-axis} then exceeds the required floor.
Section~\ref{sec:native-uniform-polynomial} supplies the graph
for this anchor radius and proves near-real clearance at the
source's computed width.

\subsection{Clearance to the right of one}
\label{sec:right-clearance}

Let $s=a+\mathrm i h$ with $a>1$, and set $v=a-1$,
$D=32/v^2$, and $B=13+2D$. For the samples at $s$ use
superscripts for rates when needed. Their exponential arguments
at rate $s$ are the sums of the arguments for rates $1$ and
$s-1$, at the same $\ell_j(r)$.
By the finite exponential product estimate and the damped bound,
\[
 |u_j^{s}(r)-u_j^{1}(r)u_j^{s-1}(r)|\le4r \qquad(0<r\le1).
\]
Here the approximation of the sum argument contributes $r$.
The product of the two native approximants differs from the
product of polynomials by at most $(2+1+4\eta)r$. The
polynomial product contributes an extra $\eta$. The squared
disk comparison removes $(1+4r)\eta$. The identity also holds
at $j=1$ with zero error.
By the rational reciprocal comparison in the proof of
Lemma~\ref{lem:native-at-one}, $|u_j^1(r)-1/j|\le3r$.
For $2^k\le j<2^{k+1}$ with $k\ge2$ the improved decay
bound for rate $s-1$ is $D/k^2+r$. Its unit bound is $2$.
Applying the decay to $1/j$ and the unit bound to the
$3r$ error gives
\begin{equation}\label{eq:normalized-node-tail}
   |u_j^s(r)|\le \frac{D}{k^2 j}+11r .
\end{equation}

For a finite prefix through $n$ use scalar radii $2$ for
$j\le3$ and the bound in \eqref{eq:normalized-node-tail}
for $j\ge4$. Let $\mathcal B$ be their sum. Let
$\mathcal T$ be the sum of those radii with $j\ge2^k$,
where $k\ge2$. Divide the indices into dyadic blocks.
As in Lemma~\ref{lem:dyadic-prefix}, each complete or partial
block has reciprocal mass at most one. The telescoping estimate
for reciprocal squares gives
\[
 \mathcal B\le6+D+11nr,\qquad
 \mathcal T\le D/(k-1)+11nr.
\]
Each sum is finite, with the empty sum equal to zero.

Use the arithmetic Möbius coefficient $\mu(j)$ for positive
integer $j$. Its defining computation factors by trial division,
returns zero if a prime repeats, and otherwise returns $(-1)^m$
for $m$ prime factors. Factorization exists by strong induction
on $j$ using the least divisor greater than one, which is prime.
Euclid's lemma follows by Euclidean division and the Bézout
identity from Section~\ref{sec:graphs}. Applying it to a factor
of a product and canceling one prime at a time proves uniqueness
up to order. Bounded divisor tests and finite lists of factors
give the computation graph. For $j>1$ the divisors with
nonzero $\mu$ correspond exactly to subsets of the distinct
prime factors of $j$. Pairing each subset with the subset that
changes membership of the first prime cancels their contributions.
Thus $\sum_{d\mid j}\mu(d)=0$ for $j>1$, and the sum equals one
at $j=1$. These arguments use only ordinary induction.

\begin{lemma}[Right-region clearance]\label{lem:right-clearance}
In $\PA$, if $a>1$ and $0<e\le v/(32B)$, then
$V(a+\mathrm i h,e)$ lies outside the disk of radius $v/(32B)$.
\end{lemma}
\begin{proof}
Put $f=v/(32B)$. Choose $k\ge2$ so that
$D/(k-1)\le1/(32B)$, for example
$k=2+\lceil32BD\rceil$.
Take
\[
 b=\max\{2k,\ K(f),\ 2+\lceil 2D/f\rceil\},\qquad
 n=2^b-1 .
\]
Then $D/b^2\le f/2$. Use a positive rational request $r$
at most each of
\[
 1,\quad Q(T_b,f),\quad
 \frac{1}{32B(1+11n)},\quad
 \frac{1}{8(1+n^2(6H+4))},\quad
 \frac{f}{2(1+11(n+1))}.
\]
It follows that $\mathcal B\le B$, $\mathcal T\le1/(16B)$,
and the endpoint term $(n+1)u_{n+1}^s(r)$ has disk radius
at most $f$ by \eqref{eq:normalized-node-tail} at block $b$.

Write $Z=\sum_{j=1}^n u_j^s(r)$ and
$J=\sum_{j=1}^n\mu(j)u_j^s(r)$. The radius of $J$ is at
most $\mathcal B$. In the product $ZJ$, restrict first to
pairs $(i,j)$ with $ij\le n$. Replacing $u_i^s(r)u_j^s(r)$
by $u_{ij}^s(r)$ costs at most $n^2(6H+4)r\le1/8$ by
the general product estimate. Grouping these terms by $ij$
and using the divisor identity gives exactly $1$.
For each remaining pair one index is at least $2^k$,
since $(2^k-1)^2\le n$. Their products have combined
radius at most $2\mathcal B\mathcal T\le1/8$.
The coefficient $\mu(j)$ has absolute value at most one.
Thus $|ZJ-1|\le1/4$. If $|Z|\le1/(2B)$, then
$|ZJ|\le1/2$, which would place $1$ in the disk of radius
$3/4$. This proves $Z$ lies outside that disk of radius $1/(2B)$.

If $|V(s,e)|\le f$, the computed fixed-truncation bound
gives $|P_b(r)|\le f+e/4+f$. Subtract the endpoint term
to get $|(s-1)Z|\le(13/4)f$. Yet the squared norm of
$s-1$ is at least $v^2$, and the squared norm of $Z$ exceeds
$1/(4B^2)$. Their product exceeds the square of $v/(2B)=16f$,
by the squared norm product rule. This is the contradiction.
\end{proof}

The source floor \code{lw-euler-right-floor} is exactly
$v/(32B)$. It increases with $a$ because $v$ increases and
$B=13+64/v^2$ decreases. For $a\ge1+g$ with rational $g>0$,
the lemma therefore gives uniform clearance at the request
and floor computed at $1+g$. This discharges the right-region
bound required in Section~\ref{sec:contour}.

\subsection{The native uniform polynomial and near-real radius}
\label{sec:native-uniform-polynomial}

Fix a rational complex center $s_0=a_0+\mathrm i h_0$ and a
rational radius $0\le R<a_0$. For $|z|\le R$ put $s=s_0+z$.
The disk gives $\re s\ge a_0-R$ and $d(s)\le H_*$, where
$H_*=a_0+|h_0|+2R$. Thus the surrogate point with real part
$a_0-R$ and height $|h_0|+3R$ bounds the tail cap, cutoff and
transport cost of $s$ by their defining formulas.

We first record the finite approximation for the exponential
mixtures used by \code{lw-euler-uniform-polynomial}. For a list
of $m$ nodes $(\ell_i,w_i)$ with $\ell_i\ge0$ rational and
$w_i$ complex rational, the coefficient at degree $j$ is
$\sum_i w_i\ell_i^j/j!$ by the source recurrence. By the
distributive laws and finite induction, its sum through degree
$N$ evaluated at $v$ equals $\sum_i w_i S_N(\ell_i v)$.
Replacing $v$ by $-s_0-z$ is the \code{lw-polynomial-affine}
operation. Expanding each factor by the binomial formula
proves equality of the two evaluations. The coefficient lists
use natural length or degree recursions.

At mixture error $\varepsilon>0$, for node $i$ the native
degree is at least the degree of the scalar tail at radius
$\ell_i H_*$ and request
\[
        t_i=\frac{\varepsilon}{2(m+1)(1+d(w_i))}.
\]
The source takes the maximum over the nodes and the supplied
natural target degree. For $|v|\le H_*$, compare
$S_N(\ell_i v)$ and $U(\ell_i v,r)$ to a common polynomial
of degree at least $N$ and the native degree of $U$.
The tail formula gives a difference of at most $t_i+r$.
Multiplying by $w_i$ and summing shows the difference between
the mixture polynomial and the sum of native samples has
radius at most
\[
        \varepsilon+r\sum_i d(w_i).
\]
Indeed $d(w_i)t_i\le\varepsilon/(m+1)$ and one can add
these nonnegative bounds by induction. The empty mixture
has zero polynomial and zero sample sum.

\begin{lemma}[Uniform polynomial comparison]\label{lem:native-poly}
Let $P_\varepsilon(z)$ be the native uniform polynomial at
center $s_0$, radius $R$, any natural target degree, and
positive model request $\varepsilon$. For every rational
$z$ in the $R$-disk and $e>0$, $\PA$ proves
\[
      |P_\varepsilon(z)-V(s_0+z,e)|\le\varepsilon+e/4 .
\]
\end{lemma}
\begin{proof}
In the source, $b$ is $K(\varepsilon/2)$ at the surrogate
point, $n=2^b-1$, and the common point request $r$ is the
minimum of $Q(T_b,\varepsilon/2)$ at the surrogate and
$Q(C,\varepsilon/4)$, where
$C=(H_*+1)(n-1)+(n+1)$. This positive rational satisfies
$Cr\le\varepsilon/8$. The ordinary mixture has nodes
$(\ell_j(r),1)$ for $2\le j\le n$ with error
$\varepsilon/(8(H_*+1))$. The boundary mixture is the
single node $(\ell_{n+1}(r),n+1)$ with error $\varepsilon/8$.
The source adds $1$ to the ordinary mixture, multiplies
it by $s_0+z-1$, and adds the boundary mixture. These
are the operations in \code{lw-euler-prefix-polynomial}.
Their evaluation identities follow by distributivity and
induction on the coefficient lists.

The error against $P_b(r)$ at $s$ is therefore at most
\[
 (H_*+1)\frac{\varepsilon}{8(H_*+1)}
          +\frac{\varepsilon}{8}+Cr\le3\varepsilon/8 .
\]
By the defining surrogate bounds, $b\ge K(\varepsilon/2)$
at $s$ and $r\le Q(T_b,\varepsilon/2)$ at $s$.
Lemma~\ref{lem:fixed-truncation} bounds the remaining
difference by $\varepsilon/2+e/4$. Adding the estimates
proves the lemma with the native cutoffs and requests.
\end{proof}

For any two model requests $\varepsilon,\varepsilon'>0$
and two target degrees at the same center and radius,
the polynomial difference has disk radius
$\varepsilon+\varepsilon'+e/2$ at every rational point of
the disk, by comparing both models to $V(s_0+z,e)$.
The arbitrary positive $e$ can be removed by the squared
disk comparison. For $R>0$, Lemma~\ref{lem:finite-cauchy}
then bounds the difference of the degree-$j$ coefficients
after multiplication by $R^j$ by $\varepsilon+\varepsilon'$,
with absent coefficients treated as zero. This is the
normalized coefficient convention in
\code{euler-coefficient-cauchy.lisp}.

We apply the comparison to the source's near-real radius.
For a real center $a_0$ and outer disk radius $R_*$ satisfying
$0\le R<R_*<a_0$, put $L=(R_*-R)/2$ and $\alpha=a_0-R_*$.
Use the source coefficient cap $C_*$ on the outer disk.
It is $2+\sum_j d(c_j)R_*^j$, where the $c_j$ are the
coefficients of the outer polynomial with request $1$ and
target degree $0$. This expression has a total rational
graph by Lemma~\ref{lem:native-poly} and the coefficient
construction, and $C_*\ge2$. That polynomial has disk
radius $C_*$ at every point in the outer disk, since
each term has radius at most $d(c_j)R_*^j$.
By Lemma~\ref{lem:native-poly}, the native evaluator
has radius at most $C_*+2$ there for request $1$.

Let $t$ be real with $|t-a_0|\le R$. Every point in
the $L$-disk about $t$ lies inside the outer disk.
For any $0<\eta\le1$, take the native local polynomial
$A_\eta$ at $t$ and $L$ with model request $\eta$.
Its values on that disk have radius at most $C_*+4$,
by comparison with the native value at request $1$.
Lemma~\ref{lem:finite-cauchy} gives coefficient radii
$(C_*+4)/L^j$. On $|z|\le L/2$, the nonconstant terms
then have total radius at most $2(C_*+4)|z|/L$,
by the finite geometric sum. Compare the two evaluations
at $z$ and $0$ to the native evaluator at request $e$.
Removing the arbitrary margin $2\eta$ gives
\begin{equation}\label{eq:near-real-native-transfer}
 |V(t+z,e)-V(t,e)|\le e/2+S|z|,
 \qquad S=4(C_*+5)/L ,
\end{equation}
with $|z|$ on the right interpreted as any rational
nonnegative upper radius at most $L/2$.

Let $\rho$ be the source anchor radius for the point $1$.
In its defining \code{lw-euler-neighborhood-radius},
the slope has a total rational graph with nonnegative value
by Lemma~\ref{lem:native-poly}: it is $2$ divided by the
positive outer radius, times the maximum of the nonconstant
rescaled coefficient sizes. The pointwise estimate at $1$
in Lemma~\ref{lem:native-at-one} gives a positive
\code{lw-euler-value-clearance-size} at request $1/64$,
since the norm of the value is at least $1-1/512$.
The minimum of half the radius and the positive request
$Q$ is therefore positive. The source choice of outer
radius $1/2$ gives $0<\rho\le1/4$.

The source real-axis floor is $F=\min\{1/8,\rho g(\alpha)/32\}$.
On its real-axis request
$e\le\min\{1/64,\rho g(\alpha)/64\}$,
Lemma~\ref{lem:real-axis} gives a real-axis value of norm
at least $F$ at $t\ge\alpha$.
The source near-real width is $w=\min\{L/2,Q(S,F/64)\}$.
This is positive and satisfies $Sw\le F/128$.
The near-real request additionally requires $e\le F/32$.
Then \eqref{eq:near-real-native-transfer} gives a disk radius
at most $3F/128$ for the change from $t$ to $t+\mathrm i h$
when $|h|\le w$. If the value at the latter point had radius
$F/8$, the value at $t$ would have radius at most
$F/8+3F/128<F$. This proves the exclusion with the exact
width and request of \code{euler-near-real-clearance.lisp}.

\section{Rational approximate factorization}
\label{sec:native-factorization}

This section proves that the native root search succeeds and
that its recursive factor list meets the requested error on a
disk. The root search uses squared norms on a computed rational
grid. An auxiliary finer grid proves that its test can succeed.

\subsection{A direction with negative real part}

Write $d(z)=|\re z|+|\im z|$. We use
$|z|^2\le d(z)^2\le 2|z|^2$, where $|z|^2$ denotes the
rational squared norm. We also use the rational inverse
$z^{-1}=\bar z/|z|^2$ when $z\ne0$. Its squared norm is
$1/|z|^2$.

\begin{lemma}[A finite direction choice]\label{lem:negative-direction}
For a nonzero polynomial $G(X)=\sum_{k=1}^n c_k X^k$
with coefficients in $\QQi$ and $n\ge1$, set $H=\sum_k d(c_k)$.
There is a computed rational unit point $u$ such that
\[
    \re G(u)\le -H/(16n).
\]
\end{lemma}
\begin{proof}
Use the rational unit point $\omega=(3+4\mathrm i)/5$ from
Lemma~\ref{lem:finite-cauchy}. For each nonzero signed exponent
$k$, the average of $u^k$ over the list
$u=1,\omega,\ldots,\omega^{N-1}$ has disk radius at most
$2d((1-\omega^k)^{-1})/N$, by the finite geometric sum.
Put $f(u)=\re G(u)$. At unit points $f$ is the Laurent
polynomial $(G(u)+\sum_{k=1}^n\bar c_k u^{-k})/2$.
Its constant coefficient is zero. The constant coefficient
of $f^2$ is $\tfrac12\sum_k |c_k|^2$, by expanding the finite
product. It is at least $H^2/(4n)$ since
$d(c_k)^2\le2|c_k|^2$ and $(\sum_k d(c_k))^2\le n\sum_k d(c_k)^2$.
The latter inequality follows by summing the nonnegative
$(d(c_i)-d(c_j))^2$ over all pairs.

Compute the Laurent coefficients of $f$ and $f^2$ by finite
convolution. For either polynomial the bound for the average
of its nonconstant part is $B/N$, where $B$ is the sum of the
coefficient radii times the corresponding geometric sum bounds
with the factor $1/N$ omitted. Take $N$ large enough that the
two errors are respectively at most $H/(32n)$ and $H^2/(32n)$.
Such an $N$ is computed by rational division and the ceiling
operation. Hence the average of $f^2$ is at least $7H^2/(32n)$.
If there were no desired point, every $f$ would lie between
$-a$ and $H$ for $a=H/(16n)$. Expanding $(H-f)(f+a)\ge0$
gives $f^2\le(H-a)f+Ha$. Its average would be at most
$3H^2/(32n)$, a contradiction. Comparing the real parts along
this finite orbit gives the point. All identities and bounds
use finite sums and the rational disk rules.
\end{proof}

\subsection{Completeness of the native polynomial grid}

\begin{lemma}[Computed approximate root]\label{lem:native-poly-root}
Let $P=\sum_{j=0}^n a_j X^j$ have rational complex
coefficients with $n\ge1$ and $a_n\ne0$. For rational
$\varepsilon>0$, both versions of the native global root
search succeed and return the same point $v$, with
\[
     |P(v)|\le\varepsilon,\qquad |v|\le 2R ,
\]
where $R$ is \code{lw-polynomial-leading-radius}.
\end{lemma}
\begin{proof}
Put $\lambda=|a_n|^2/(1+|a_n|^2)>0$ and
$M=\sum_{j<n}d(a_j)$. The source takes
$R=1+16M2^n/\lambda$ and evaluates $Q(X)=P(RX)$ on
the rational square $[-1,1]^2$. The inequality
$\lambda^2\le |a_n|^2$ follows from
$|a_n|^2\le(1+|a_n|^2)^2$.
For a point $x$ in the square having a coordinate of absolute
value at least $1/2$, the rational inverse of $Rx$ has disk
radius $2/R$. Hence each factor $(Rx)^{j-n}$ with $j<n$
has disk radius at most $2^n/R$, using $R\ge1$.
Dividing the polynomial by its leading term at $x$ gives
\[
  \left|\frac{Q(x)}{a_n(Rx)^n}-1\right|
       \le \frac{M2^n}{\lambda R}\le 1/16.
\]
The triangle inequality, or its squared form, gives the
lower bound $15/16$ for the norm of the ratio.
Since the leading term has norm at least $\lambda R^n/2^n$,
$Q(x)$ has norm at least $\lambda R^n/2^{n+1}>0$.
This bound is at least $8M$, whereas $|Q(0)|\le M$.
It follows also when $M=0$ that any point with
$|Q(x)|^2\le |Q(0)|^2$ has both coordinates strictly
between $-1/2$ and~$1/2$.

Write $b_j=a_j R^j$. Let
\[
 U=1+\sum_{j=0}^n d(b_j)4^j,\qquad
 L=1+\sum_{j=1}^n j\,d(b_j)4^{j-1}.
\]
On the square $|Q(x)|\le U$ and
$|Q(x)-Q(y)|\le Ld(x-y)$. Indeed
$x^j-y^j=(x-y)\sum_{\ell=0}^{j-1}x^\ell y^{j-1-\ell}$
and both arguments lie in the disk of radius $4$.
The squared value changes by at most $2ULd(x-y)$,
using the real part of $(Q(x)-Q(y))\overline{(Q(x)+Q(y))}$.
Expand $Q(q+w)=Q(q)+\sum_{k=1}^n A_k(q)w^k$
at $q$ in the square by the binomial formula.
The sum of the $d(A_k(q))$ for $k\ge1$ is at most $U$,
by adding the binomial terms and using $d(q)\le2$.
The coefficient $A_n(q)$ is $b_n$.

Set $\varepsilon_0=\varepsilon/2$ and
$t=\min\{1/4,\varepsilon_0/(128nU)\}$. If
$|Q(q)|^2>\varepsilon_0^2$, apply Lemma~\ref{lem:negative-direction}
to $G(u)=\overline{Q(q)}\sum_{k=1}^n A_k(q)t^ku^k$.
Put $S=\sum_k d(A_k(q))t^k>0$. The coefficient sum $H$ of
$G$ satisfies $H\ge\varepsilon_0 S/2$. For each term this
follows by the squared product identity and
$|A_k(q)|^2\ge d(A_k(q))^2/4$.
The perturbation $W=Q(q+tu)-Q(q)$ has radius $S$ at a unit
point and $S\le tU\le\varepsilon_0/(128n)$, so
$|W|^2\le H/(64n)$. At the direction supplied by the lemma,
\[
 \begin{split}
 |Q(q+tu)|^2-|Q(q)|^2
   &=2\re G(u)+|W|^2\\
   &\le -H/(16n)\le -\Delta,
 \qquad \Delta=\frac{\varepsilon_0\,d(b_n)t^n}{32n}>0 .
 \end{split}
\]
Here $\Delta$ depends on $P$ and $\varepsilon$, and the
lower bound for $H$ uses $S\ge d(b_n)t^n$.

Choose an auxiliary square grid with
$2(1+\lceil8UL/\Delta\rceil)$ subdivisions of a side.
It contains zero and has step $\sigma\le\Delta/(8UL)$.
By a finite scan choose $q$ of least squared value on this
grid. The preceding leading-term comparison puts $q$ inside
the square with coordinates of absolute value less than $1/2$.
Thus $q+tu$ from the descent lies in $[-1,1]^2$.
Rounding each coordinate to a grid coordinate within $\sigma$
gives a squared value at most $|Q(q)|^2-\Delta+4UL\sigma$.
This contradicts minimality when $|Q(q)|^2>\varepsilon_0^2$,
so the auxiliary grid provides a point of radius at most
$\varepsilon_0$ in value.

The native grid \code{lw-polynomial-global-root-grid} has
$N=4(F+J)$ subdivisions of a side. Here $F$ is the unit-square
mesh factor at the source outer step, and
$J=1+\lceil8(1+K)/\varepsilon\rceil$, with $K$ the native
polynomial Lipschitz cap for $Q$ at radius $7$.
In the unit-square mesh factor, the zero branch returns one
and the other branch adds one to a natural from the ceiling loop.
Thus $F$ is positive and $N\ge J$.
Unrolling the native Lipschitz recurrence gives
$K=\sum_{j=1}^n j\,d(b_j)7^{j-1}$. It bounds value
differences by $K d(x-y)$ on the square by the difference
of powers used above. The mesh step satisfies
$2(2/N)K\le\varepsilon/2$ by the formula for $J$.
Round the auxiliary point to this native grid within $2/N$
in each coordinate. Its value has radius at most
$\varepsilon_0+\varepsilon/2=\varepsilon$.

Both searches select the first passing entry by row and
column. The row scan stops at the first squared norm at most
$\varepsilon^2$. Induction on the remaining columns proves
its result is the least passing column or reports failure
when all those columns fail. Induction on the remaining rows
then gives the first passing address. In the streaming search
the next row is evaluated only when the preceding row fails,
and its values are the same coefficient evaluations as in
the full grid. Thus the returned address agrees. Rescaling
by $R$ gives the source point $v$. Each coordinate before
rescaling has absolute value at most one, giving the radius
$2R$. All scans, polynomial evaluations, and grid roundings
have total graphs using bounded loops and the arithmetic
quotient and remainder.
\end{proof}

\subsection{The native factor list}

\begin{lemma}[Approximate factorization on a disk]\label{lem:native-factors}
For a nonzero polynomial $P$ after trailing-zero trimming,
and rational $B\ge0$ and $\varepsilon>0$, the native
normalized factor list has length $\deg P$ and satisfies
\[
     \left|P(x)-a_n\prod_{j=1}^n (x-v_j)\right|
           \le\varepsilon\qquad (|x|\le B).
\]
For the zero polynomial the native factor polynomial is zero.
\end{lemma}
\begin{proof}
Trailing-zero removal preserves evaluations by induction on
the list. If the list is nonempty after trimming, its leading
coefficient is nonzero. At degree zero the factor list is
empty and the product has value $a_0$.
For degree $n>0$ take the native root search at request
$\varepsilon/2$ and call its point $v$. Let $D$ be the source
divided-difference polynomial. For coefficient index
$0\le j<n$, its coefficient is
$\sum_{\ell=j+1}^n a_\ell v^{\ell-j-1}$, by the source
recursion on the coefficient tail. In particular it has
degree $n-1$ and leading coefficient $a_n$.
Multiplying by $x-v$ and canceling consecutive terms gives
$P(x)=P(v)+(x-v)D(x)$. The next request in
\code{lw-polynomial-factor-next-error} is
$\varepsilon/[2(1+B+2R)]>0$, with $R$ for this step.
Apply induction on the degree to the recursive list of $D$.
For $|x|\le B$ the new error from its product is at most
$(B+2R)\varepsilon/[2(1+B+2R)]\le\varepsilon/2$,
since $|v|\le2R$ by Lemma~\ref{lem:native-poly-root}.
Adding $|P(v)|\le\varepsilon/2$ proves the asserted bound.
The list length increases by one from the recursive length.
The degree decrease also establishes totality by
Lemma~\ref{lem:recursion}.
For the empty trimmed list, the leading coefficient defaults
to zero and the algorithm returns the empty root list,
so the factor polynomial is zero. For all lists the evaluated
factor polynomial equals the displayed product by induction
on the linear multiplications. This proves the error statement
for \code{lw-polynomial-disk-factor-polynomial}.
\end{proof}

\section{Constructing the checked contour}
\label{sec:construct-cover}

\subsection{A logarithmic bound for reflected factors}

Use the positive rational logarithm intervals from
Section~\ref{sec:log-rectangles}, with the source signed convention
for values below one. In this subsection a logarithm expression
denotes a finite sum of these interval names with specified rational
coefficients. To assert that such an expression is at most a
rational means that every lower endpoint is at most that rational.
The upper bound then holds for the upper endpoints up to the
sum of the widths. The analogous meaning is used for a lower
bound. Bounds established at arbitrarily fine common requests
imply bounds in this convention by pairwise overlap and rational
margin removal. For a finite sum the common request can be
divided by one plus the sum of the absolute coefficients.

Here are the arithmetic facts about these interval names used
below. Orient the reciprocal rectangles from $1$ to any positive
rational $x$, using minus the reversed sum if $x<1$.
Positive scaling preserves every rectangle, since its width and
its sample point are multiplied by the same rational.
For $x<1$, scaling the reversed interval by $1/x$ proves
agreement of the signed convention with
\code{lw-signed-log-bounds}. Concatenate the interval from $1$
to $x$ with the scaled interval from $x$ to $xy$. Refine at all
their breakpoints and those for the interval from $1$ to $xy$.
The overlapping reversed cells cancel and the rest are the cells
from $1$ to $xy$. The lower and upper sums on each partition
enclose the corresponding refined sum. This proves the
product relation $\log(xy)=\log x+\log y$ with arbitrary
rational accuracy. The same partition comparison gives
monotonicity. A single coarse cell on $[1,x]$ or $[1,1/x]$ gives
$-1/x\le\log x\le x$ for $x>0$, and $\log 2\ge1/2$.
For any rational $\xi$ with $|\xi|\le1/2$, the reciprocal
on the interval from $1$ to $1+\xi$ differs from one by at
most $2|\xi|$. The oriented rectangles give
$|\log(1+\xi)-\xi|\le 2\xi^2$. Every use of refinement here
is a finite list refinement with rational endpoints.

For $|w|\le q<1$, $q$ rational and nonnegative, let
$Y(t)=|1-tw|^2$ for $0\le t\le1$, and set
$D(t)=2\re(w/(1-tw))$. The exact identity
\[
 \frac{Y(t+h)-Y(t)}{Y(t)}
     =-hD(t)+h^2 |w|^2/Y(t)
\]
holds for each subdivision step $h$. Since
$Y(t)\ge(1-q)^2$, the absolute value of the left side
is at most $h K_1$, with $K_1=2/(1-q)+(1-q)^{-2}$.
Take a uniform mesh of width $h$ so that $hK_1\le1/2$.
The estimate for $\log(1+\xi)$ shows that the logarithm
increment differs from $-hD(t)$ by at most
$K_2h^2$, where $K_2=2K_1^2+(1-q)^{-2}$.
The product relation then telescopes these logarithms,
with total error at most $K_2h$.

The finite geometric sum writes $D(t)$ as
$2\re\sum_{k=1}^N w^k t^{k-1}$ with error at most
$2q^{N+1}/(1-q)$. For $k\ge1$ the rectangle sum of
$t^{k-1}$ on $[0,1]$ differs from $1/k$ by at most $h$.
Indeed $(b^k-a^k)/k$ lies between $(b-a)a^{k-1}$
and $(b-a)b^{k-1}$ by the difference of powers. Sum
over the mesh and use the difference of the upper and lower
rectangles, which is at most $h$. Thus the finite sum for
$D$ differs from $2\re\sum_{k=1}^N w^k/k$ by at most the
geometric error plus $2h/(1-q)$. Refine the logarithm
requests after fixing the mesh. Take arbitrarily small
positive $h$ by rational ceilings and remove the terms
proportional to $h$. Since $Y(0)=1$, this proves the
arithmetic interval estimate
\begin{equation}\label{eq:finite-factor-log}
  \left|\log |1-w|^2+2\re\sum_{k=1}^N w^k/k\right|
          \le 2q^{N+1}/(1-q).
\end{equation}
All bounds are uniform in the rational $w$ in this disk.

\begin{lemma}[A bound from positive real parts]
\label{lem:real-part-coefficients}
If $T(X)=C+\sum_{j=1}^n c_jX^j$, $C$ real rational and
$C\ge0$, has $\re T(u)\ge0$ at every rational unit point,
then $|c_j|\le2C$ for every $j\ge1$.
\end{lemma}
\begin{proof}
On the rational orbit of Lemma~\ref{lem:finite-cauchy} write
$f(u)=\re T(u)$ as a finite Laurent polynomial using
$\bar u=u^{-1}$. Its constant coefficient is $C$.
The constant coefficient of $f(u)u^{-j}$ is $c_j/2$.
The same finite geometric sum bounds put their orbit averages
within rational errors of these constants, each error a fixed
rational divided by the orbit length. The norm of the second
average is at most the first average since $f\ge0$ there.
The disk comparison and arbitrary orbit length yield
$|c_j/2|\le C$ by rational margin removal.
\end{proof}

\begin{lemma}[The reflected product lower bound]
\label{lem:reflected-lower}
Let $P(z)=A\prod_\nu(z-v_\nu)$ be a rational factor product.
Suppose $|P(z)|\le B$ on the rational unit circle and
$|P(0)|^2\ge c_0>0$. For $0<s^2<c<1$ with $s>0$,
replace each factor with $|v_\nu|^2\le c$ by
$1-\bar v_\nu z$. Call the result $F$.
Put
\[
 B_1=\left(\frac{1+B}{1-s}\right)^2,\quad
 C=2+B_1+1/c_0,\quad
 b(r)=1+1/c_0+(2C+4)r/(1-r).
\]
For $0\le r<1$ and $|u|\le r$ the squared norm
$|F(su)|^2$ strictly exceeds
\code{lw-signed-log-value-lower} at $b(r)$.
\end{lemma}
\begin{proof}
At rational unit points the replacement preserves squared
norm, by expansion of both squared norms. At zero it can
only increase it, since the squared norms of replaced
roots are at most $c<1$. Thus $X=|F(0)|^2\ge c_0$.
By Lemma~\ref{lem:finite-cauchy} each coefficient of $F$
has radius $B$. The finite geometric sum bounds
$|F(z)|$ by $(1+B)/(1-s)$ on $|z|\le s$.

Write $F(su)=F(0)\prod_\nu(1-\theta_\nu s u)$, where
$\theta_\nu=\bar v_\nu$ for a replaced factor and
$\theta_\nu=1/v_\nu$ for an unchanged factor.
Since $s^2<c$, the rational number
$s'=s+(c-s^2)/4$ satisfies $s<s'$ and $(s')^2<c$.
Indeed $s<1$, $c-s^2<1$, and
$(s')^2-s^2=(c-s^2)(s/2+(c-s^2)/16)<c-s^2$.
At $|u|\le1$ all $\theta_\nu s u$ have radius at most
$q=s/s'<1$. For replaced factors this follows from
$|v_\nu|<1$ and $s'<1$, and for unchanged factors from
$|v_\nu|^2>c>(s')^2$. In particular every factor is
nonzero there and $Y=|F(su)|^2>0$.

Use \eqref{eq:finite-factor-log} and product comparison.
For the polynomial
\[
 T_N(u)=C+2\sum_\nu\sum_{j=1}^N(\theta_\nu su)^j/j
\]
the real part differs from $C+\log X-\log Y$ by at most
$2m q^{N+1}/(1-q)$, where $m$ is the number of roots.
Take $N$ large enough that the error is at most one and
also at most any smaller prescribed rational request.
Such degrees are computed with the geometric budgets of
Lemma~\ref{lem:geometric}. On the unit circle,
$C+\log X-\log Y\ge2$ by $X\ge c_0$ and $Y\le B_1$.
Thus $\re T_N\ge0$ there. Lemma~\ref{lem:real-part-coefficients}
and the finite geometric sum give
$\re T_N(u)-C\le2Cr/(1-r)$ on $|u|\le r$.
Remove the arbitrary comparison error to obtain
$-\log Y\le1/c_0+2Cr/(1-r)\le b(r)-1$.

The source lower bound is $2^{-(1+j)}$ with
$j=2+\operatorname{num}(2b(r))>2b(r)$.
If $Y\le2^{-(1+j)}$, monotonicity and the finite product
relation would give $-\log Y\ge(1+j)\log 2>b(r)$.
This contradicts the preceding bound. The formulas for
$B_1,C,b(r)$ and the dyadic lower bound are exactly those
in \code{circle-reflected-boundary-budget.lisp},
\code{circle-reflected-harnack-budget.lisp}, and
\code{signed-log-positive-lower.lisp}.
\end{proof}

\subsection{The anchor at the center}

\begin{lemma}[Native center request]\label{lem:circle-center-anchor}
For $s=a+\mathrm i h$ with rational $a>3$, put
$K=1+d((s-1)^{-1})$ and $e_0=1/(64K)$.
The native check \code{lw-euler-value-clearance-check} passes
at $(s,e_0)$.
\end{lemma}
\begin{proof}
For the native samples at the same rational index and request
$0<r\le1$, comparison with finite exponential products gives
\[
 |u_j^s(r)-u_j^1(r)u_j^{s-1}(r)|\le4r,\qquad
 |u_j^{s-1}(r)-u_j^1(r)u_j^{s-2}(r)|\le4r.
\]
These use the sum of the exponential arguments at the same
logarithm endpoint, as in Section~\ref{sec:right-clearance}.
The unit bound is $1+r\le2$ at each rate, and
$|u_j^1(r)-1/j|\le3r$ by Lemma~\ref{lem:native-at-one}
and its proof. The two product comparisons cost at most
$4r+2\cdot4r=12r$. Replacing $(u_j^1(r))^2$ by $1/j^2$
costs at most $3r(2+1)$ before multiplication by the last
sample of radius at most $2$. Hence
$|u_j^s(r)-j^{-2}u_j^{s-2}(r)|\le30r$ and
$|u_j^s(r)|\le j^{-2}+31r$. The first sample equals one.

For $n\ge3$, $\sum_{j=2}^n j^{-2}\le25/36$. Keep the
term $1/4$ at $2$, majorize the others by $1/[j(j-1)]$,
and subtract the difference $1/6-1/9$ at $3$.
The sum of the majorants from $3$ is at most $1/2$ by
telescoping. For the prefix with $n=2^b-1$ the endpoint
term has radius at most $1/(n+1)+31(n+1)r$.
Given any $\tau>0$, choose $b\ge K_s(\tau)$, where
$K_s$ is the cutoff $K(\cdot)$ of
Lemma~\ref{lem:fixed-truncation} at $s$, and $b$ large
enough that $1/(n+1)\le\tau$. For every sufficiently
small positive $r$ the comparison to $V(s,e_0)$ has radius
$e_0/4+\tau$. Multiply the prefix identity by $(s-1)^{-1}$,
whose radius is at most $K$. Add the bounds and remove the
error proportional to $r$, then remove $\tau$. This gives
\[
 |(s-1)^{-1}V(s,e_0)-1|\le25/36+K e_0/4<3/4 .
\]
Thus the normalized value has norm strictly greater than
$1/4$. Writing $\sigma=d(V(s,e_0))$, it follows that
$K\sigma>1/4$, so $\sigma>0$ and $e_0\le\sigma/16$.
These are the clearance checks at the positive rational request
$e_0$ returned by \code{lw-euler-zeta-clearance-request}.
\end{proof}

\subsection{The selected radius and its value floor}

For a real band $(a,b)$ with $0\le a<b$ and a list of
squared norms, the native \code{lw-circle-finite-radius-bracket}
keeps a closed quarter of the current band for each entry.
With midpoint $m$ and width $w$, it keeps $[m+w/4,b]$
if the tested square is at most $m^2$, and $[a,m-w/4]$
otherwise. In either case the square is at distance at
least $(w/4)^2$ from every squared radius in the kept band.
This follows from the difference of squares and the
nonnegativity of both endpoints. Previous separations are
preserved under restriction. By induction, after $k$ entries
the width is $(b-a)4^{-k}$ and all tested squares have the
separation given by its square. This is the graph and
separation proof for the native bracket and gap formulas.

The following construction uses the exact arithmetic of
\code{lw-euler-circle-generated-cover}. Let the center be
$s_0=A+\mathrm i h$ with $A>3$, let $0<U<A$ be the outer
radius, and take a normalized band $(a,b)$ with
$0\le a<b<s<1$ and $s^2<c<1$. The source first runs
the bracket for the single point $(1-s_0)/U$.
In the rest of the argument $a,b$ denote the resulting
endpoints. All physical radii $r$ between $Ua$ and $Ub$
are then separated in square from the pole point $1-s_0$
by the native gap $g_p=U^2(b-a)^2>0$.

Write $C_*$ for the uniform Euler coefficient cap at
$s_0,U$ and let $\sigma=d(V(s_0,e_0))$ from
Lemma~\ref{lem:circle-center-anchor}. Set $c_0=(\sigma/8)^2$ and
$B=1+C_*$. Write $M$ for
\code{lw-circle-inner-factor-count-limit} at $(c,c_0,B)$.
It is the native resolvent truncation count at ratio $c$ and
error $\epsilon=c_0/[2(1+B)^2]$, which is one plus
\code{lw-geometric-index} at $1/c$ and $1/[\epsilon(1-c)]$.
The fuel inequality and Bernoulli's inequality
$(1/c)^j\ge1+j(1/c-1)$ show $c^M\le\epsilon$.
Both are finite induction estimates as in
Lemmas~\ref{lem:fixed-truncation} and~\ref{lem:geometric}.
Define
\[
 \gamma=\min\{((b-a)4^{-M})^2,c-b^2\},\qquad
 \tau=\frac14\left(\frac{\gamma}{5+\gamma}\right)^2.
\]
If $b_0$ is the lower bound in Lemma~\ref{lem:reflected-lower}
with $r=b/s$, the native uniform square bound is
$b_0\tau^M$. Its conversion to distance is
\[
       d_0=\frac{b_0\tau^M}{1+b_0\tau^M}>0 .
\]
We have $d_0\le1$ and $d_0^2\le b_0\tau^M$ by multiplying
out the positive denominator.

Take the source model and factor requests
$\varepsilon,\varepsilon_f$, both bounded by
$\min\{1,\sigma/64,d_0/16\}$.
Let $P$ be the normalized Euler polynomial on the unit
disk and $P_f$ its native factor polynomial.
Lemmas~\ref{lem:native-poly} and
\ref{lem:native-factors} give
\[
 |P_f(z)-V(s_0+Uz,e)|\le\varepsilon+\varepsilon_f+e/4
       \qquad (|z|\le1),
\]
At zero, $|V(s_0,e_0)|^2\ge(\sigma/2)^2$.
The comparison error is at most
$\sigma/64+\sigma/64+\sigma/64$ because $e_0\le\sigma/16$.
Hence $|P_f(0)|^2\ge c_0$. This also ensures its leading
factor and each of its roots are nonzero.
On the whole unit disk the coarse Euler polynomial of target
degree zero and request one has radius $C_*-2$, by its
coefficient mass. Compare $P$ to that coarse polynomial at
error $\varepsilon+1\le2$, using the polynomial comparison
after Lemma~\ref{lem:native-poly}, and then add the factor
error. The resulting cap for $P_f$ is $B$.

Reflect the roots with squared norm at most $c$, as in
Lemma~\ref{lem:reflected-lower}. If their number is $k$, the
reflected polynomial $F$ has
$|P_f(0)|^2\le c^k |F(0)|^2$ by the product norm identity.
On the unit circle $|F(z)|=|P_f(z)|$ in squared norms,
so $|F(0)|\le B$ by Lemma~\ref{lem:finite-cauchy}.
If $k\ge M$, then $c_0\le c^M B^2\le
c_0B^2/[2(1+B)^2]<c_0$, a contradiction. Hence
$k\le M$.

Run the native bracket on just these $k$ roots and take
the midpoint $R$ of the result. This is the normalized
radius from \code{lw-euler-circle-bounded-radius}.
It lies strictly between $a$ and $b$. By the bracket
invariant, each reflected root is separated in squared
radius from $R$ by at least $((b-a)4^{-M})^2$.
Every other root has squared norm greater than $c$ and
therefore gap at least $c-b^2$. At a rational point $z$
with $|z|^2=R^2$, each reflected root satisfies $|z+v|\le2$.
If $|z-v|\le\gamma/(5+\gamma)$, the squared radius
difference would be at most $2\gamma/(5+\gamma)<\gamma$.
Thus $|z-v|^2\ge(\gamma/(5+\gamma))^2$.
Also $|1-\bar vz|\le2$. Consequently the original
factor square is at least $\tau$ times the reflected
factor square. Multiply the bounds over the $k$ replaced
factors. All other factors are identical, and $0<\tau\le1$,
so $|P_f(z)|^2\ge\tau^M|F(z)|^2$.
Apply Lemma~\ref{lem:reflected-lower} to $u=z/s$,
which has radius at most $b/s$. This proves the native
distance floor $|P_f(z)|^2\ge d_0^2$ on the selected circle.

\subsection{Verification of the native arc list}

\begin{proposition}[Native checked cover]\label{prop:native-cover}
For rational height $h$ and $0\le l<u<252$, the native
\code{lw-euler-window-circle-cover} returns a radius
strictly between $l$ and $u$ and two lists satisfying
\code{lw-euler-circle-coverp}. Its central anchor check
also holds at the native request.
\end{proposition}
\begin{proof}
Use $s_0=252+\mathrm i h$, $U=(252+u)/2$, the
normalized original band $(l/U,u/U)$,
$s=(1+u/U)/2$, and $c=(1+s^2)/2$. Their strict
relations follow by multiplication by the positive $U$.
The previous construction gives the physical radius
$r=UR$ in the requested band. It proves the value
floor and the pole gap $g_p$ at this radius. The source
arc lists use local outer radius $L=(U-r)/2$ and
anchor request $e_a=d_0/32$ at each midpoint.
For every point $p$ on the physical circle, its evaluator
differs from the factor model there by at most
$2d_0/16+d_0/128<d_0/4$. Hence
$d(V(p,e_a))\ge d_0/2$, and the anchor check holds since
$e_a\le(d_0/2)/16$. The local radius $L$ is positive
and smaller than $\re p$, because $\re p\ge252-r>L$.
Thus the list of four patch fields satisfies the native
patch predicate.

The pole separation also yields
$|p-1|^2\ge[g_p/(1+g_p+4r)]^2$.
Indeed if the norm of $p-1$ had radius
$\xi=g_p/(1+g_p+4r)\le1$, the sum of the two offsets
$p-s_0$ and $1-s_0$ would have radius at most $2r+\xi$.
The product identity for the difference of their squared
norms would bound the gap by $\xi(2r+\xi)<g_p$.
For the origin, $d(p)\ge\re p\ge252-U$.
These facts rule out the exceptional patch center $1$.

For the local uniform polynomial of model request $e_a$,
every point of its $L$-disk lies in the global $U$-disk.
The global coarse model bounds native values at request
one by $C_*+2$ as in
Section~\ref{sec:native-uniform-polynomial}. Therefore
the local polynomial has radius at most $C_*+4$ on
its disk. Lemma~\ref{lem:finite-cauchy} bounds every
rescaled coefficient by $C_*+4$ in norm and by
$2(C_*+4)$ in coordinate sum. The source neighborhood
slope, twice their nonconstant maximum divided by $L$,
is at most $S=4(C_*+5)/L$.
By the formula for the normalized disk radius, the
actual patch radius is at least
\[
 W=\min\left\{L/2,\ Q(S,(d_0/2)/32),\
                  (252-U)/4,\ \xi/4\right\}.
\]
Here $Q$ is the rounding request in
Lemma~\ref{lem:fixed-truncation}. Its monotonicity in
both arguments follows from its formula and is used
with the slope upper bound and anchor-size lower bound.
This is exactly the uniform grid reach of
\code{euler-circle-regular-grid.lisp}.
The native grid count $N$ is one plus \code{lw-precision-fuel}
at $(2r,W)$. Thus $2r/N\le W$ by the fuel inequality.

The two native charts use
$\phi(t)=(4-t^2+4\mathrm i t)/(4+t^2)$ and its
negative with $-2\le t\le2$. The squared norm of
$\phi(t)$ is one by expansion. For each sign the $i$th
arc, $0\le i<N$, has parameter endpoints
$-2+4i/N$ and $-2+4(i+1)/N$. Its patch is centered
at the chart's midpoint and its required reach is $2r/N$.
Thus the midpoint disk comparison is the zero difference,
and the reach condition holds by the bound on $W$.
The endpoints join consecutively and end at $2$.
Induction on the remaining length proves the two
native chart predicates and hence the cover predicate.
The central anchor check is Lemma~\ref{lem:circle-center-anchor}.
\end{proof}

\section{Stability of the native interior count}
\label{sec:native-count}

\subsection{Comparing counts by finite averages}

Use the formal polynomial derivative
$(\sum_j a_j X^j)'=\sum_{j\ge1}j a_j X^{j-1}$.
For a monomial product the derivative multiplier $i+j$
is the sum of the multipliers from the two factors.
Summing over coefficient pairs gives the product rule.
Thus if a nonzero polynomial $F$
is given by a rational linear factor list with no root on
$|z|^2=r^2$, then
\[
      zF'(z)/F(z)=\sum_v z/(z-v).
\]
This is obtained by the product rule and division by
the nonzero evaluated product. All derivative and
product graphs are finite coefficient loops.

For a root of squared norm less than $r^2$, the summand
is $1/(1-v/z)$. Its geometric polynomial in $v/z$
through degree $N$ has constant term one, with remainder
of radius at most $q^{N+1}/(1-q)$ for any rational
$q<1$ bounding $|v/z|$ on the circle.
For a root of squared norm greater than $r^2$, the
summand is $-(z/v)/(1-z/v)$, and the geometric sum
$-\sum_{j=1}^N(z/v)^j$ has zero constant term and the
same error bound with a cap for $|z/v|$.
To compute $q$ in either case, use $(1+x)/2$ for the
corresponding squared ratio $0\le x<1$, since
$((1+x)/2)^2-x=(1-x)^2/4$. Use the maximum of
these caps over the finite list, or zero for an empty list.
The sum of remainders can be made smaller than any
specified positive rational by the geometric budget.

At $z=ru$ this gives a Laurent polynomial in the rational
unit point $u$, with constant coefficient the number of
interior roots. Now average over the rational orbit of
$\omega$ in Lemma~\ref{lem:finite-cauchy}. The nonconstant
terms have averages bounded by their finite geometric sums
divided by the orbit length. Thus after first choosing the
truncation order and then the orbit length, the average
of $zF'(z)/F(z)$ lies within any prescribed positive
rational error of the integer interior count. For two
factor products one may use the same truncation order
and orbit length by using the larger bounds from both
finite lists. Consequently if their two functions
$zF'(z)/F(z)$ and $zG'(z)/G(z)$ differ by at most
$1/4$ on the rational circle, the two counts are equal.
Indeed the orbit bounds with errors at most $1/8$ show
that their integer difference has absolute value at most $1/2$.

\begin{lemma}[Count through a boundary-safe comparison]
\label{lem:count-comparison}
Suppose $F_0,F_1$ are nonzero rational factor products
and $0<r<T$ are rational. If for every rational $t\in[0,1]$
the polynomial $F_t=(1-t)F_0+tF_1$ has boundary norm
at least $b>0$ on $|z|^2=r^2$, then the two factor
products have the same interior count at $r$.
\end{lemma}
\begin{proof}
Take $B\ge1$ a rational upper bound for $|F_0|,|F_1|$
on the $T$-disk, obtained by the sum of the coefficient
radii there. The disk rules give the same cap for $F_t$.
Let $L=T-r>0$, and put $\delta=b/2$.
For any polynomial $Q$ of radius at most $d$ on the
$T$-disk, its derivative has radius $d/L$ on $|z|^2=r^2$.
Apply Lemma~\ref{lem:finite-cauchy} on the $L$-circle of
the translated polynomial $Q(z+w)$. Its coefficient at $w$
is $Q'(z)$ by the binomial formula.

For two polynomials $F,G$ of radius at most $B+1$ on
the $T$-disk, with boundary norms at least $\delta$ and
difference bounded by $d$ on the disk, we have
\[
 |zF'(z)/F(z)-zG'(z)/G(z)|\le d C_1,\qquad
 C_1=r\left(\frac1{L\delta}+\frac{B+1}{L\delta^2}\right).
\]
Use $F'/F-G'/G=(F'-G')/F+G'(G-F)/(FG)$
and the derivative estimate. Choose a positive rational
$\eta\le\min\{1,b/2,1/[12(1+C_1)]\}$ and an integer
$m\ge1$ with $2B/m\le\eta$.
At each $t=i/m$ with $0<i<m$, apply
Lemma~\ref{lem:native-factors} to $F_t$ at request $\eta$
on the $T$-disk. Each $F_t$ is nonzero since its value
at the rational point $r$ has norm at least $b$. Use
the original factor products at the two endpoints.
The new products have cap $B+1$, boundary norm at
least $\delta$, and adjacent differences at most
$2B/m+2\eta\le3\eta$. Thus adjoining products have
weighted quotient difference at most $1/4$.
The orbit count comparison proves equality of each pair
of counts. Induction on $i$ gives equality at the endpoints.
All interpolations and factor choices in this proof are
finite arithmetic operations.
\end{proof}

\subsection{The generated cover and its refinement requests}

We apply the comparison to the exact native count on each
cover returned by Proposition~\ref{prop:native-cover}.
Use its physical radius $r$ and its boundary floor $d_0$.
Every anchor field of its arcs is $e_a=d_0/32$.
All patch centers avoid $1$, so the anchor returned by
\code{lw-euler-contour-anchor} is $e_a/2$.
The source anchor minimum, starting from one on each list,
is therefore $\alpha=d_0/64$ because both lists have
positive length and $0<d_0\le1$.

For the count refinement models let
$T=(252+r)/2$ and $g=(T-r)/2$. These are the outer
radius and gap of \code{euler-circle-clearance.lisp}.
The source count request at any positive rational
request $a$ is the minimum of
\[
    a,\quad \alpha/4,\quad \alpha/(1+g),\quad
           1/[16(1+r)(1+C_2)],
\]
where $C_2$ is the sum of the native global quotient cost
and the factor quotient cost. These costs are positive
rational expressions in the positive radii, gaps, anchor and
global derivative cap. In particular the last cap has a
rational graph. The derivative polynomial is the formal
derivative of the source uniform polynomial at model error
$g$ for its request of one. Its coefficient mass on the
specified disk gives that cap after adding the source
positive constants. This verifies totality of the request
on the generated covers. Write $e$ for its value.
Then $0<e\le\alpha/4$ and $ge\le\alpha$.

At target degree $j$, the native count model is the
factor polynomial of the global uniform model of error $ge$
on the $T$-disk, factored on that disk at error $e$.
At an offset $z$ with $|z|\le T$ its error against
$V(s_0+z,\varepsilon)$ is at most
$(1+g)e+\varepsilon/4$ by
Lemmas~\ref{lem:native-poly} and~\ref{lem:native-factors}.
Here $s_0=252+\mathrm i h$. At the fixed value request
$\varepsilon=e_a=2\alpha$ this is bounded by
$7\alpha/4$. On the rational $r$-circle
the construction in Section~\ref{sec:construct-cover}
has $|V(s_0+z,e_a)|^2\ge(3d_0/4)^2$.
For any two native count models their convex
interpolations obey the same $7\alpha/4$ comparison on
the circle. This follows by multiplying the two difference
vectors from the same evaluator value by $1-t$ and $t$.
The triangle inequality then bounds their boundary norms
below by $d_0/2$. Apply Lemma~\ref{lem:count-comparison}.
The integer count is the same for every positive request
and target, in particular the native reference target zero
and request one. This proves the instance of
\code{lw-euler-circle-request-count-exact} used for the
generated covers. This proof compares the actual two
factor products and uses auxiliary factorizations of their
interpolation to establish equality of the returned counts.

\subsection{Interior probes and exclusion at count zero}

For a positive native reference count use a positive
tolerance $a$ and the source request $a/[4(1+g)]$
as input to the count request. The resulting $e$ satisfies
$(1+g)e\le a/4$. By count stability the native interior
list is nonempty. Its first root $v$ satisfies $|v|^2<r^2$,
and the factor polynomial there is zero. Thus the
native \code{lw-euler-circle-interior-probe-offset}
has $|V(s_0+v,\varepsilon)|\le a/4+\varepsilon/4\le a$
for every $0<\varepsilon\le a$. This gives the interior
probe at any required positive tolerance.

Conversely, for an inner physical radius $0\le R_0<r$
put $m=(r+R_0)/2$, $s=m/T$, $t=R_0/m$ and
$c=(s^2+(r/T)^2)/2$. These are the relations
$0<s^2<c<(r/T)^2<1$ and $0\le t<1$ used in
\code{euler-window-interior-sample.lisp}.
Let $\sigma=d(V(s_0,e_0))$ for the native center
request and use $c_0=(\sigma/8)^2$, $B=1+C_*$
for the outer $T$-disk. Let $b_0$ be the squared lower
bound from Lemma~\ref{lem:reflected-lower} at
these $s,t,c_0,B$, and set $D=b_0/(1+b_0)$.
These are exactly the native inner floor and its
conversion to distance. In particular $D>0$.
The source \code{lw-euler-window-interior-sample-request}
is $D/8$.

Suppose the native count is zero. For the model request use
the count request applied to
$\min\{1,\sigma/64,D/8\}/(1+g)$, the source
\code{lw-euler-circle-interior-value-request}. Denote its
value by $e$. Its global model error is $ge$ and its factor
error is $e$. Both are at most one and $\sigma/64$, and
their sum is at most $D/8$. The factor model therefore
has center squared norm at least $c_0$ and unit boundary
cap $B$ when rescaled by $T$, by the same comparison as
in Section~\ref{sec:construct-cover}. Its actual factor
list, divided by $T$, has no root of squared norm at
most $c$, since such a root would be strictly inside
the physical $r$-circle. Here we used that every root
count at the chosen request equals the reference
count. The rescaled product thus equals its reflected
product for cutoff $c$.

At an offset $v$ of radius at most $R_0$, take $u=v/m$
so $|u|\le t$ and $Tsu=v$. Lemma~\ref{lem:reflected-lower}
puts the factor polynomial outside the disk of radius
$D$ at this point. If $0<\varepsilon\le D/8$ and the
native evaluator there had radius at most $D/8$, the
factor polynomial would have radius at most
$D/8+D/8+(D/8)/4<D$ by its comparison. This contradiction
shows that the small sample forces positive count at
the exact native interior request.

\subsection{The prescribed collar}

Let $P_0$ be the fixed global polynomial of target zero at
the reference count request divided by sixteen, as in
\code{euler-circle-interior-collar.lisp}. Let $K_0$ be its
native polynomial Lipschitz cap on radius $r$.
The difference of powers shows that this cap bounds
$|P_0(x)-P_0(y)|$ by $K_0$ times any common rational
upper radius for $x-y$, on the $r$-disk.
In detail the native cap dominates
$\sum_{j\ge1}j\,d(a_j)r^{j-1}$, by unrolling its
recurrence and using $\max(1,r)\ge r$ for the disk caps
on the coefficient tails. The native collar width is
$w=\min\{r/2,\alpha/[8(1+K_0)]\}>0$.
The model error of $P_0$ is at most $\alpha/16$.
For any native count model at its refinement request,
its factor polynomial differs from $P_0$ by at most
$\alpha+\alpha/4+\alpha/16$ on the $r$-disk by
uniform polynomial and factor comparisons.

Suppose an interior root $v$ had $|v|^2\ge(r-w)^2$.
There is a rational point $z$ on the $r$-circle with
$|z-v|\le2w$. To construct it, multiply $v$ by a sign
so its real part $a$ is nonnegative, and write the
signed value as $a+\mathrm ib$. Bisect the rational
square-root bracket $[r-w,r]$ for $a^2+b^2$ to width
at most $w/4$ and take its upper endpoint $\rho>0$.
Set $u=2b/(\rho+a)$ in the chart $\phi$ of
Proposition~\ref{prop:native-cover}. Direct substitution
gives
\[
 \left|\rho\phi(u)-(a+\mathrm ib)\right|^2
   =\frac{(\rho^2-a^2-b^2)^2}{(\rho+a)^2+b^2}.
\]
The numerator before squaring is at most $2r(w/4)$
and the denominator is at least $\rho^2\ge(r/2)^2$.
Hence the norm of the difference is at most $w$.
Replace $\rho$ by $r$ in the multiplier of $\phi(u)$,
which changes the point by radius at most $w$, and
restore the sign. This gives the indicated $z$.

Since the factor polynomial vanishes at $v$, the
comparison and Lipschitz estimates give
$|P_0(z)|\le21\alpha/16+2w K_0\le25\alpha/16$.
But at this boundary point
$P_0$ compares to $V(s_0+z,e_a)$ with error at most
$\alpha/16+e_a/4$, and the latter has norm at least
$3d_0/4=48\alpha$. This is a contradiction.
Therefore every interior root is in the disk of radius
$r-w$, with strict inequality in squared norm.

These roots are also nonzero. In Lemma~\ref{lem:circle-center-anchor}
the constant $K\le2$ at $s_0=252+\mathrm ih$, since the
squared norm of $(s_0-1)^{-1}$ is at most $1/251^2$
and its coordinate sum has square at most twice that.
The proof there gives $|V(s_0,e_0)|^2\ge(1/8)^2$.
At zero a native count factor polynomial is within
$5\alpha/4+e_0/4\le3/128$ of this value, using
$\alpha\le1/64$ and $e_0\le1/64$. Thus it does
not vanish there.

\section{The paired contours}
\label{sec:contour}

The following scan uses the native covers and count tests of
Sections~\ref{sec:construct-cover} and \ref{sec:native-count}.
Sections~\ref{sec:right-clearance} and
\ref{sec:native-uniform-polynomial} supply the clearance tests for
its real-centered probes.

\begin{lemma}[A horizontal gap]\label{lem:horizontal}
Suppose $u=-a+\mathrm i v$ is a probe in a real-centered disk of
radius $R\le c<252$. Suppose the exclusion estimates give
$a>251-g$ and $|v|>w$, where $0<g\le1/4$ and $8192g\le w^2$.
Then
\[
 a+8g<R,\qquad 251+7g<R.
\]
\end{lemma}
\begin{proof}
The disk implies $a\le c<252$, and $a>251-g$ implies $a>0$.
Expanding the square gives
\[
 (a+8g)^2-a^2=16ag+64g^2
       \le 4032g+16g<8192g\le w^2<v^2.
\]
Adding $a^2$ and comparing with the disk inequality shows
$(a+8g)^2<R^2$. Both quantities being squared are positive.
Combining the result with $a>251-g$ gives the second inequality.
\end{proof}

Use $c$ and $b$ from $\Good(D)$, the comparison outer radius
$U=(252+c)/2$, and the native width $w$ of the near-real
exclusion on the $c$-disk. At any $g>0$ the native real-probe
tolerance is the minimum of the right-region request at $1+g$,
the near-real request at $252-U$, and $1/16$ of the real-axis
floor at $252-U$. This positive minimum is the definition
in \code{euler-real-circle-probe-geometry.lisp}. If the
probe value and its request are at most that tolerance,
the two exclusion results force $a>251-g$ and $|v|>w$.
Define
\begin{align*}
 g &= \min\{1/4,w^2/8192,(c-251)/8,b/8\}, \qquad u_0=251+2g,\\
 N &= \left\lceil\frac{4(c-u_0)}g\right\rceil,\qquad
 s = \frac{c-u_0}{N+1}.
\end{align*}
These are positive where required, with $s\le g/4$. For
$j=N,\ldots,0$ use the open band
\[
       (c-(j+1)s,\ c-(j+\tfrac12)s).
\]
The first starts at $u_0$ and ends below $251+7g$. If its
count were positive, an arbitrarily fine probe would satisfy
both clearance tests and contradict Lemma~\ref{lem:horizontal}.
Scan the finite bands from this empty circle, and stop
immediately before the first positive successor, if it exists.
Let $E$ be the radius of the selected empty circle. Consecutive
circles have radius difference below $2s\le g/2$.
The empty and positive branches of the scan are proved by
induction on the remaining length.

Use the real normalization radius $\varrho=E-g/4$ and a
modulated band $(E-3g/4,E-g/2)$ with tested inner radius
$E-g$. These numbers satisfy
\[
 251<\varrho<E<c,\qquad
 E-g < E-3g/4 < E-g/2 < \varrho .
\]
When the scan reaches its last band empty, $E>c-s$.
The source uses the seed height $h$ and takes the interior
sample request of the modulated cover as its output request.
The disk radius $c-b$ for the original probe satisfies
$c-b\le c-8g<E-g$.
Thus the original probe lies in the smaller test disk at
height~$h$. If the value at one quarter of the output
request is in its prescribed disk, the native interior
test in Section~\ref{sec:native-count} forces positive count.

In the successor branch let $P$ be the positive real radius.
The probe in the $P$-disk can be taken fine enough for
Lemma~\ref{lem:horizontal}. For its height $v\in[-c,c]$,
take
\[
 h'=-c+g\left\lceil\frac{v+c}{g}\right\rceil .
\]
The set of possible heights is a finite grid with indices
from $0$ to $\lceil2c/g\rceil$. The relocated offset has
absolute imaginary part at most $g$ and absolute real part $a$.
Its squared norm is at most $(a+g)^2$.
Since $a+8g<P<E+g$ we have $a+g<E-g$.
For each possible $h'$ construct a checked modulated contour
and its admissible request. Take the minimum of the positive
requests over this finite grid together with the real-probe
tolerance. The minimum is positive by induction on the grid
length. The source uses this minimum as both the tolerance
and value request of the native interior probe of the positive
successor. At the relocated offset $-a+\mathrm i(v-h')$
the argument of the evaluator is exactly the same as on the
real chart, including the rational height. It therefore passes
the native interior test at height $h'$. In this branch the
output sample request is one and the modulated count is
already positive. In both branches the output request is
positive and computed before admission of the seed probe.
These are the two branches of \code{lw-euler-window-prepare-circles}.

In either branch the positive modulated radius $R$ satisfies
$R<\varrho<c<252-p/q$. Its interior collar lowers $R$ to a
strict pole cutoff $d<R$. These are the inequalities that allow
a geometric comparison with the exterior radius $R$, the real
normalization radius $\varrho$, and the ceiling distance
$252-p/q$. The source proof packages this scan in
\code{euler-window-circle-pair.lisp}. The geometric argument
also explains why the radius and request can be computed before
the original probe point is provided.

\section{Quotients on the modulated contour}
\label{sec:circle-quotient}

\subsection{Derivative coefficients at different centers}

Fix a generated cover from Section~\ref{sec:construct-cover}.
Use $r,T,g,\alpha,e_a$ as in Section~\ref{sec:native-count},
and put $P_e$ for the global Euler polynomial at any
natural target degree and model error $ge$ on the $T$-disk.
For any two global errors $e,e'>0$, the derivative
comparison on the $r$-disk is
\[
          |P'_e(z)-P'_{e'}(z)|\le(e+e')/2 .
\]
Indeed the polynomial difference has radius $g(e+e')$
on the $T$-disk by Lemma~\ref{lem:native-poly} and rational
margin removal. Apply the derivative comparison in
Lemma~\ref{lem:count-comparison} with shell width $T-r=2g$.
The native global derivative cap $D_*$ is $5$ plus
the disk coefficient mass of $P'_1$ at target zero and
radius $r$. Its formula uses $\max(1,r)$ for the disk
mass. The derivative comparison at $e\le1$ therefore
gives $|P'_e(z)|\le D_*$.

For a rational point $p=s_0+z$ with $|z|^2=r^2$, take
the local radius $R=\re(p)/2>0$. The native local
coefficient value $A_q$ at index one and request $q>0$
is the coefficient at degree one of the polynomial
centered at $p$ with disk radius $R$, target one and
model error $\min(1,Rq)$. Its coefficient normalization
multiplies by $R$ and then divides by $R$, so cancels
exactly. By the coefficient comparison after
Lemma~\ref{lem:native-poly},
$|A_q-A_{q'}|\le q+q'$.
Comparing to the local coarse polynomial at model
error one and target zero gives $|A_q|\le D_p$,
where $D_p$ is the local uniform coefficient cap divided
by $R$. Indeed the rescaled coefficient error is at most
two, and the coarse coefficient radius is at most the
coarse coefficient mass. This $D_p$ is the numerator cap
in \code{euler-log-derivative.lisp}.

The two systems of polynomials give
\begin{equation}\label{eq:local-global-derivative}
                     |P'_e(z)-A_q|\le e/2+q .
\end{equation}
To see this, introduce positive requests $e',q'$ and
compare $P'_{e'}(z)$ to $A_{q'}$ in the disk about $p$
of radius $L=\min(T-r,R)$. Translate the global polynomial
to that center by the finite binomial formula.
The two model values on this disk compare to the same
native evaluator, giving a polynomial difference of
radius $ge'+Rq'$ after removing the evaluator request.
Lemma~\ref{lem:finite-cauchy} bounds the coefficient
difference by $(ge'+Rq')/L$. Together with the two
same-center derivative comparisons this gives
$e/2+q+e'/2+q'+(ge'+Rq')/L$. Positive $e',q'$
are arbitrary, so rational margin removal gives
\eqref{eq:local-global-derivative}.

\subsection{The native quotient comparison}

The native circle log value uses the local coefficient
at the point $p$ and the anchor from the arc containing
that point. Every rational unit pair $u$ belongs to
one of the two charts used in
Proposition~\ref{prop:native-cover}. Choose a sign to
make the real part $a$ of the signed pair nonnegative,
and let $b$ be its imaginary part. Then
$t=2b/(1+a)\in[-2,2]$, and substituting $a^2+b^2=1$
shows that $\phi(t)$ is the signed pair. This is the
source inverse parametrization. Since the joined arc
intervals start at $-2$ and end at $2$, selecting the
first right endpoint at least $t$ gives a valid index by
induction on the arc list. All these arcs return anchor
$\alpha$ at the selected index.

From the boundary floor at $e_a=2\alpha$ and the
pointwise comparison of native evaluator values,
the norm of $V(p,\alpha)$ is at least
$3d_0/4-3\alpha/4\ge47\alpha$.
Hence the source check at $(p,\alpha)$ holds.
Put $\Lambda=d(V(p,\alpha))/4\ge\alpha$. For
$0<q\le\Lambda/4$, the norm of $V(p,q)$ is at least
$\Lambda$. Indeed $V(p,\alpha)$ has norm at least
$2\Lambda$, and the error is at most
$(q+\alpha)/4\le5\Lambda/16$.
If $\xi>0$ is the desired log-value request, the source
uses
\[
 q_\xi=\min\{\Lambda/4,Q(L_p,2\xi)\},\qquad
 L_p=1/\Lambda+D_p/\Lambda^2 .
\]
Thus $q_\xi\le1$, $q_\xi\le\xi$, and $L_p q_\xi\le\xi$,
by the formula for $Q$. Write
$\mathcal L_\xi(p)=A_{q_\xi}/V(p,q_\xi)$
for this native value.
The quotient difference identity
\[
 \frac{A}{B}-\frac{A'}{B'}
    =\frac{A-A'}{B}+\frac{A'(B'-B)}{BB'}
\]
and the lower and upper caps just proved give a difference
at most $L_p(q_\xi+q_{\xi'})\le\xi+\xi'$ for two
log requests. Use the coefficient comparison and the
value comparison $(q_\xi+q_{\xi'})/4$ in this identity.

Suppose $0<e\le1$ and $ge\le\alpha$. On the boundary,
$P_e$ has norm at least $\alpha$ by comparison to
$V(p,e_a)$, with error at most $ge+e_a/4$.
The denominator $V(p,q_\eta)$ of the local log value
at an arbitrary request $\eta>0$ also has norm at least
$\alpha$. By \eqref{eq:local-global-derivative},
Lemma~\ref{lem:native-poly}, and the same quotient identity,
\[
 \left|\frac{P'_e(z)}{P_e(z)}-\mathcal L_\eta(p)\right|
 \le \left(\frac1\alpha+\frac{gD_*}{\alpha^2}\right)e
      +\left(\frac1\alpha+\frac{D_*}{4\alpha^2}\right)q_\eta .
\]
The two coefficients are the native global quotient
cost $C_G$ and node cost, respectively. Compare
$\mathcal L_\eta(p)$ to $\mathcal L_\xi(p)$ and remove
the terms proportional to $\eta$, using $q_\eta\le\eta$.
The resulting bound is $C_Ge+\xi$, the comparison to
the native log value on the generated cover.

Factor $P_e$ on the $T$-disk at $0<\varepsilon\le\alpha/4$,
and call its factor polynomial $F$. Its approximation error
is at most $\varepsilon$ by Lemma~\ref{lem:native-factors},
so both boundary denominators have norm at least
$\delta=\alpha/4$. The derivative difference has radius
$\varepsilon/(T-r)$ on the $r$-circle, by finite Cauchy
comparison. The same quotient identity gives
\[
 \left|\frac{F'(z)}{F(z)}-\frac{P'_e(z)}{P_e(z)}\right|
 \le C_F\varepsilon,\qquad
 C_F=\frac1{(T-r)\delta}+\frac{D_*}{\delta^2}.
\]
This is the native factor quotient cost.
No factor vanishes on the boundary. The product identity
in Section~\ref{sec:native-count} gives $F'/F$ as the
full factor sum $\sum_v 1/(z-v)$. Its difference from
the native $\mathcal L_\xi(p)$ is at most
$C_F\varepsilon+C_G e+\xi$.
On the boundary the full factor sum consequently has
the native cap
\[
           B_F=D_*/\alpha+(\alpha/4) C_F
\]
by the bound for the global quotient and the factor
comparison. All reciprocals here have their nonzero
denominators justified by rational disk bounds.

\subsection{Exterior coefficients and their budget}

With the same model let $v$ run over its actual
factor list, split at $|v|^2=r^2$. At $z=ru$ the
inside terms $1/(z-v)$ expand as finite sums of strictly
negative powers of $u$, with an arbitrarily small
uniform geometric remainder. This is the count
calculation in Section~\ref{sec:native-count}, divided
by $z$. The exterior terms expand in nonnegative powers
of $u$ by
\[
       \frac{1}{ru-v}=-\frac1v\sum_{j=0}^N (r/v)^j u^j
                        +E_N(u),\qquad |u|^2=1 ,
\]
where $E_N$ has disk radius at most
$|1/v|$ times the geometric remainder bound.
One can replace $|1/v|$ by its rational coordinate sum.
For a fixed index $j\ge0$, take $N\ge j$ so all
remainders are arbitrarily small. Multiply the finite
Laurent approximation to the full sum by $u^{-j}$
and average on the rational orbit. The constant term
is exactly the native normalized exterior coefficient
$-\sum_{|v|^2>r^2}v^{-1}(r/v)^j$, since the interior
terms contribute no nonnegative powers.
The original boundary sum has radius $B_F$. Removing
the truncation and averaging margins therefore gives
this same cap for the normalized exterior coefficient.
Dividing by $r^j$ gives the physical bound $B_F/r^j$,
at all orders for the chosen native factor list.

For two native models, their whole factor sums compare
on the same boundary to $\mathcal L_\xi$ for every
positive $\xi$. Their difference has radius at most
$C_F(\varepsilon+\varepsilon')+C_G(e+e')$ by
removing $2\xi$. Applying the same Laurent projection
to the difference bounds the difference of their
normalized exterior coefficients by this radius.
If each model uses its source exterior request at
precision $a>0$, its common model and factor request
is at most $a/(1+C_F+C_G)$ by definition, hence its
contribution is at most $a$. This proves the coefficient
comparison on the native refinements as well as the
uniform exterior cap. The interior roots themselves
have the common count and collar of
Section~\ref{sec:native-count}.

\section{The real-centered polynomial model}
\label{sec:real-polynomial}

\subsection{The Euler constant at the pole}

Let $\gamma_\xi$ be the lower endpoint returned by the
source Euler-constant interval at $\xi>0$.
In this section the real cover has center $s_0=252$,
radius $r>251$, and outer radius $T=(252+r)/2$ for
the count models. Put $g=(T-r)/2$, as before.
Let $P_e$ be its global polynomial of model request $ge$.
At offset $z_1=1-s_0=-251$ we have
\begin{equation}\label{eq:gamma-native-derivative}
     |P_e(z_1)-1|\le ge,\qquad
     |P'_e(z_1)-\gamma_\xi|\le e/2+\xi .
\end{equation}
The first estimate follows from the uniform polynomial
comparison to $V(1,e')$ and Lemma~\ref{lem:native-at-one},
with the arbitrary error proportional to $e'$ removed.

For the derivative estimate fix any $0<\tau\le1$.
Choose a dyadic prefix through $n=2^b-1$ with a cutoff
$b$ adequate for target error $\tau$ uniformly on the
half-unit disk about one. Use the surrogate real part
$1/2$ and height $3/2$ of
Lemma~\ref{lem:native-poly}. Increase $b$ so that
$1/N\le\tau$, where $N=n+1$. Choose a positive point
request $q$ below both the uniform transport request
for this prefix and $\tau/(32N^2)$.
For all $1\le j\le N$, form $\ell_j(q)$ by the native
log rectangles. Their arguments $-s\ell_j(q)$ on the
half-unit disk have coordinate sum at most $2N$.
Take a common finite exponential degree $J$ so large
that the scalar tails at both $J-1$ and $J$ for radius
$2N$ are at most $q$. This uses the explicit exponential
tail and halving calculations of
Section~\ref{sec:euler-truncations}, and one may take
$J$ one more than the degree at radius $2N$ and request
$q$. Comparing to common larger degrees, the
polynomials $S_{J-1}$ and $S_J$ differ from the native
sample with request $q$ at their arguments by at most
$2q$. The index-one value is exactly one.

As a polynomial in $s$ define the auxiliary approximation
\[
A(s)=(s-1)\sum_{j=1}^n S_J(-s\ell_j(q))
                  +N S_J(-s\ell_N(q)).
\]
It differs from the native prefix on the half-unit disk
by at most $3Nq\le\tau$, by multiplying the ordinary
sum error by $1/2$ and adding the endpoint error.
Lemma~\ref{lem:fixed-truncation} then gives
$|A(s)-V(s,e')|\le2\tau+e'/4$ on that disk for
every $e'>0$.
Termwise coefficient differentiation gives
\[
 A'(1)=\sum_{j=1}^n S_J(-\ell_j(q))
          -N\ell_N(q)S_{J-1}(-\ell_N(q)).
\]
Every sample in this formula differs from its
corresponding reciprocal by at most $5q$. This adds the
polynomial comparison $2q$ to the reciprocal comparison
$3q$ in Lemma~\ref{lem:native-at-one}. Since
$\ell_N(q)\le N-1$, the expression differs from
$H_n-\ell_N(q)$ by at most $10N^2q\le\tau$.
The interval
$[H_N-U(N)-1/N,H_N-L(N)]$, built from the request-$q$
log bounds, contains $H_n-\ell_N(q)$ and has width
at most $1/N+q$. It overlaps the native Euler-constant
interval at $\xi$, by the finite harmonic comparison in
Section~\ref{sec:log-rectangles}. Hence
$|A'(1)-\gamma_\xi|\le \xi+3\tau$.

Let $L=\min\{1/2,T-r\}>0$ and take an arbitrary finer
global request $e''>0$. On the $L$-disk about $s=1$,
compare the translated $P_{e''}$ to $A$ using the
evaluator with an arbitrary value request. The resulting
polynomial difference has radius at most $ge''+2\tau$
after removing the value request. Finite Cauchy
comparison bounds the derivative difference at one by
$(ge''+2\tau)/L$. The global derivative comparison in
Section~\ref{sec:circle-quotient} bounds the other
difference $|P'_e(z_1)-P'_{e''}(z_1)|$ by
$(e+e'')/2$. Removing the arbitrary $e'',\tau$ gives
the second bound in \eqref{eq:gamma-native-derivative}.

For $0<e\le1$, $ge\le\alpha$ and factor request
$0<\varepsilon\le\alpha/4$, the same global factor
polynomial $F$ has norm at least $1/8$ at $z_1$, by
\eqref{eq:gamma-native-derivative} and the factor
approximation. The global polynomial itself has norm
at least $1/4$ there. For $0<\xi\le1$, we have
$|\gamma_\xi|\le2$ by overlap with the harmonic
interval $[0,1]$ from index one and the width at most
$\xi$. The quotient identity thus gives
\[
\left|\frac{P'_e(z_1)}{P_e(z_1)}-\gamma_\xi\right|
       \le 4(e/2+\xi)+8ge .
\]
At $z_1$ the factor quotient comparison also holds with
clearance $1/8$, using the derivative comparison on
the $r$-disk. Its error is
$\varepsilon(8/(T-r)+64D_*)$.
The resulting estimate for $F'(z_1)/F(z_1)$ is bounded
by the native \code{lw-euler-circle-at-one-log-error}
from \code{euler-circle-log-at-one.lisp}, whose
non-factor terms are $4(e+3\xi)+8(ge+\xi/2)$.

\subsection{The cap after polynomial cancellation}

Suppose the generated real cover has count zero.
Take the factor polynomial at the native joint
request, which is a count request applied to a
minimum of positive rational upper bounds. Its factor
list is all exterior by Section~\ref{sec:native-count}.
On the $r$-circle the factor sum has radius $B_F$
from Section~\ref{sec:circle-quotient}.

At each exterior root $v$ the source geometric ratio
cap is $(1+r^2/|v|^2)/2<1$. The native pole polynomial
in $u$ is the partial geometric sum with factor $-1/v$
and ratio $r/v$ evaluated at $u$. Its truncation
count is at least the resolvent count at the ratio cap
and request equal to the pole error divided by
$1+d(1/v)$. As follows from the fuel formula for
\code{lw-resolvent-truncation-count}, the ratio tail
divided by one minus the ratio cap is at most this
request. The actual ratio and multiplier respect these
caps, so the pole polynomial approximates
$1/(ru-v)$ on the unit disk to the specified pole error.
The source allocates half the current error to the first
root and half to the rest. Induction on the list gives
error at most $\varepsilon$ for the sum on the unit disk
at approximation request $\varepsilon>0$.
This proves the required evaluation bound for
\code{lw-euler-circle-empty-log-polynomial}.

Use $Q(x)$ for that polynomial rescaled by $-1/r$ in
its argument, as in
\code{lw-euler-circle-reflected-log-polynomial}.
At $|x|^2=r^2$, it differs from $F'(-x)/F(-x)$ by at
most $\varepsilon$. When $0<\varepsilon\le1$, every
coefficient at degree $j$ has radius at most
$(B_F+1)/r^j$ by Lemma~\ref{lem:finite-cauchy}.
For the inner radius $\varrho$ with $251<\varrho<r$,
the finite geometric bound gives
\[
        |Q(x)|\le C=(1+B_F)/(1-\varrho/r)
                      \qquad (|x|\le\varrho).
\]
This is the reflected log cap of the source.
The subtraction at one uses the polynomial numerator
$1+Q(251)+(Q(x)-Q(251))/(x-251)$, where the quotient
is represented by the divided-difference polynomial.
On the $\varrho$-circle its norm is at most
$1+C+2C/(\varrho-251)$, since $|x-251|$ has
lower radius $\varrho-251$.
At $x=251$ the same polynomial uses the finite
divided-difference formula and has a value there by
its construction. The source multiplies by the
wide linear inverse for $1/(252-x)$, computed by the
geometric polynomial of the exterior root $-252$ and
the same rescaling argument. Its norm on this circle is
at most $1+1/(252-\varrho)$, by the finite pole error
bound and $\varepsilon\le1$.

Consequently the real normalized polynomial in
\code{euler-circle-real-normalized-model.lisp}
has each degree-$j$ coefficient bounded by
\[
 \frac{(1+C+2C/(\varrho-251))(1+1/(252-\varrho))}
      {\varrho^j}.
\]
This gives the exact geometric cap used to compute the
real contribution to the threshold.

\section{Comparison at the native coefficient requests}
\label{sec:native-normalization}

Use a center $a$ with $\re a=252$ and an outer disk radius
$T<252$, as supplied by the contour. The argument below
also uses the outer radius $126$ for the public Mellin
coefficient. In either case the center anchor uses a
positive request $b$ passing the native clearance check.
Write $C=V(a,b)$, $\sigma=d(C)$, and $b\le\sigma/16$.
When comparing two polynomials only the actual disk radii
of their stated approximations are used.

\subsection{Local inverse and log polynomials}

Write $P_u$ for the uniform Euler polynomial at model
request $u$ on the $T$-disk, with the target specified by
its caller. Let $S$ be $2/T$ times the maximum of the
nonconstant coefficient sizes of $P_b(Tz)$ at target
zero, with maximum zero for an empty list.
For $|z|\le T/2$, the finite geometric comparison gives
$|P_b(z)-P_b(0)|\le S|z|$.
The source neighborhood has radius
\[
       n=\min\{T/2,Q(S,\sigma/32)\}.
\]
Thus $Sn\le\sigma/64$. By Lemma~\ref{lem:native-poly},
on this neighborhood and for $0<u,\eta\le b$,
\begin{align*}
 |P_u(z)-C| &\le u+b+Sn+5b/4 \le 14\sigma/64,\\
 |V(a+z,\eta)-C| &\le b+\eta/4+Sn+5b/4
                              \le 11\sigma/64 .
\end{align*}
For the first line the difference of two polynomials
has radius $u+b$ after comparison to a common evaluator
and removal of its request. In particular all such
denominators have norm at least $\lambda=\sigma/4$.
Also $d(V(a+z,b/2))\ge\sigma-2(11\sigma/64)>\sigma/2$,
so the shifted anchor at $b/2$ passes the check.

Set $K=2/\sigma$, so $|1/C|\le K$, and put
$G=1/\lambda^2$. In \code{euler-uniform-inverse.lisp}
the inverse at error $v$ uses a model at
$u=\min\{b,Q(G,v)\}$. Its residual $1-P_u/C$ has radius
$1/2$. The source inverse is $C^{-1}$ times the partial
geometric sum of this residual. The native resolvent
count at ratio $1/2$ uses tolerance $v/(2(1+K))$.
Multiplying the finite geometric identity by $P_u/C$
shows that the inverse polynomial $I_v$ differs from
$1/P_u$ by at most $v/2$. Its cap is at most $2K$,
and hence at most the source cap $4K$.
Using the reciprocal difference and $Gu\le v/2$ gives
the source comparison
\begin{equation}\label{eq:arithmetic-inverse}
   |I_v(z)-1/V(a+z,\eta)|\le v+G\eta/4
                    \qquad(0<\eta\le b).
\end{equation}

For derivative error $v$, the local uniform derivative
polynomial $D_v$ is the derivative of a model with
error $(T-n)v/2$. Its difference from the native
first coefficient $A_\eta$ at the shifted center,
at any admissible coefficient radius, is at most
$v/2+\eta$. This follows from the proof of
\eqref{eq:local-global-derivative} with $n$ in place of
$r$. The native slope cap $D$ is $5$ plus the disk
coefficient mass of $D_1$ at target zero. Comparing
derivatives at the same center gives
$|D_v(z)|\le D$ for $v\le1$.
The native uniform log polynomial $L_e$ is the exact
product $D_v I_v$ with
\[
     v=Q(M,e),\qquad
     M=1/\lambda+D,\qquad H=1/\lambda+DG/4 .
\]
The product and \eqref{eq:arithmetic-inverse} give
\begin{equation}\label{eq:arithmetic-uniform-log}
   |L_e(z)-A_\eta/V(a+z,\eta)|\le Mv+H\eta
                             \qquad(0<\eta\le b).
\end{equation}
Its source cap is $L_*=4KD$.
These are the model and node costs in
\code{euler-uniform-log-derivative.lisp}.

For a shifted point $s=a+z$, let $W_\xi$ be the native
log-derivative value with any admissible local anchor
and coefficient radius at request $\xi>0$.
The fraction estimate used on the contour also works
with that anchor. If $\Lambda$ is one fourth of its
clearance size and $D_s$ is the native first-coefficient
cap, the source uses the request
\[
         q_\xi=\min\{\Lambda/4,Q(1/\Lambda+D_s/\Lambda^2,2\xi)\}.
\]
The evaluator denominator for requests up to $\Lambda/4$
has radius at least $\Lambda$, by the pointwise
comparison with the anchor. For $0<\eta\le\min(b,\Lambda/4)$
the quotient at $\eta$ therefore differs from $W_\xi$
by at most $\xi+(1/\Lambda+D_s/\Lambda^2)\eta$.
Together with \eqref{eq:arithmetic-uniform-log} and
arithmetic margin removal this gives
\begin{equation}\label{eq:arithmetic-log-local}
                       |L_e(z)-W_\xi|\le e/2+\xi.
\end{equation}
For two different coefficient radii in the native
quotients, compare their first coefficients using the
finite Cauchy bound on the smaller disk and take
arbitrarily fine model requests. This gives a difference
at most the sum of their coefficient requests, by
the same argument as \eqref{eq:local-global-derivative}.
In particular \eqref{eq:arithmetic-log-local} holds
for the anchor and radius used by either of two models
on their common disk.

\subsection{Normalization and physical coefficient transport}

Let $h=\min\{n,d(a)/4,d(a-1)/4\}$ be the native normalized
disk radius. For $|z|\le h$ put $s=a+z$ and
\[
    \lambda_0=d(a)/4,\qquad \lambda_1=d(a-1)/4 .
\]
These are lower radii for $|s|$ and $|s-1|$.
The source linear inverse at anchor $a$ and request $v$
is $a^{-1}$ times the truncated geometric polynomial
in $-z/a$. The modulus of this residual is at most
$1/2$. Its source anchor cap is $1/(2\lambda_0)$,
so the native count ensures error at most $v$, and
the polynomial has norm at most $2/\lambda_0$.
The same comparison applies with anchor $a-1$.

Write $c=1+d(a)+h$, $B=L_*+2c$. If $\gamma_v$ is the
lower endpoint of the native Euler-constant interval
at request $v$, then $|\gamma_v|\le2$ and
$|\gamma_v-\gamma_w|\le v+w$ by the intervals used in
Section~\ref{sec:real-polynomial}.
For $s$ in this disk, write the normalization as
\[
               N(s,L,\gamma)=\frac1s
                   -\frac{L-\gamma s}{s(s-1)} .
\]
Its native polynomial $N_e(z)$ uses $L_v(z)$, $\gamma_v$,
and both linear inverse polynomials at the same
$v=Q(I+U,e)$, where
\[
  I=1+B(1/\lambda_1+2/\lambda_0),\qquad
  U=(1+c)/(\lambda_0\lambda_1).
\]
These are the native inverse and transport costs.
Indeed, the numerator has cap $B$. Replacing both
inverses by exact reciprocals changes their product by
at most $(1/\lambda_1+2/\lambda_0)v$, by expanding
it as two differences. Replacing the outside inverse
costs at most $v$. Thus
\begin{equation}\label{eq:arithmetic-normalize}
        |N_e(z)-N(s,L_v(z),\gamma_v)|\le Iv .
\end{equation}
The native polynomial cap follows from the same
inverse caps and the numerator bound.

The native normalized point value at request $\xi$ is
$N(s,W_w,\gamma_w)$ with $w=Q(U_s,\xi)$, where $U_s$
is the transform cost in \code{euler-log-transform.lisp}.
The transport of $L,\gamma$ in this disk, by
\eqref{eq:arithmetic-log-local} and the Euler intervals,
costs at most
\[
      \frac{v/2+w+c(v+w)}{\lambda_0\lambda_1}\le U(v+w).
\]
The portion \(Uw\) can instead be bounded by \(U_s w\).
In fact the exact coefficient on \(w\) in the transport
is at most \((1+|s|)/|s(s-1)|\), which is bounded by
\(U_s\) using the source inverse caps
\(1+d((s-1)^{-1})\) and \(1+d(s^{-1})\).
The native point request avoids \(s=1\) since
\(|s-1|\ge\lambda_1\).
Consequently the difference is at most
\((I+U)v+U_s w\le(e+\xi)/2\). The point comparison
holds on the overlap of two such disks with a common
native point value, since the anchor \(b/2\) from
either disk passes its check.

For two native normalized polynomials with the same
center, this gives a pointwise difference at most
\((e+e')/2\) on the smaller normalized disk by
margin removal. The finite Cauchy lemma converts this
to a degree-$j$ difference at most
\((e+e')/(2\min(h,h')^j)\).
The physical coefficient at request $\xi$ uses target
$j$ and model request $\min(1,h^j\xi)$. Scaling and
unscaling the coefficient cancel exactly.
In the Mellin coefficient program there is an
additional factor $(-1)^j$, since it uses the
coordinate \(x=a-s\). In particular it differs from
the degree-$j$ coefficient of $N_e(-x)$ by at most
\begin{equation}\label{eq:arithmetic-physical}
                   \xi/2+e/(2h^j).
\end{equation}
For distinct outer radii and anchors at the same
center, put a third, arbitrarily fine model in each
chart and compare those third models on the
smaller disk. This leaves a difference at most
$(\xi+\xi')/2$ between the Mellin coefficients.
This comparison applies to the public coefficient
at outer radius \(126\) with the right-side
clearance request justified in
Section~\ref{sec:construct-cover}.

\subsection{A comparison on the center disk}

Write $r_c=\min(r,n)$ and $\delta_c=\sigma/16$,
the native center disk radius and margin. In addition
to the contour budgets, require $ge\le b$ and factor
accuracy $f\le\delta_c$. The joint request for a
given positive bound $u$ is the count request with input
\[
  \min\{u,\delta_c,b/(1+g),\alpha/[8(1+g)]\}.
\]
It is positive, is at most $u$, and ensures these
conditions when used for both $e$ and $f$. These
inequalities follow from the native count minima
in Section~\ref{sec:native-count}.
On the $r_c$-disk, the neighborhood bounds give
$|P_{ge}|>2\delta_c$ and $|V(a+z,\eta)|\ge\delta_c$
for $\eta\le b$. The factor polynomial is within
$f$, so has clearance $\delta_c$ there. Its full
root list therefore consists of exterior roots for
this disk.

For any such list let $F$ be the factor polynomial.
The difference between $F'/F$ and the native
same-request quotient $A_\eta/V(a+z,\eta)$ has the
bound
\begin{align}
          &C_f f+ E_c(e,\eta), \label{eq:center-quotient-error}\\
 C_f&=\frac1{(T-r_c)\delta_c}+\frac{D_*}{\delta_c^2},
 &E_c(e,\eta)&=\frac{e+\eta}{\delta_c}
                    +\frac{D_*(ge+\eta/4)}{\delta_c^2}.
\notag
\end{align}
Indeed the factor derivative difference costs at most
$f/(T-r_c)$, and the model derivative comparison
costs at most $e/2+\eta$. The remaining terms follow
by subtracting fractions and using the respective
denominator differences $f$ and $ge+\eta/4$.
Thus \eqref{eq:center-quotient-error} has the exact
source constants of the center quotient comparison.
The costs in \eqref{eq:arithmetic-uniform-log} give
the discrepancy with $L_v$ by adding $M Q(M,v)+H\eta$.
The exterior resolvent on the center disk incurs the
specified extraction error $t$. It is bounded there
by the finite geometric identity in
Section~\ref{sec:real-polynomial}. The resulting
center polynomial difference has radius
\begin{equation}\label{eq:native-center-cost}
             E(t,e,f,v,\eta)
                 =t+C_f f+E_c(e,\eta)+M Q(M,v)+H\eta.
\end{equation}
This is \code{lw-euler-circle-center-coefficient-error}.
Finite Cauchy on the $r_c$-disk bounds its degree-$j$
coefficient by $E/r_c^j$. After reflection the
coefficient of the whole root sum is
\(\sum_\nu(-1)^j(-v_\nu^{-j-1})\). Separating the
root list at radius $r$ divides this finite sum into
the interior and exterior portions used by the
modulated comparison.

\subsection{Real coefficient transport}

For the real cover the interior count is zero. Let $e$
be the native joint request at $u$ and use a target-zero
factor model at both errors $e$. Let $Q_r(x)$ denote
its exterior log polynomial reflected and approximated
with error $e$, as in Section~\ref{sec:real-polynomial}.
The notation distinguishes it from the request function
$Q$. By the same geometric comparison, on the disk
of radius \(h_0=\min(\varrho,r_c,h)\),
\[
              |Q_r(x)-L_v(-x)|\le E(e,e,e,v,e).
\]
Here \(v=Q(I+U,e)\) is the request in \(N_e\).
At \(x=251\), set \(G_1=Q_r(251)\). The at-one
comparison and the geometric approximation give
\[
                |G_1-\gamma_v|\le e+E_1(e,e,v),
\]
where \(E_1\) is the native at-one log error proved
in Section~\ref{sec:real-polynomial}.

The exact rational normalization
\(N(252-x,Q_r(x),G_1)\) equals the real normalized
polynomial numerator divided by \(252-x\) on
the \(h_0\)-disk. This is obtained by collecting
terms over \(s(s-1)\), using \(s-1=251-x\).
For its inverse approximation the gain on this
disk is at most the source
\[
        C_0=\frac{(1+c_0')(1+C)}{\lambda_1},
             \qquad c_0'=1+252+h_0 ,
\]
since the divided numerator is
\(((1+G_1)s-1-Q_r(x))/(s-1)\) and
\(|G_1|,|Q_r(x)|\le C\).
The wide linear inverse incurs error at most $e$.
Combining this with \eqref{eq:arithmetic-normalize}
gives the exact bound
\[
 I v+
 \frac{E(e,e,e,v,e)+c_0'[e+E_1(e,e,v)]}{\lambda_0\lambda_1}
 +C_0 e
\]
for the difference of the reflected local polynomial
and the real normalized polynomial. This is
\code{lw-euler-circle-real-native-error-at} at $e$.
Each request inside it is linear in \(e\) for
\(0<e\le1\), since \(Q(B,e)=eQ(B,1)\) for every
nonnegative cost \(B\). Write \(B_0\) for this
whole gain at \(e=1\), and \(M_0\) for the real
normalized cap from Section~\ref{sec:real-polynomial}.

Finite Cauchy on the \(h_0\)-disk gives the coefficient
difference at most \(eB_0/h_0^j\).
The source geometric cap bounds the real model
coefficient by \(M_0/\varrho^j\). For a public
coefficient request \(\xi>0\), transport first
from the public to the contour coefficient at
request \(e\), then use \eqref{eq:arithmetic-physical}
to compare that coefficient to the model at \(e\).
The difference to the real polynomial is at most
\[
       \xi/2+e\bigl(1+1/(2h^j)+B_0/h_0^j\bigr).
\]
The native requested upper bound on \(e\) is the
minimum of \(1/4\) and \(\xi/2\) divided by one
plus the parenthesized gain. Applying the joint
request preserves this upper bound. This proves
the source real native coefficient bound
\(M_0/\varrho^j+\xi\).

\subsection{Modulated coefficient transport}

Use \(x=a-s\) and put \(J(x)=(a-x)(a-1-x)\).
Multiplying the exact local normalization by \(J(x)\)
gives \((1+\gamma_v)(a-x)-1-L_v(-x)\).
In degree \(j\ge2\) the affine part disappears.
Thus by \eqref{eq:arithmetic-normalize}
and finite Cauchy, the degree-$j$ coefficient of
\(J(x)N_e(-x)+L_v(-x)\) has radius at most
\[
          \frac{(1+d(a)+h)(1+d(a-1)+h) I v}{h^j}.
\]
By \eqref{eq:arithmetic-physical}, replacing the
three coefficients of $N_e$ by the native Mellin
coefficients at request $\xi$ adds the costs
\[
    d(a(a-1))\left(\frac{\xi}{2}+\frac{e}{2h^j}\right)
   +d(1-2a)\left(\frac{\xi}{2}+\frac{e}{2h^{j-1}}\right)
              +\frac{\xi}{2}+\frac{e}{2h^{j-2}}.
\]
Their sum with the preceding radius is the native
quadratic log comparison at index $j-2$.
Adding the center comparison \eqref{eq:native-center-cost}
at index $j$ gives the native circle Mellin
coefficient error. The right-hand side is linear
in the common request when \(e,f,t,\xi,\eta\) and the
normalized model error are all set to that request
and are at most one. This follows directly from
the displayed costs and \(Q(B,e)=eQ(B,1)\).

At the finitely many indices of a block, take the
maximum of the gains at request one by list
induction. The source takes the minimum of one and
the tolerance divided by one plus this maximum.
Using the joint request at or below this amount
keeps the actual candidate models in both the
contour and the center budgets. The comparison
error in each index is at most the tolerance.
This proves the precision rule in
\code{euler-circle-mellin-precision.lisp}.

\section{Finite Mellin quadrature}
\label{sec:finite-mellin}

\subsection{Rational interpretation of the quadrature identities}
\label{sec:quadrature-interpretation}

We use \(E(z)\) as shorthand for the finite exponential comparisons of
Section~\ref{sec:euler-truncations}. Each use on a bounded domain is an
estimate with arbitrary positive rational slack, using a finite Taylor
degree computed for that domain and slack. Products, sums, and
reciprocals with stated positive denominator bounds use the
corresponding rational error rules. The same convention applies to
the logarithm and Euler-constant intervals of
Section~\ref{sec:log-rectangles}. A bound on a fixed rational output is
obtained by the squared-disk margin argument in
Section~\ref{sec:euler-truncations}. In particular, the formulas below are
comparisons of finite rational outputs in PA.

Here are the elementary facts used for these comparisons. For real
rational \(u\ge0\), \(E(u)\ge1+u\) by the nonnegative Taylor terms,
and \(E(-u)=1/E(u)\) by the product estimate. The same estimate gives
\(E(z+w)=E(z)E(w)\) and \(E(\mathrm i u)\) of squared norm one in
the comparison sense. For any nonnegative rational \(B\), let \(A(B)\)
be the upper endpoint of the native exponential interval at \(B\)
and request one. On \([0,B]\) the Taylor polynomials and their
derivatives are bounded by \(A(B)\), since the derivatives are the
Taylor polynomials with degree reduced by one. The difference
of their values is therefore at most \(A(B)|u-v|\) by
the difference of powers. This proves the same bound for \(E\).

For the logarithm comparison of \(E(u)\), put \(v=u/n\) with \(n>0\).
Comparison of Taylor coefficients gives
\(0\le E(v)-1-v\le v^2E(v)\).
Telescoping the difference of the \(n\)th powers bounds
\(E(u)-(1+v)^n\) between zero and \(u^2 A(u)/n\).
On the logarithm side, \(v/(1+v)\le\log(1+v)\le v\) follows from
the reciprocal rectangles. Multiplying by \(n\) and comparing the log
at \((1+v)^n\) to that at the exponential approximant shows that
\(\log E(u)=u\) in the comparison convention.
For this comparison the logarithm has slope at most one on
\([1,\infty)\). Its additive law follows by splitting the reciprocal
rectangles at the first factor and rescaling the remaining cells, as
in Section~\ref{sec:construct-cover}. For \(1\le x\le y\le B\)
the reciprocal lower sum gives
\(\log y-\log x\ge(y-x)/B\). This permits inversion of the logarithm
comparison on each bounded positive range. In particular,
\(E(k\log x)=x^k\) for every integer \(k\) and \(x\ge1\), by first
using the log comparison and then taking powers and reciprocals.
All errors here can be made arbitrarily small by selecting a natural
above a specified rational for \(n\) and then increasing the Taylor
degree and rectangle count. The proof also applies uniformly when
the rational \(x,u\) range over fixed bounded intervals.

The native phase at angle \(u\) is the \(m\)th power of
\(R(u/m)=(1+\mathrm i u/(2m))/(1-\mathrm i u/(2m))\).
The source step count gives \(|u|/m\le1\) and \(2u^2/m\) at most
the phase request. The identity \(R(v)-(1+\mathrm i v)
=-(v^2/2)/(1-\mathrm i v/2)\) bounds the residual by \(v^2/2\).
For \(|v|\le1\), the exponential residual after the linear term has
radius at most \(3v^2/4\), by comparison to the geometric ratio
\(1/3\) starting at \(v^2/2\). Since \(R(v)\) and \(E(\mathrm i v)\)
have unit norm, telescoping the \(m\)th powers bounds the phase error
by \(2u^2/m\). The step count is one plus the maximum of precision
fuel for \((|u|,1)\) and for \((2u^2,\text{request})\), so this bound
holds at its exact output. Moreover
\[
 |R(u)-R(v)|\le |u-v|
\]
by subtracting the fractions and using the two denominators of norm
at least one. Applying this to grids with the same arbitrary \(m\),
and then removing their errors, proves
\(|E(\mathrm i u)-E(\mathrm i v)|\le|u-v|\).

Write \(g(x)\) for the rational signal comparison given by the
unsubtracted expression in Section~4 with the upper count covering
\(x\). Its \(k\)th summand is \((x/k-1)_+\Lambda(k)\), where
\((v)_+=\max(v,0)\). The enclosure uses the same finite summands
with prime-log intervals. On \(1\le x\le B\) its slope cap is
\(S(B)=\operatorname{num}(B)+2B+3\) for the positive rational \(B\)
in reduced form. The proof in Section~4 gives this bound for
comparisons, including when \(x\) is itself an approximation.
For \(x\ge1\), each \(\Lambda(k)\) lies in \([0,k]\),
at most \(x\) index terms contribute, and each ramp term is at
most \(x\). The two other terms \(x\log x\) and
\(x(1+\gamma)\) have caps \(x^2\) and \(2x\). Thus
\begin{equation}\label{eq:mellin-signal-growth}
                         |g(x)|\le4x^2 .
\end{equation}
The estimate extends to \(x=E(t)\) by using common rational
approximants in a capped interval with the given slope.

For rate \(\sigma\ge0\), height \(h\), and natural \(j\), set
\begin{equation}\label{eq:mellin-sample-target}
    k_{\sigma,j}(t)=\frac{t^j}{j!}E(-\sigma t),\qquad
    f_{\sigma,h,j}(t)=g(E(t))k_{\sigma,j}(t)E(-\mathrm i ht).
\end{equation}
The expression \(g(E(t))\) on \([0,B]\) uses a finite prefix up to
the numerator of \(2+A(B)\), which is at least \(2+A(B)\).
Thus it has uniformly computable approximants from the exponential,
the signal intervals and the slope bound.

For the exact native sample at \(t\) and request \(e\), the signal
component is the lower log-signal endpoint at
\(\epsilon_g=\min(1,e/[3(1+K_t)])\), where \(K_t\) is one plus the
unit-weight kernel upper bound at request one. To see that the
log-signal interval encloses \(g(E(t))\), its rational address is
the lower exponential endpoint plus \(\delta\), for
\(\delta=\min(1,\epsilon_g)/(4(1+S_t))\).
Here \(S_t=S(1+A(t))\), the signal slope used in the source.
The address lies between \(E(t)\) and \(E(t)+\delta\)
in the comparison sense and has cap \(1+A(t)\).
The base signal interval at that address has width at most
\(\min(1,\epsilon_g)/2\). Its enlargement by \(S_t\delta\)
at both ends gives width at most \(\epsilon_g\)
and encloses \(g(E(t))\) by the slope bound.
The kernel lower bound is obtained by multiplying the reciprocal
exponential interval at \(\sigma t\) by \(t^j/j!\).
The source exponential request divides the kernel request by
\(1+t^j/j!\), so the latter bounds the width, also at \(t=0\).
If \(G_t\) is one plus the sum of the absolute endpoints of the
log-signal interval at one, then \(G_t\) bounds the signal and
the refined lower endpoint for \(\epsilon_g\le1\). This follows
from overlap and width at most one. The kernel and phase
requests are respectively \(e/[3(1+G_t)]\) and
\(e/[3(1+G_tK_t)]\). The kernel is nonnegative and at most
\(K_t\), and both phases have unit norm.
Expanding the difference of the three factors in
\eqref{eq:mellin-sample-target} now bounds the discrepancy of the
native sample by \(e\). In this expansion the costs of the three
component differences are \(K_t,G_t,G_tK_t\) respectively.
This gives the native approximation at the specified point
request.

\subsection{Finite partition comparisons and the improper cutoff}
\label{sec:mellin-partitions}

On the bounded interval \(0\le t\le B\), the source uses the
following caps and slopes for the three components:
\begin{align*}
 K&=A(B), & K'&=K(1+\sigma A(\sigma B)),\\
 G'&=S(2+A(B))A(B), &
 G&=2+|\ell_g(0,1)|+|u_g(0,1)|+G'B .
\end{align*}
Here \([\ell_g(0,1),u_g(0,1)]\) is the native log-signal interval
at \(t=0\) and request one. The kernel power \(t^j/j!\)
and its formal derivative have bounds \(K\), with derivative zero
when \(j=0\). The exponential factor has cap one. Its difference
at \(\sigma u,\sigma v\) is at most \(\sigma A(\sigma B)|u-v|\)
by the positive exponential slope and reciprocals of numbers at
least one. This proves the kernel bounds \(K,K'\).
The signal slope on log addresses was proved above, giving \(G'\),
and comparison to the interval at zero gives \(G\).
Thus the target sample has slope at most
\[
              S_f=KG'+GK'+GK|h|.
\]
This is the source spatial slope for step one. Its spatial
error factor is \(E_f=1+K+G+GK\ge1\).

Use the notation \(\mathcal R_{[a,b]} f\) for the common
comparison of finite partition sums for a given \(f\) with a
computable slope bound on \([a,b]\). For a mesh of maximum cell
width \(\mu\), compare its left-endpoint sum to that for any
refining mesh. Each fine cell contributes a difference bounded
by its length times \(S_f\mu\), by the slope estimate. Comparing
two meshes through the sorted union of their endpoints
then costs at most \((b-a)S_f(\mu+\mu')\).
Sample errors of sizes \(e,e'\) add at most \((b-a)(e+e')\).
The union, partition length identities and sums are finite
operations on lists. An arbitrarily fine uniform mesh has a
computed natural number of cells, obtained by ceiling division
of the length by the requested width. Consequently any
comparison with \(\mathcal R\) can be instantiated with a
PA-computable mesh at the desired rational slack. Concatenation,
linearity, positivity and the disk triangle rule follow on a
common refinement. A piecewise cap with finitely many rational
breakpoints can be used by splitting the mesh at those points.
At a boundary, use each side on its own interval, with a separate
cell-end comparison if needed. The cost of boundary cells with
length tending to zero is bounded by their total length times a
cap on the enclosing interval.

In the native dyadic sum, induction on depth gives the mesh of
\(2^d\) equal leaf intervals and the left tag of each.
For the native bounded integral on \([0,B]\) at request \(e\),
its depth is the halving depth for width \(B^2S_f\) and step
\(e/4\). Its mesh \(\mu\) satisfies \(B S_f\mu\le e/4\).
Indeed the halving loop terminates by the fuel estimate proved in
Section~\ref{sec:euler-truncations}, with initial width \(B^2S_f\).
Its point request is \(e/[4(1+B E_f)]\), so the point
error contribution is at most \(e/4\). This gives a discrepancy
at most \(e/2\) from \(\mathcal R_{[0,B]} f_{\sigma,h,j}\).
In particular the native finite improper value \(I_j(h,\eta)\)
uses bounded request \(\eta/2\) and therefore differs by at most
\(\eta/4\) from the comparison on its cutoff interval.

The source improper cutoff for \(\sigma>2\) uses rate \(r=\sigma-2\)
and mass four. From \eqref{eq:mellin-signal-growth} and the
single term of degree \(j+2\) in the positive Taylor sum at
\(rt\) we have
\[
 |f_{\sigma,h,j}(t)|\le 4 k_{r,j}(t)
      \le \frac{4(j+1)(j+2)}{r^{j+2}t^2}\qquad(t>0).
\]
Write \(H=4(j+1)(j+2)/r^{j+2}\). The comparison on
\([B,B']\), \(0<B\le B'\), has radius at most \(H/B\).
For a left-tagged mesh of width at most \(\mu\), use
\[
  \frac{b-a}{a^2}
     \le (1+\mu/B)\left(\frac1a-\frac1b\right)
       \qquad (B\le a<b).
\]
The sum telescopes, and a bounded-mesh comparison removes the
extra positive term. At the source cutoff
\[
                B=1+(j+2)/r+\lceil 2H/\eta\rceil
\]
this tail is at most \(\eta/2\). In particular the native
improper sample is within \(3\eta/4\) of arbitrarily fine
quadratures at any larger cutoff, in the comparison sense.
The same bounds are uniform over any rationally bounded range of
heights and rates with lower rate \(>2\), by using that lower rate
in the tail coefficient. They also apply to any fixed finite
list of moment orders by taking the maximum of the cutoffs.

\subsection{The kernel mass on a failed-search window}
\label{sec:mellin-window-bound}

For \(r>0\) and interval \([a,b]\subseteq[0,\infty)\),
the finite comparison for the positive kernel is
\begin{equation}\label{eq:mellin-kernel-primitive}
 \mathcal R_{[a,b]} k_{r,j}=H_{r,j}(a)-H_{r,j}(b),\qquad
 H_{r,j}(t)=E(-rt)\sum_{i=0}^j\frac{t^i}{i!\,r^{j+1-i}} .
\end{equation}
To verify it, replace \(E\) in \(H_{r,j}\) by a sufficiently
fine Taylor polynomial and use the formal derivative.
After cancellation between neighboring terms the resulting derivative
is uniformly within arbitrary slack of \(-k_{r,j}\).
Indeed the derivative of the exponential polynomial uses the
same Taylor sum with degree reduced by one, and both exponential
errors are uniformly bounded on \([a,b]\).
For any fixed polynomial \(P\), the binomial expansion of
\(P(v)-P(u)-(v-u)P'(u)\) on a bounded interval has an explicit
bound \(D(v-u)^2\) there, by summing the absolute coefficients
of the terms of degree at least two in \(v-u\).
For a refined partition, the sum of these remainders is bounded
by \(D(b-a)\) times the mesh width. First choose the polynomial
for the desired error, then choose a sufficiently small mesh.
Telescoping proves \eqref{eq:mellin-kernel-primitive}.
Since \(H_{r,j}(0)=r^{-(j+1)}\) and \(H_{r,j}(b)\ge0\), the
mass on any bounded subinterval is at most \(r^{-(j+1)}\).

Let \(Y>1\) be rational, let \([L,U]\) be its native
logarithm interval at request \(1/1000\), and assume
\[
     L\ge64,\qquad (6061/20)U\le j\le(625/2)L,\qquad
                      (4977/5000)^j\le1/128.
\]
Under these assumptions put \(A=U\) and \(B=(3/2)L-1\).
They satisfy \(0<A<B\) since \(U-L\le1/1000\)
and \(L\ge64\).
For \(A\le t\le B\), the log comparison gives
\(\log Y\le t\) and \(2t\le3\log Y-2\).
Consequently a sufficiently fine rational approximation to \(E(t)\)
can be taken in the window \(x\ge Y,\ x^2\le Y^3\).
For the lower condition, taking the maximum with \(Y\) preserves
arbitrary closeness because \(E(t)\ge Y\).
For the upper condition use the strict gap
\(E(2t)\le Y^3 E(-2)\le Y^3/3\).
This gap follows from \(E(\log Y)=Y\), the product identity and
\(E(2)\ge3\). It permits a uniform positive margin for the
approximations throughout \([A,B]\).

Under failure of the rational window search, Section~4 bounds
the signal at each such rational \(x\) by \((C+1)x^{p/q}\).
At \(x=E(t)\) the power comparison is \(E((p/q)t)\).
In fact the \(q\)th power of the latter is \(E(pt)\), which
equals the \(p\)th power of \(E(t)\), by the finite product
identity. For numbers at least one the difference of the
\(q\)th powers has absolute size at least the difference of
the bases. This follows by factoring it and using \(q\ge1\).
Thus the rational power brackets at approximating addresses
converge in the comparison sense, with an explicit error
from the power factorization and the widths of those brackets.
Together with the signal slope this gives
\[
                g(E(t))\le(C+1)E((p/q)t)\qquad(A\le t\le B).
\]
On the rest of the half-line use the growth cap
\(4E(2t)\) from \eqref{eq:mellin-signal-growth}.

Both tails of the rate-\(250\) kernel admit stronger
estimates at these cutoffs. On \(t\le A\), write
\(k_{250,j}(t)=E(50t)k_{300,j}(t)\). Its mass there is
at most \(E(50A)300^{-(j+1)}\).
Since \(A\le j/300\), the multiplier is at most
\((237/200)^j\). Indeed
\[
 E(1/6)\le1+\tfrac16+
               \frac{(1/6)^2}{2(1-1/18)}<237/200
\]
by the geometric bound on the Taylor tail starting at the
quadratic term. On \(t\ge B\), write
\(k_{250,j}(t)=E(-50t)k_{200,j}(t)\).
Here \(B\ge(95/64)L\), so \(50B\ge(19/80)j\).
The first four terms at \(19/80\) have sum greater than
\(100/79\), and therefore \(E(-50B)\le(79/100)^j\).
The mass of that tail up to any finite endpoint is at most
\((79/100)^j200^{-(j+1)}\).
These comparisons use positivity and
\eqref{eq:mellin-kernel-primitive} on the truncated intervals.
Since
\((250/300)(237/200)=(250/200)(79/100)=79/80\),
the combined tail mass is at most
\[
    \left(\tfrac56+\tfrac54\right)(79/80)^j\,250^{-(j+1)}
       \le \frac{(25/12)(252/250)}{128}\,252^{-(j+1)}.
\]
The last bound uses \((79/80)(252/250)=4977/5000\).

With \(\Delta=252-p/q\), the weighted cap for the real
signal on any finite interval \([0,b]\) is thus at most
\begin{align*}
 M_{\rm cap}
   &=(C+1)\Delta^{-(j+1)}
             +4\frac{(25/12)(252/250)}{128}\,252^{-(j+1)}\\
   &\le (4C+5)\Delta^{-(j+1)} .
\end{align*}
The first part uses the mass of the rate-\(\Delta\) kernel on
the whole truncated interval, which only enlarges the cap.
The second uses the two rate-\(250\) tails.
For a real number \(g\) bounded above by \(c\ge0\),
\(|g|\le 2c-g\) follows by the two signs of \(g\).
Apply this row by row to approximating sums on a common
refinement, with \(c\) the applicable signal cap.
The phases have unit radius and the kernel is nonnegative.
The triangle rule and the uniform comparison remove the row
and mesh errors. In particular
\[
  \left|\mathcal R_{[0,b]} f_{252,h,j}\right|
      \le 2M_{\rm cap}
                   -\re\mathcal R_{[0,b]} f_{252,0,j}.
\]
Choose \(b\) to be the native cutoff for \(\eta\), which
does not depend on height. Each native integral has discrepancy
at most \(\eta/4\) on this interval by
Section~\ref{sec:mellin-partitions}. Thus we obtain
\begin{equation}\label{eq:proved-window-quadrature}
 |I_j(h,\eta)|\le(8C+10)\Delta^{-(j+1)}-\re I_j(0,\eta)+\eta .
\end{equation}
The available \(\eta\) allowance includes both discrepancies.
This proves the failed-search quadrature bound at the exact
native improper request.

\subsection{The normalized Euler expression on a right-hand disk}
\label{sec:mellin-normalized-identity}

Fix \(a=252+\mathrm i h\) and use outer radius \(126\) and
the standard clearance request at this center. The clearance
check was proved in Section~\ref{sec:construct-cover}.
Let \(r_0>0\) be the native normalized disk radius and
\(\rho=\min(1,r_0)\).
The auxiliary estimates below use the larger \(2\)-disk
about \(a\), on which \(\re s\ge250\).
Put \(\ell_n=\log n\) in the rational comparison convention
and define finite expressions
\begin{align*}
  Z_N(s)&=\sum_{n=1}^N E(-s\ell_n),\\
  D_N(s)&=-\sum_{n=1}^N\ell_n E(-s\ell_n),\\
  B_N(s)&=\sum_{n=1}^N\Lambda(n)E(-s\ell_n).
\end{align*}
These have uniformly computable polynomial approximations
on the \(2\)-disk, by first taking log enclosures and then
Taylor degrees for the bounded arguments. For \(D_N\),
take log approximants as constant rational coefficients.
For a Taylor approximation to \(Z_N\), differentiation
uses the degree reduced by one. Take the degree large
enough for both approximation bounds. This gives the
formal derivative comparison with \(D_N\).

On this disk we have
\(|E(-s\ell_n)|\le n^{-4}\), since the exponential modulus
uses the rate \(\re s\ge4\) and \(E(4\ell_n)=n^4\).
For \(n\ge2\), \(n^{-4}\le n^{-2}/4\), and
\(\sum_{n=2}^M n^{-2}\le1\) by comparison with
\(\sum_{n=2}^M (1/(n-1)-1/n)\).
Thus \(Z_N(s)\) has a uniform lower radius \(3/4\).
The same telescope gives \(\sum_{n=1}^M n^{-2}\le2\).
It also bounds the tails of \(Z_N,D_N,B_N\) since
\(\ell_n\le n\) and \(\Lambda(n)\le n\). For instance
\(\sum_{n=N+1}^M n^{-3}\le1/N\).

The arithmetic log-derivative comparison is finite.
For every integer \(k\ge1\),
\[
                 \sum_{d\mid k}\Lambda(d)=\ell_k .
\]
To see it, factor \(k\) into prime powers by ordinary
induction and use uniqueness from Euclid's lemma.
For each prime with exponent \(u\) exactly its first \(u\)
powers contribute, and the additive logarithm law gives
the equality. The factorization code can be obtained by
bounded divisions by the least prime divisor, decreasing
the integer at each step.
The exponential product law and \(\ell_m+\ell_n=\ell_{mn}\)
give
\[
 Z_N(s)B_N(s)+D_N(s)
   =\sum_{\substack{u,v\le N\\uv>N}}\Lambda(v)E(-s\ell_{uv}).
\]
The terms at \(uv\le N\) cancel exactly by the divisor sum.
For \(uv>N\) the norm of the summand is bounded by
\(u^{-4}v^{-3}\le u^{-3}v^{-2}/N\).
The product of the finite sums thus gives a uniform error
of at most \(4/N\) on the \(2\)-disk. In particular
\(-D_N/Z_N\) differs from \(B_N\) by at most \(16/(3N)\)
in the comparison sense, with the lower radius of \(Z_N\)
available for rational approximations.

We next compare the normalized point evaluator in
Section~\ref{sec:native-normalization} with these finite
expressions. For \(N=2^k\) the source prefix consists of
\((s-1)\) times the sum up to \(N-1\) and the endpoint
\(N E(-s\ell_N)\), at the specified sample requests.
After refinement the endpoint and the change from \(N-1\)
to \(N\) cost at most \((N+d(s-1))N^{-4}\).
Lemma~\ref{lem:fixed-truncation} gives the uniform comparison
with the native evaluator \(V(s,\eta)\) at error at most
\(\eta/4\) plus an arbitrarily small margin.
For a prescribed margin choose \(k\) above the uniform
tail threshold at \(\re s\ge250\) and the bounded
\(d(s)\) on the \(2\)-disk, and then choose finite sample
requests small enough for its prefix. This proves that
\((s-1)Z_N(s)\) compares to the native evaluator as
the margin and \(\eta\) decrease.

For the first coefficient at \(s\) in the \(\rho\)-disk,
compare a sufficiently fine native polynomial centered
at \(s\) with a Taylor approximation to
\((s-1)Z_N(s)\) on the unit disk about \(s\).
This disk is inside the \(2\)-disk about \(a\).
The evaluator comparison and the finite Cauchy bound on
that unit disk give, up to the uniform margins, the
comparison of derivative coefficients.
The native first coefficient at request \(\eta\) differs
from the refined one by at most \(\eta\) plus its request
by Section~\ref{sec:circle-quotient}.
On the finite-expression side the coefficient is
\[
                        Z_N(s)+(s-1)D_N(s)
\]
up to the uniform Taylor margin, by the formal derivative
of the explicit polynomial.
Choosing \(N\), the native refinement and the Taylor degrees
with successively smaller errors gives the desired comparison
to this expression at arbitrary slack. All denominators
in the normalized point model on this disk have positive
lower bounds. For the native log derivative this follows
from the shifted anchor at \(b/2\) in
Section~\ref{sec:native-normalization}, and for the finite
Dirichlet comparison it follows from the bound on \(Z_N\)
and \(|s-1|\ge249\).
By the quotient rule for rational errors and the divisor sum
estimate, the native log-derivative expression can therefore
be compared to
\[
                          \frac1{s-1}-B_N(s)
\]
with arbitrarily small slack as the requests are refined
and \(N\) is taken sufficiently large.
Inserting this in the normalization, the rational identity
\begin{equation}\label{eq:dirichlet-normalized}
  N\!\left(s,\frac1{s-1}-B_N(s),\gamma\right)
      =\frac{B_N(s)}{s(s-1)}
                    -\frac1{(s-1)^2}+\frac{1+\gamma}{s-1}
\end{equation}
is obtained by clearing the nonzero denominators. Here
\(\gamma\) is represented by the same Euler-constant
intervals in the signal and in the normalization.

\subsection{Integration of the finite prime-power prefix}
\label{sec:mellin-ramp-integration}

For complex \(s\) in the \(\rho\)-disk about \(a\),
consider the finite comparison
\[
                       \mathcal R_{[0,B]} g(E(t))E(-st).
\]
The primitive comparison used in
\eqref{eq:mellin-kernel-primitive} extends to complex rates
\(r\ne0\) by the same derivative calculation and triangle
rule. In particular the primitives for \(E(-rt)\) and
\(tE(-rt)\) are \(-E(-rt)/r\) and
\(-E(-rt)(t/r+1/r^2)\), respectively.
At every positive rate floor, their endpoint terms at a
large rational \(B\) are arbitrarily small by the linear
or quadratic Taylor lower bound at \((\re r)B\).
The bounds are uniform in the stated disk.

For each \(n\le N\), the ramp
\((E(t)/n-1)_+\) vanishes below \(\ell_n\).
Above \(\ell_n\), multiplying by \(E(-st)\) gives
\[
                    n^{-1}E(-(s-1)t)-E(-st).
\]
To split the partition at the log comparison, use a
rational endpoint of an arbitrarily fine interval for
\(\ell_n\). The cost of its difference from the stated
breakpoint is bounded by the interval width times a
cap for the two expressions on the bounded interval.
This cap is computed from the exponential estimates.
Refining all \(N\) intervals and summing their errors
makes that cost arbitrarily small. Since
\(E(\ell_n)=n\), the contribution of the two primitives
at the lower breakpoint is
\[
                   E(-s\ell_n)
                          \left(\frac1{s-1}-\frac1s\right)
                     =\frac{E(-s\ell_n)}{s(s-1)} .
\]
For a fixed finite \(N\), choose \(B\) above all the log
intervals and large enough for the sum of the endpoint
errors. For the other two signal terms,
\(-E(t)\log E(t)\) and \((1+\gamma)E(t)\),
use \(\log E(t)=t\) and the two primitives at \(s-1\).
Their combined lower-endpoint contribution is
\(-1/(s-1)^2+(1+\gamma)/(s-1)\).
Thus the finite-prefix signal quadrature agrees with
\eqref{eq:dirichlet-normalized} up to any prescribed slack
at a sufficiently large cutoff.

To compare with the full signal, first fix a rational
tail cutoff \(B_0\) whose tail after \(B_0\) is smaller
than the prescribed tolerance, using the inverse-square
bound in Section~\ref{sec:mellin-partitions} at order
zero and a uniform positive lower rate \(>2\).
Choose \(N\ge 1+A(B_0)\) so that every omitted ramp
vanishes on \([0,B_0]\). The prefix and full signal
both have the growth bound \eqref{eq:mellin-signal-growth},
so their quadrature difference after \(B_0\) is bounded
by twice the same uniform tail. This choice of \(N\)
can also meet the tail and divisor-sum constraints in
Section~\ref{sec:mellin-normalized-identity}.
Now choose \(B\) large enough for the finite-prefix
primitive estimate. The triangle rule proves that
the finite quadrature for the shifted order-zero
expression and the native normalized point evaluator
agree as their tolerances decrease. Concretely, given
any positive slack, take \(B_0\) with the tail smaller
than a fixed fraction of the slack, take a common
dyadic \(N\) large enough for the prefix and the
Dirichlet comparison, and then take all sample
requests and mesh widths small enough for these fixed
finite sums. Any two sufficiently large cutoffs give
the same comparison up to the tail estimate.
Each comparison has a finite rational witness at the chosen slack.

\subsection{Comparison with the native moment coefficient}
\label{sec:mellin-moment-comparison}

Fix a natural order \(j\), native integral request
\(\eta>0\), and coefficient request \(e>0\).
For an arbitrary small positive margin \(\tau\) we
construct a finite polynomial on the \(\rho\)-disk
with coordinate \(x=a-s\). Use a cutoff \(B\) at or
above the native moment cutoff and large enough for
the uniform tail bound on the shifted order-zero
expression to be at most \(\tau/3\).
Use a common finite mesh of \([0,B]\) with left tags
\(t_i\) and lengths \(w_i\).
Choose its maximum width to give partition discrepancy
at most \(\tau/3\) both for order \(j\) at the center
and order zero on the shifted disk. The caps in
Section~\ref{sec:mellin-partitions} are uniform here:
replace \(\sigma\) in the slope by \(253\) and
\(|h|\) by \(|h|+1\). At each tag choose a rational
approximation \(A_i\) to \(g(E(t_i))E(-at_i)\) and
form the polynomial
\[
                Q_\tau(x)=\sum_i w_i A_i S_M(xt_i)
\]
with \(M\ge j\).
The Taylor degree is taken so that the error for
\(E(xt_i)\) is at most
\(\tau/[24(1+B)(1+A(B)^2)]\)
uniformly for \(|x|\le\rho\), using the argument
cap \(2B\).
The error in each \(A_i\) can be taken at most
\(\tau/[24(1+B)(1+A(2B)+A(B))]\).
The estimates \( |S_M(xt_i)|\le A(2B)\),
\(|g(E(t_i))E(-at_i)|\le4A(B)^2\),
and \(\sum_i w_i=B\) give total polynomial
evaluation error at most \(5\tau/24<\tau/3\) from
these approximations.
Together with the mesh and tail choices, this bounds
the value discrepancy of \(Q_\tau(x)\) from the
shifted order-zero quadrature comparison by \(\tau\).
At degree \(j\), the coefficient of \(Q_\tau\) is exactly
\(\sum_i w_i A_i t_i^j/j!\), since \(M\ge j\).
Its additional point error is at most \(\tau/24\)
using \(t_i^j/j!\le A(B)\).
The mesh comparison and the native cutoff estimate
therefore give
\begin{equation}\label{eq:auxiliary-mellin-moment}
     \big|[x^j]Q_\tau-I_j(h,\eta)\big|\le \tau+3\eta/4 .
\end{equation}
All choices here are finite bounds, even when \(\tau\)
is so small that the chosen mesh extends far beyond the
native cutoff.

Let \(N_v\) be the native local normalized polynomial
at any model request \(v>0\), in the offset \(s-a\).
The normalization comparison in
Section~\ref{sec:native-normalization} bounds it
against a normalized point value at request \(\xi\)
by \((v+\xi)/2\).
The finite-prefix identity and partition comparison
above connect that point value to the shifted
order-zero quadrature at arbitrarily fine requests.
Taking \(\xi\) sufficiently small for each prescribed
margin and removing the margin yields
\[
          |N_v(-x)-Q_\tau(x)|\le v/2+\tau
                              \qquad(|x|\le\rho).
\]
The finite Cauchy bound therefore compares their
degree-\(j\) coefficients within
\((v/2+\tau)/\rho^j\).
At request \(e\), the native public coefficient
\(T_j(h,e)\) differs from the degree-\(j\) coefficient
of \(N_v(-x)\) by at most \(e/2+v/(2r_0^j)\) by
\eqref{eq:arithmetic-physical}.
Together with \eqref{eq:auxiliary-mellin-moment}, the
total comparison radius between the two fixed rational
outputs is at most
\[
       \frac{3\eta}4+\frac e2+
         \tau(1+\rho^{-j})+\frac v2(r_0^{-j}+\rho^{-j}).
\]
The margin and the model request are arbitrary positive
rationals, so the squared-disk argument removes the
last two terms. In particular we obtain the required
native bound
\begin{equation}\label{eq:proved-mellin-moment}
                 |I_j(h,\eta)-T_j(h,e)|\le 5\eta+e/2 .
\end{equation}

\section{The window coefficient input and the final deduction}
\label{sec:composition}

Section~\ref{sec:native-normalization} connects the contour
models to the finite estimators at their native requests.
The finite quadratures of Section~\ref{sec:finite-mellin}
bound those estimators when the window search fails.
We combine the bounds at a common order and request.

\subsection{The window coefficient bound}
\label{sec:window-estimate}

Write \(T_j(h,e)\) for the public Mellin coefficient at
center \(252+\mathrm i h\), degree \(j\), and request \(e\).
Use the notation \(C,p,q\) of \(\Good(D)\) and put
\(\Delta=252-p/q\).
Let \(Y>1\) be rational, \([L,U]\) its native logarithm
interval at request \(1/1000\), and \(G_0\) the native
geometric index of \(5000/4977\) at target \(128\).
The bound uses a natural \(j\) with
\begin{equation}\label{eq:window-band}
     64\le L,\qquad j\ge G_0,\qquad
                   (6061/20)U\le j\le (625/2)L .
\end{equation}
If the rational window search fails, then at every
positive rational request \(e\),
\begin{equation}\label{eq:window-mellin-bound}
      |T_j(h,e)|
       \le (8C+10)\Delta^{-(j+1)}-\re T_j(0,e)+e .
\end{equation}
The geometric index gives \((4977/5000)^j\le1/128\).
Apply \eqref{eq:proved-mellin-moment} at heights \(h,0\)
and use \eqref{eq:proved-window-quadrature} between the
two quadratures. This gives the proposed radius in
\eqref{eq:window-mellin-bound} plus \(11\eta\) for every
rational \(\eta>0\). If that proposed radius were negative,
a sufficiently small \(\eta\) would give a negative upper bound
on a norm. The radius is therefore nonnegative, and rational
margin removal gives the claim. It is the same-request case of the
bound in \code{euler-window-coefficients.lisp}.

\subsection{The three-order algebra}

The basic quadratic identity is finite. Let
$B_j=\sum_\nu b_\nu t_\nu^{j+1}$ with $t_\nu=(a-\rho_\nu)^{-1}$,
and let $Q(z)=z(z-1)$. For $j\ge2$,
\begin{equation}\label{eq:quadratic}
 a(a-1)B_j+(1-2a)B_{j-1}+B_{j-2}
   =\sum_\nu b_\nu Q(\rho_\nu)t_\nu^{j+1}.
\end{equation}
Indeed, the two terms of lower order insert respectively
$a-\rho_\nu$ and $(a-\rho_\nu)^2$ into the bracket after
factoring $b_\nu t_\nu^{j+1}$. Expanding $Q$ gives
\[
 a(a-1)+(1-2a)(a-\rho_\nu)+(a-\rho_\nu)^2
      =\rho_\nu(\rho_\nu-1).
\]
Finite induction gives the equality for the entire sum.
In the normalized principal part, $b_\nu=-m_\nu/Q(\rho_\nu)$,
where $m_\nu$ is the positive weight. With $\rho_\nu\ne0,1$,
the right side of \eqref{eq:quadratic} is the negative
logarithmic-derivative pole coefficient. Thus a bound on three
consecutive normalized coefficients bounds the pole coefficient
using only addition and multiplication of finite rational data.
The exact shifts and scale factors are incorporated in
\code{euler-window-quadratic-coefficients.lisp}.

\subsection{The exact decay index and logarithmic threshold}

The source decay index implements the geometric comparison
of Lemma~\ref{lem:geometric} using a resolvent count.
For a rational \(0<q_0<1\), natural \(k\), rational \(a_0\ge1\),
mass \(H\ge0\), and tolerance \(\tau>0\), put
\[
  \delta=\frac{1-q_0}{2(1+k)},\quad
  B=\frac1{1-\delta},\quad
  q_1=\frac{2q_0}{1+q_0},\quad
  A=a_0+1+\frac1{B-1}.
\]
Bernoulli's inequality gives
\((1-\delta)^k\ge1-k\delta\ge(1+q_0)/2\)
and \(i+a_0\le AB^i\) for every natural \(i\).
Hence
\[
       H(i+a_0)^k q_0^i\le H A^k q_1^i .
\]
For \(i\) at or above the native resolvent count at
\(q_1\) and tolerance \(\tau/(1+H A^k)\),
the right side is at most \(\tau\).
This gives the bound of
\code{lw-polynomial-decay-index} at its exact output.

Use the modulated circle radius \(R\), real
normalization radius \(\varrho\), and pole cutoff
\(0<d<R<\varrho<\Delta\). Write \(m>0\)
for the stable modulated count and \(M_0\) for the
native real normalization cap. Let \(a=252+\mathrm i h\)
and put
\[
    W(t)=d(a(a-1))+d(1-2a)t+t^2 .
\]
Here \(d(u)\) in \(W\) denotes the sum of the component
absolute values of a complex rational, as above.
Use the three masses and radii
\begin{align*}
   H_1&=B_F, & R_1&=R,\\
   H_2&=(8C+10)W(\Delta)/\Delta, & R_2&=\Delta,\\
   H_3&=M_0W(\varrho), & R_3&=\varrho .
\end{align*}
For a window with first order \(J\), let
\[
   A_J=\frac{m}{8[3(J+m+3)]^{m-1}},\qquad
   L_j=\frac{A_J}{d^{j+3}},\qquad
   F_J=\frac{A_J}{\max(1,d)^{J+m+2}} .
\]
For \(J\le j<J+m\) we have \(0<F_J\le L_j\).
For a mass \(H_\nu\) at radius \(R_\nu>d\), the
ratio of \(H_\nu/R_\nu^{j+2}\) to \(L_j\) is bounded by
\[
       \frac{8\,3^{m-1}H_\nu d}{m}
               (j+2+3+m)^{m-1}(d/R_\nu)^{j+2},
\]
using \(J\le j\).
The native \code{lw-circle-pole-tail-minimum-order}
applies the decay index with this leading constant,
degree \(m-1\), offset \(3+m\), and tolerance \(1/8\).
Let \(N\) be the maximum of the three resulting
indices. For \(j\ge N\) all three terms are bounded
by \(L_j/8\). These computations are the source
minimum-index function in
\code{euler-circle-window-budget.lisp}.

The geometric index \(G_0\) equals one plus precision
fuel for width \(128\) at step \(5000/4977-1>0\).
The precision fuel bound and Bernoulli give
\((5000/4977)^{G_0}\ge128\).
The source threshold takes
\(K=\max(64,G_0,N,m+1)\) and returns \(2^{2K+9}\).
For \(Y\) at or above it, the rational log bounds
give \(L\ge(2K+9)/2-1/1000\ge K+4\).
For this bound, refine the logarithm comparison
as in Section~\ref{sec:construct-cover}, split the
interval at the successive powers of two and rescale
each subinterval to \([1,2]\). Every such subinterval
has lower integral comparison \(1/2\), and the log
interval at \(Y\) has width at most \(1/1000\).
The native first order is \(J=\lceil(615/2)L\rceil\).
Then \(J\ge N,G_0\) and \(J\ge(6061/20)U\), since
\(U\le L+1/1000\) and \(L\ge64\).
Moreover \(J+(m+1)\le (625/2)L\), since
\(m+2\le L\) and \(J\le (615/2)L+1\).
Every order \(j,j+1,j+2\) needed for
\(J\le j<J+m\) therefore satisfies
\eqref{eq:window-band}.

\subsection{Composition from the window estimate}

The tolerance in the source is
\[
                    \frac{F_J}{8(1+3W(1))}.
\]
Compute the native block precision with this tolerance
and the \(m\) indices starting at \(J\).
Apply the joint request to obtain the common error \(e\).
Section~\ref{sec:native-normalization} shows the
resulting model keeps its count and collar and pays
comparison error at most this tolerance at each
index. The block precision divides the tolerance
by one plus a nonnegative gain, so \(e\) is also
at most the tolerance.

At this joint request the source chooses the first
large index in the finite block scan of
Corollary~\ref{cor:late}. The test is on the power
sums scaled by \(d\), with start \(J+2\) and \(m\)
successive tests beginning at \(J+3\). The cap of
the scan is \(A_J\), so the scan succeeds.
The source subtracts one to get the pole coefficient
index and two more to get the selected window index
\(j\in[J,J+m-1]\). Its pole contribution
\(P_{j+2}\) satisfies \(|P_{j+2}|>L_j\).

Suppose the window search failed. Then
\eqref{eq:window-mellin-bound} holds for the
three selected orders. The real native bound gives
\(-\re T_\ell(0,e)\le M_0/\varrho^\ell+e\)
at each such order. Transporting the public
modulated coefficient to the contour model's native
coefficient at the same request costs at most \(e\)
by Section~\ref{sec:native-normalization}.
Thus each of the three modulated coefficients has
radius at most
\((8C+10)\Delta^{-\ell-1}
      +M_0/\varrho^\ell+3e\).
The three weights of the quadratic combination sum
to \(W(1)\). Applying them to the two geometric
terms gives precisely
\(H_2/\Delta^{j+2}+H_3/\varrho^{j+2}\).
The modulated circle comparison then bounds
\(|P_{j+2}|\) by those terms plus the exterior bound
\(H_1/R^{j+2}\), the rounding term \(3eW(1)\),
and the model comparison error.
The last two together are at most \(F_J/8\).
The three geometric terms are at most \(3L_j/8\).
Their sum is at most \(L_j/2\), contradicting
the strict pole bound.

The search therefore succeeds with the
exact threshold, joint error and
selected index of the native program. The search
verification in Section~4 supplies the returned
address, both window inequalities, and the
specified positive interval row.

\begin{lemma}[Admission of a finer probe]\label{lem:admit}
Suppose the evaluator at a fixed argument has the Cauchy estimate
$|V(e_1)-V(e_2)|\le e_1+e_2$, and $r>0$ is the computed
sample request. If $0<e\le r/4$ and $|V(e)|\le e$, then
$|V(r/4)|\le r$.
\end{lemma}
\begin{proof}
The triangle rule gives
$|V(r/4)|\le r/4+2e\le3r/4\le r$.
All comparisons are rational disk inequalities.
\end{proof}

\begin{theorem}[Arithmetic windowed ceiling reduction]
\label{thm:pa-window}
The arithmetic target \(\mathcal W\) of \eqref{eq:target}
has a derivation in first-order Peano arithmetic.
\end{theorem}
\begin{proof}
The paired scan of Section~\ref{sec:contour} constructs
both covers and the sample request before a point is
provided. Its construction uses \(c>251\) and \(b>0\).
This includes the data with \(b\ge c\).
The selected real count is zero, the proposed modulated
count is a natural \(m\), and the computed cutoff and radii
have \(0<d<R<\varrho<\Delta\).
The prepared threshold is a positive natural since it
is a power of two with the computed natural exponent.
Its graph is total also at \(m=0\). In that branch the
source rational power returns one for the exponent
\(m-1=-1\), and rational reciprocal returns zero at zero,
the arithmetic completion used in these formulas.
The pole-tail mass multiplying the decay scale is then
zero, so its resolvent count uses the positive input
tolerance and the same rate in \((0,1)\).
The search proof below will use the pole floor only for
\(m>0\).

Every graph used to produce these outputs has a
natural-number construction of the kind in
Lemma~\ref{lem:recursion}. For the evaluator and root
models the measures are the stated fuel, degree, and
grid and factor-list lengths. The dyadic sums decrease
depth, the real and modulated cover scans decrease their
remaining band and grid lengths, and the arithmetic
tail counts and precision maxima use fixed natural
fuel or a finite list. All code comparisons and
reciprocals have rational graphs. These constructions
give unique codes \(T,r\) in \(\mathsf{Prepared}\);
the paired scan proves \(r>0\).

Now take an eligible \(Y,u,e\).
If \(c-b<0\) the disk test is false. If \(c-b=0\)
it forces \(u=0\), which also satisfies the smaller
test disk inclusion of the paired scan. In every
eligible case the probe admission lemma gives the
small-value test at the computed request.
In the seed branch it forces positive modulated
count by the interior test. In the successor branch
the relocated real probe already forces that count.
Thus the paired contours satisfy the hypotheses of
the coefficient argument at the prepared threshold.
By \eqref{eq:window-mellin-bound} and the composition
above, failure of the native search is impossible.
The interval-search and mesh derivations of
Section~4 give the output \(x,L,U\) with
\(\Cert(D,Y,x,L,U)\). These prove the universal
implication in \eqref{eq:target}.

All the steps are rational graph comparisons, finite
inductions and first-order deductions on the codes.
The definitions and computation graphs on their
indicated domains are eliminated by the arithmetic
interpretation in Section~\ref{sec:graphs}.
This gives the claimed PA derivation. If
\(\Zero(D)\) is given, instantiate it at one quarter
of the prepared request and use the same witness
for every \(Y\ge T\).
\end{proof}

\section{Relation to replay of the ACL2 proof}
\label{sec:transfer}

The ACL2 proof record is a sequence of certification events.
Its underlying logic uses first-order equality, recursive definitions
and an induction rule whose general form uses ordinal measures
\cite[Sections 10--13]{acl2logic}. The procedure for choosing and
combining induction schemes is described in \cite{acl2induction}.
To use an ACL2 proof inside this paper one needs a translation of the
\emph{proof} into the arithmetic fragment, including its primitive
axioms and induction justifications.

\begin{theorem}[Finite replay criterion]\label{thm:replay}
Consider a finite derivation in a language of coded objects. Suppose
each primitive symbol has a graph which $\PA$ proves total and
single-valued on the interpreted domain, each used initial axiom
has a $\PA$ proof of its translation, and every extra definition
or induction step meets the graph and measure conditions of
Section~\ref{sec:graphs}. Suppose the other logical steps are
first-order consequence and equality rules. Then the translation
of the last formula has a finite proof in $\PA$.
\end{theorem}
\begin{proof}
Introduce each function graph in derivation order. For a function
already given by a total single-valued graph $F(\mathbf a,y)$,
replace a term $f(\mathbf t)$ by a bound variable $y$ together
with the graph relation for its arguments. Every atomic occurrence
can be flattened this way. In $\PA$, existence and uniqueness
permit the usual substitution and equality rules on the flattened
formulas. Guard all sorted quantifiers by code validity. Interpret
the domain operations and truth values by their graph formulas.
The graph predicates for recursive definitions and their totality
and uniqueness are justified by Lemma~\ref{lem:recursion}.
An ordinary first-order rule maps to its logical derivation under
this interpretation, and a natural-measure induction maps to the
strong-induction derivation proved after Lemma~\ref{lem:recursion}.
Paste the translated derivations at each reference to a previous
line. Since the original derivation is finite, repeated pasting
produces a finite $\PA$ derivation with the desired last formula.
\end{proof}

The arithmetic fragment can be presented as a definitional
extension of $\PA$ and subsequently eliminated by this graph
translation. The criterion applies to a fixed finite derivation.

The checked dependency closure has 558 books in the repository
snapshot used for this manuscript. The project audited recorded
termination justifications of its 347 exported recursive
definitions and found natural-valued measures. The configured
book suite separately certified 647 books and 3649 exported theorem
events. These facts locate the available proofs and establish
what was checked by ACL2. To apply Theorem~\ref{thm:replay} one must
still discharge the initial formulas and proof steps actually used,
including local induction schemes and proof simplifications. A
termination audit on the exported definitions alone does not perform
this replay.

The derivation of Theorem~\ref{thm:pa-window} follows
the finite computations and derives their inequalities directly.
In particular it uses finite primitive evaluations and a finite
divisor-sum argument for the Mellin estimate. This is a
mathematical PA proof of the stated target. The criterion above
describes a separate way to obtain a proof from the full ACL2
event sequence.

\section{External comparison with the classical statement}
\label{sec:comparison}

This section works with ordinary real and complex analysis to identify
the arithmetic sentence's meaning. Its limits and compactness argument
are external to the PA derivation target.

Let $Z(s)$ be the pole-removed zeta function, so that $Z(s)=(s-1)\zeta(s)$
for $s\ne1$ and $Z(1)=1$. At a rational point $s=\sigma+\mathrm i t$
with $\sigma>0$, the finite ACL2 evaluator at request $e$ has the
following external comparison:
\begin{equation}\label{eq:analytic-bound}
      |V(s,e)-Z(s)|\le e/4.
\end{equation}
Here is a way to read the certified data. Reciprocal rectangles converge
to the real logarithm as their widths go to zero. At fixed integer
$n\ge1$ the Dirichlet sample thus tends to $\exp(-s\log n)$, using the
Taylor remainder bound for the exponential and continuity in its
argument. For fixed dyadic cutoff $X=N+1=2^k$, the finite model tends to
\[
 U_N(s)=(s-1)\sum_{n=1}^N n^{-s}+X^{1-s}.
\]
The estimate in \code{euler-fixed-truncation.lisp} bounds the
discrepancy from $V(s,e)$ by $e/4+A(s)/(k-1)+B(s,k)\eta$ with
$\eta$ the primitive request. For fixed $k$, let $\eta$ tend to
zero. On the unit intervals the integral of
$s\lfloor v\rfloor v^{-s-1}$ telescopes, giving
\[
 \sum_{n=1}^N n^{-s}
  =N X^{-s}+s\int_1^X\lfloor v\rfloor v^{-s-1}\,dv .
\]
Splitting $\lfloor v\rfloor=v-\{v\}$ and integrating the first term
gives
\[
 U_N(s)=s-(s-1)X^{-s}
                -s(s-1)\int_1^X \{v\}v^{-s-1}\,dv .
\]
This equality also holds at $s=1$ directly. The absolute value of
the integral tail is at most $X^{-\sigma}/\sigma$.
The fractional-part identity for $Z$ then shows that $U_N(s)\to Z(s)$.
The identity is established in the Lean development
\code{Soft/ZetaReal.lean}. Sending $k$ to infinity after the
primitive limit proves \eqref{eq:analytic-bound}.

If a zero of $Z$ is in the interior of the coded disk, continuity and
density give a rational point with $|Z(s)|<3e/4$ for every $e>0$.
By \eqref{eq:analytic-bound} it supplies an encoded probe at request
$e$. Conversely, take a sequence of probes with requests tending to
zero inside the fixed closed disk. A convergent subsequence exists
by compactness of this disk in the right half-plane. Along the
subsequence $|Z(s)|\le 5e/4$, so its limiting point is a zero of $Z$.
It cannot be at one because $Z(1)=1$. This argument interprets the
existential zero condition regardless of multiplicity.

The rational signal brackets name
\[
 G(x)-C x^{p/q},\qquad
 G(x)=\sum_{n\le x}\Lambda(n)(x/n-1)-x\log x+(1+\gamma)x.
\]
The log and power comparisons follow from their enclosing
rectangles and bisections. The harmonic intervals use
$H_n-\log n-1/n\le\gamma\le H_n-\log n$, obtained by summing the
same reciprocal bounds and taking the standard harmonic limit.
The prime-power calculation is exact. Consequently the returned
$L>0$ witnesses the strict inequality. A scaled margin
$m=L/x>0$ also works since $x^{p/q}\le x$ on this range.

Let $0\le\theta<1$ and let a zeta zero have real part $\beta>\theta$.
The right-half-plane clearance combined with \eqref{eq:analytic-bound}
excludes a zero with $\beta>1$:
at rational points converging to such a zero choose a fixed rational
lower rate greater than one and a sufficiently small fixed value
request. The uniform positive native lower bound and
\eqref{eq:analytic-bound} give a positive lower bound for $|Z|$
at those points and at their limit. With $\beta\le1$, choose rational
$\theta'\in[\theta,\beta)$ and take $c$ strictly between
$\max(251,252-\beta)$ and $252-\theta'$. A rational height
and positive rational slack place the zero strictly inside the
disk. For $C'>\max(C,0)$ rational the arithmetic reduction gives
a violation with ceiling $C'x^{\theta'}$ at every sufficiently
large rational start. To handle a fixed sufficiently large real
start $Y$, take nearby rational starts tending to $Y$. Extract a
convergent subsequence of the witness points from a compact
surrounding interval. The limit lies in $[Y,Y^{3/2}]$ by continuity
of the two endpoints. The smoothed prime sum is locally a
finite sum of continuous ramps and therefore continuous.
The limit satisfies $G(x)\ge C'x^{\theta'}>Cx^\theta$.
This proves the comparison with the classical recurring-window
formulation.

\appendix
\section{Source index}

The following index identifies the finite proof interfaces.
The book names refer to \code{acl2/} in the repository.

\begingroup
\raggedright
\begin{description}
\item[Arithmetic entry point.]
 \code{euler-window-encoded-zero.lisp} defines the approximate-zero
 predicate and states the final universally quantified witness
 theorem. \code{euler-window-ceiling-reduction.lisp} defines the
 threshold and composes the two generated contours with the search.

\item[Enclosures and window.]
 \code{log-bounds.lisp}, \code{euler-bounds.lisp},
 \code{rpow-bounds.lisp}, \code{signal-bounds.lisp}, and
 \code{signal-continuity.lisp} contain the independent precision
 budgets and transport bound. \code{rational-window.lisp}
 proves the finite mesh and endpoint coverage.

\item[Euler evaluation.]
 \code{complex-exp.lisp}, \code{euler-continuation.lisp}, and
 \code{euler-fixed-truncation.lisp} give finite exponential and
 Euler comparisons. The right-half-plane inverse uses
 \code{mobius-sieve-inversion.lisp},
 \code{dirichlet-right-inverse-budget.lisp}, and
 \code{euler-right-clearance.lisp}. The real and near-real estimates
 use \code{euler-real-axis-clearance.lisp} and
 \code{euler-near-real-clearance.lisp}.

\item[Contour and count interfaces.]
 \code{euler-window-circle-cover.lisp},
 \code{euler-circle-positive-count-witness.lisp},
 \code{euler-circle-generated-cover.lisp}, and
 \code{euler-window-real-scan.lisp} supply the checked counts and
 the scan. The two branches are composed in
 \code{euler-window-circle-pair.lisp}.

\item[Coefficient contradiction.]
 \code{power-sum-domination.lisp}, \code{power-sum-reversal.lisp},
 and \code{power-sum-search.lisp} supply the block estimate.
 \code{euler-circle-joint-model.lisp} supplies the common factor
 requests and collar. \code{kernel-window-index-block.lisp},
 \code{euler-window-quadratic-coefficients.lisp}, and
 \code{euler-circle-window-budget.lisp} connect the block to the
 search. \code{euler-circle-real-native-bound.lisp} supplies the
 real-centered coefficient bound used by
 \code{euler-circle-pair-window-certificate.lisp}.

\item[Mellin quadrature.]
 \code{mellin-partitions.lisp} and \code{mellin-integral.lisp}
 give the finite sums and cutoff. The failed-search estimate is
 in \code{mellin-dyadic-normalization.lisp}. The comparison with
 the normalized coefficient is in
 \code{euler-mellin-native-coefficients.lisp}.
\end{description}
\endgroup

\begin{thebibliography}{9}
\bibitem{companion}
The \texttt{landau-windows} project,
\emph{Landau's ceiling step on windows, with explicit constants},
companion manuscript, \code{classical/paper.tex}.

\bibitem{acl2logic}
M. Kaufmann and J S. Moore,
\emph{A Precise Description of the ACL2 Logic}, 1998.
\url{https://www.cs.utexas.edu/~moore/publications/km97a.pdf}.

\bibitem{acl2induction}
J S. Moore, \emph{ACL2 Induction Heuristics}, 2020.
\url{https://www.cs.utexas.edu/~moore/publications/acl2-induction-heuristics.pdf}.
\end{thebibliography}
\end{document}
